% Leçon 31 — Grammaire : Les connecteurs 首先… 然后……其次……最后… ; 既…..又 ; 只有….才…. ; 只要……就….. — Chinois Tle A — Cours signé OnBuch+
% Programme officiel MINESEC. Compiler avec : tectonic ZH31-grammaire-les-connecteurs.tex
\newcommand{\DOCMATIERE}{Chinois}
\newcommand{\DOCNIVEAU}{Tle A}
\newcommand{\DOCMODULE}{Module 4 : Activités économiques et monde du travail}
\newcommand{\DOCLECON}{ZH31}
\newcommand{\DOCTITRE}{Leçon 31 — Grammaire : Les connecteurs 首先… 然后……其次……最后… ; 既…..又 ; 只有….才…. ; 只要……就…..}
% OnBuch+ — préambule « cours signés » ENRICHI (compilé avec tectonic / XeLaTeX)
% Système visuel « Pop » d'OnBuch+ : fond crème (jamais blanc), bordures encre
% épaisses, ombres « dures » décalées, titres Archivo Black en capitales.
% Version MATHÉMATIQUES : graphes (pgfplots/tikz), boîtes pédagogiques
% (définition, propriété, méthode, attention, le savais-tu, exercice résolu,
% exercice, TP, bilan), page de garde et signature OnBuch+ ancrée en bas.
%
% Macros attendues dans main.tex AVANT \input{preamble} :
%   \DOCMATIERE  \DOCNIVEAU  \DOCMODULE  \DOCLECON (numéro)  \DOCTITRE
\documentclass[11pt]{article}

\usepackage{fontspec}
\usepackage[french]{babel} % français = langue principale ; le chinois via \zh{..} / textebox (xeCJK)
\usepackage{xeCJK}
\usepackage{pifont} % \ding{51} \ding{55} utilisés par les modèles
\usepackage[a4paper,margin=2cm,top=3.4cm,bottom=2.8cm,headheight=1.4cm,footskip=1.3cm]{geometry}
\usepackage{amsmath,amssymb,amsfonts}
\usepackage{mathtools} % \xmapsto, \coloneqq... souvent utilisés par les modèles
\usepackage{cancel} % simplifications barrées (bilans de forces)
% \abs{z} (module, valeur absolue) et \norm{u} (norme), utilisés par les modèles
\providecommand{\abs}{}\let\abs\relax\DeclarePairedDelimiter{\abs}{\lvert}{\rvert}
\providecommand{\norm}{}\let\norm\relax\DeclarePairedDelimiter{\norm}{\lVert}{\rVert}
\usepackage[table]{xcolor} % [table] : fournit aussi \rowcolors (alternance de lignes)
\usepackage{pagecolor}
\usepackage{tikz}
\usetikzlibrary{arrows.meta,positioning,calc,shapes.geometric,shapes.misc,shapes.arrows,decorations.pathmorphing,decorations.pathreplacing,decorations.markings,patterns,shadows,mindmap,trees,fit,backgrounds,matrix,intersections,angles,quotes}
\usepackage{pgfplots}
\usepackage[version=4]{mhchem} % équations nucléaires \ce{^{235}_{92}U}
\usepackage[european,siunitx]{circuitikz} % schémas électriques
\pgfplotsset{compat=1.18}
\usepgfplotslibrary{fillbetween,groupplots} % aires entre courbes (calcul intégral)
\usepackage{enumitem}
\usepackage{fancyhdr}
% [most] charge toutes les bibliothèques tcolorbox (listings, minted, raster,
% documentation...) et gonfle inutilement la mémoire TeX sur un document long
% (~300 boîtes) : on ne charge que ce qui est réellement utilisé.
\usepackage[skins,breakable]{tcolorbox}
\usepackage{graphicx}
\usepackage{colortbl}
\usepackage{booktabs}
\usepackage{tabularx}
\usepackage{diagbox} % case d'en-tête diagonale des tableaux à double entrée
% Colonnes extensibles centrée / gauche / droite (C, L, R), souvent utilisées par les modèles
\newcolumntype{C}{>{\centering\arraybackslash}X}
\newcolumntype{L}{>{\raggedright\arraybackslash}X}
\newcolumntype{R}{>{\raggedleft\arraybackslash}X}
\usepackage{array}
\usepackage{multicol}
\usepackage{multirow}
\usepackage{siunitx}
% parse-numbers=false : accepte aussi \SI{2,5 \times 10^{-3}}{...}
\sisetup{locale=FR,output-decimal-marker={,},group-minimum-digits=4,parse-numbers=false}
\DeclareSIUnit\ohm{\ensuremath{\symup{\Omega}}} % U+2126 absent de la police
\DeclareSIUnit\division{div} % réglages d'oscilloscope (ms/div, V/div)
\DeclareSIUnit\tr{tr} % tours (vitesse de rotation, ex. \SI{1500}{\tr\per\minute})
\DeclareSIUnit\molar{M} % concentration molaire (ex. \SI{3}{\molar}, \SI{1}{\micro\molar})
\DeclareSIUnit\an{an} % année (le modèle confond parfois avec \year, un compteur TeX) — ex. \SI{5}{\kilo\meter\per\an}
\DeclareSIUnit\FCFA{FCFA} % franc CFA (ex. \SI{500}{\FCFA\per\kilo\gram})
\DeclareSIUnit\degreeBrix{\textdegree Brix} % teneur en sucre d'un sirop/jus (réfractométrie)
\DeclareSIUnit\month{mois} % le modèle utilise parfois \month, un compteur TeX, comme unité de durée
\usepackage{qrcode}
% \degree : commande fréquemment utilisée par les modèles pour un angle ou une
% température en dehors de \SI{}{} (non fournie par les paquets chargés).
\providecommand{\degree}{^{\circ}}
% \qed (fin de démonstration), utilisé par les modèles sans amsthm
\providecommand{\qed}{\hfill$\square$}
% \barre : double barre des valeurs interdites dans un tableau de variations (tkz-tab)
\providecommand{\barre}{\ensuremath{\|}}
% environnement proof (amsthm non chargé) utilisé par les modèles
\ifdefined\proof\else\newenvironment{proof}[1][Démonstration]{\par\noindent\textit{#1.}\ }{\qed\par}\fi
\usepackage{lastpage}
\usepackage{hyperref}
\hypersetup{hidelinks}

% ============================================================
% Palette « Pop » OnBuch+ — mêmes hex que la classe P (plus_ui.dart)
% ============================================================
\definecolor{poporange}{HTML}{F59321}
\definecolor{popdark}{HTML}{E07A0C}
\definecolor{popcream}{HTML}{FFF6E8}
\definecolor{popink}{HTML}{1C1714}
\definecolor{poppurple}{HTML}{7A5AE0}
\definecolor{popgreen}{HTML}{2E9E6A}
\definecolor{poppink}{HTML}{E0457B}
\definecolor{popblue}{HTML}{2F7DE1}
\definecolor{popgold}{HTML}{C98A12}
\definecolor{popmuted}{HTML}{8A7F76}
\definecolor{popline}{HTML}{EADFD2}
\definecolor{popgray}{HTML}{8A7F76}
\colorlet{popgrey}{popgray}
% Alias tolérants (le modèle nomme parfois les couleurs différemment)
\colorlet{popwhite}{white}
\colorlet{popblack}{popink}
\colorlet{popred}{poppink}
\colorlet{popyellow}{popgold}
% Teintes claires (fonds de tableaux/encadrés) — le modèle nomme parfois
% les couleurs avec un suffixe L (ex. popblueL) sans les avoir définies.
\colorlet{poporangeL}{poporange!20}
\colorlet{popdarkL}{popdark!20}
\colorlet{poppurpleL}{poppurple!20}
\colorlet{popgreenL}{popgreen!20}
\colorlet{poppinkL}{poppink!20}
\colorlet{popblueL}{popblue!20}
\colorlet{popgoldL}{popgold!20}
\colorlet{popredL}{popred!20}
\colorlet{popgrayL}{popgray!20}
% Teintes claires pour fonds de tableaux / zones
\colorlet{popblueL}{popblue!10!white}
\colorlet{poporangeL}{poporange!14!white}
\colorlet{popgreenL}{popgreen!12!white}
\colorlet{poppurpleL}{poppurple!12!white}
\colorlet{poppinkL}{poppink!12!white}
\colorlet{popgoldL}{popgold!14!white}

\pagecolor{popcream}

% ============================================================
% Typographie
% ============================================================
\defaultfontfeatures{Ligatures=TeX}
\setmainfont{PlusJakartaSans-Medium.ttf}[Path=../../../fonts/,BoldFont=PlusJakartaSans-Bold.ttf]
% Symboles, API et emojis absents de la police principale : repli sur DejaVu Sans (caractères actifs)
\newfontface\symfont{DejaVuSans.ttf}[Path=../../../fonts/]
\makeatletter
\newcommand\symactive[1]{\catcode#1=13 \begingroup\lccode`\~=#1 \lowercase{\endgroup\def~}{{\symfont\char#1}}}
\newcommand\symblank[1]{\catcode#1=13 \begingroup\lccode`\~=#1 \lowercase{\endgroup\def~}{}}
\newcommand\symhyph[1]{\catcode#1=13 \begingroup\lccode`\~=#1 \lowercase{\endgroup\def~}{\mbox{-}}}
\newcount\symc
\newcommand\symrange[2]{\symc=#1 \@whilenum\symc<#2 \do{\expandafter\symactive\expandafter{\the\symc}\advance\symc by 1 }}
\newcommand\symblankrange[2]{\symc=#1 \@whilenum\symc<#2 \do{\expandafter\symblank\expandafter{\the\symc}\advance\symc by 1 }}
\makeatother
\symrange{"0250}{"02FF}\symactive{"2713}\symactive{"2714}\symactive{"2717}\symactive{"2718}\symactive{"2610}\symactive{"2611}\symactive{"2612}
\symrange{"2600}{"27BF}
\catcode"2705=13 \begingroup\lccode`\~="2705 \lowercase{\endgroup\def~}{{\symfont\char"2713}}\symblank{"26BD}\symblank{"26A1}
\symrange{"2460}{"257F}\symactive{"223C}\symactive{"2248}\symactive{"2260}
\symactive{"01F8}\symactive{"01F9}\symactive{"1E3E}\symactive{"1E3F}
\symhyph{"2011}\symhyph{"2E1A}\symblankrange{"1F300}{"1FAFF}\symblank{"FE0F}
% Caractères chinois : WenQuanYi Zen Hei (sous-ensemble GB2312 + pinyin), gras/italique simulés
\setCJKmainfont{WenQuanYiZenHei-subset.ttf}[Path=../../../fonts/,AutoFakeBold=3,AutoFakeSlant=0.2]
\xeCJKsetup{CJKspace=false}
% Pinyin : voyelles à caron (ǎ ǐ ǒ ǔ ǖ ǘ ǚ ǜ) absentes de la police principale -> caractères actifs
\newfontface\pyfont{WenQuanYiZenHei-subset.ttf}[Path=../../../fonts/]
\newcommand\pyactive[1]{\catcode#1=13 \begingroup\lccode`\~=#1 \lowercase{\endgroup\def~}{{\pyfont\char#1}}}
\pyactive{"01CD}\pyactive{"01CE}\pyactive{"01CF}\pyactive{"01D0}\pyactive{"01D1}\pyactive{"01D2}\pyactive{"01D3}\pyactive{"01D4}\pyactive{"01D5}\pyactive{"01D6}\pyactive{"01D7}\pyactive{"01D8}\pyactive{"01D9}\pyactive{"01DA}\pyactive{"01DB}\pyactive{"01DC}
% Maths en sans-serif (Fira Math) avec chiffres et lettres droites de Plus
% Jakarta, pour que les formules se fondent dans le texte.
\usepackage[math-style=ISO,bold-style=ISO]{unicode-math}
\setmathfont{FiraMath-Regular.otf}[Path=../../../fonts/]
\setmathfont{PlusJakartaSans-Medium.ttf}[Path=../../../fonts/,range={up/num,up/{Latin,latin},\mathup}]
\usepackage{icomma} % virgule décimale sans espace en mode math
% Caractères mathématiques Unicode absents des polices de texte (Plus Jakarta,
% Archivo Black) : rendus par leur équivalent mathématique, en texte comme en
% formule — sinon ils s'affichent en carré vide (ex. « Arithmétique dans ℤ »).
\usepackage{newunicodechar}
\newunicodechar{ℝ}{\ensuremath{\mathbb{R}}}
\newunicodechar{ℤ}{\ensuremath{\mathbb{Z}}}
\newunicodechar{ℂ}{\ensuremath{\mathbb{C}}}
\newunicodechar{ℕ}{\ensuremath{\mathbb{N}}}
\newunicodechar{ℚ}{\ensuremath{\mathbb{Q}}}
\newunicodechar{𝔻}{\ensuremath{\mathbb{D}}}
\newunicodechar{α}{\ensuremath{\alpha}}
\newunicodechar{β}{\ensuremath{\beta}}
\newunicodechar{γ}{\ensuremath{\gamma}}
\newunicodechar{δ}{\ensuremath{\delta}}
\newunicodechar{ε}{\ensuremath{\varepsilon}}
\newunicodechar{θ}{\ensuremath{\theta}}
\newunicodechar{λ}{\ensuremath{\lambda}}
\newunicodechar{μ}{\ensuremath{\mu}}
\newunicodechar{σ}{\ensuremath{\sigma}}
\newunicodechar{τ}{\ensuremath{\tau}}
\newunicodechar{φ}{\ensuremath{\varphi}}
\newunicodechar{ψ}{\ensuremath{\psi}}
\newunicodechar{ω}{\ensuremath{\omega}}
\newunicodechar{Σ}{\ensuremath{\Sigma}}
% majuscules produites par \MakeUppercase dans les titres (α-aminés -> Α)
\newunicodechar{Α}{\ensuremath{\alpha}}
\newunicodechar{Β}{\ensuremath{\beta}}
\newunicodechar{Γ}{\ensuremath{\Gamma}}
\newunicodechar{Δ}{\ensuremath{\Delta}}
\newunicodechar{Ε}{\ensuremath{\varepsilon}}
\newunicodechar{Θ}{\ensuremath{\Theta}}
\newunicodechar{Λ}{\ensuremath{\Lambda}}
\newunicodechar{Μ}{\ensuremath{\mu}}
\newunicodechar{Τ}{\ensuremath{\tau}}
\newunicodechar{Φ}{\ensuremath{\Phi}}
\newunicodechar{Ψ}{\ensuremath{\Psi}}
\newunicodechar{Ω}{\ensuremath{\Omega}}
\newunicodechar{↦}{\ensuremath{\mapsto}}
\newunicodechar{∈}{\ensuremath{\in}}
\newunicodechar{∉}{\ensuremath{\notin}}
\newunicodechar{∧}{\ensuremath{\wedge}}
\newunicodechar{⇔}{\ensuremath{\Leftrightarrow}}
\newunicodechar{⟺}{\ensuremath{\Longleftrightarrow}}
\newunicodechar{⇒}{\ensuremath{\Rightarrow}}
\newunicodechar{⟹}{\ensuremath{\Longrightarrow}}
\newunicodechar{∘}{\ensuremath{\circ}}
\newunicodechar{∪}{\ensuremath{\cup}}
\newunicodechar{∩}{\ensuremath{\cap}}
\newunicodechar{≡}{\ensuremath{\equiv}}
\newunicodechar{⊂}{\ensuremath{\subset}}
\newunicodechar{⊥}{\ensuremath{\perp}}
\newunicodechar{∃}{\ensuremath{\exists}}
\newunicodechar{∀}{\ensuremath{\forall}}
\newunicodechar{∛}{\ensuremath{\sqrt[3]{\,}}}
\newunicodechar{∜}{\ensuremath{\sqrt[4]{\,}}}
\newunicodechar{⃗}{\ensuremath{{}^{\to}}}
\newunicodechar{ⁿ}{\ifmmode^{n}\else\textsuperscript{\ensuremath{n}}\fi}
\newunicodechar{ˣ}{\ifmmode^{x}\else\textsuperscript{\ensuremath{x}}\fi}
\newunicodechar{ᵇ}{\ifmmode^{b}\else\textsuperscript{\ensuremath{b}}\fi}
\newunicodechar{ʳ}{\ifmmode^{r}\else\textsuperscript{\ensuremath{r}}\fi}
\newunicodechar{ᵘ}{\ifmmode^{u}\else\textsuperscript{\ensuremath{u}}\fi}
\newunicodechar{ʸ}{\ifmmode^{y}\else\textsuperscript{\ensuremath{y}}\fi}
\newunicodechar{ᵃ}{\ifmmode^{a}\else\textsuperscript{\ensuremath{a}}\fi}
\newunicodechar{ᵉ}{\ifmmode^{e}\else\textsuperscript{\ensuremath{e}}\fi}
\newunicodechar{ᵏ}{\ifmmode^{k}\else\textsuperscript{\ensuremath{k}}\fi}
\newunicodechar{ᵖ}{\ifmmode^{p}\else\textsuperscript{\ensuremath{p}}\fi}
\newunicodechar{ᴺ}{\ifmmode^{N}\else\textsuperscript{\ensuremath{N}}\fi}
\newunicodechar{ᶜ}{\ifmmode^{c}\else\textsuperscript{\ensuremath{c}}\fi}
\newunicodechar{ʰ}{\ifmmode^{h}\else\textsuperscript{\ensuremath{h}}\fi}
\newunicodechar{ᵠ}{\ifmmode^{q}\else\textsuperscript{\ensuremath{q}}\fi}
\newunicodechar{ₙ}{\ifmmode_{n}\else\textsubscript{\ensuremath{n}}\fi}
\newunicodechar{ₐ}{\ifmmode_{a}\else\textsubscript{\ensuremath{a}}\fi}
\newunicodechar{ᵢ}{\ifmmode_{i}\else\textsubscript{\ensuremath{i}}\fi}
\newunicodechar{ᵣ}{\ifmmode_{r}\else\textsubscript{\ensuremath{r}}\fi}
\newunicodechar{ₖ}{\ifmmode_{k}\else\textsubscript{\ensuremath{k}}\fi}
\newunicodechar{ᵦ}{\ifmmode_{\beta}\else\textsubscript{\ensuremath{\beta}}\fi}
\newunicodechar{ₓ}{\ifmmode_{x}\else\textsubscript{\ensuremath{x}}\fi}
\newfontfamily\popdisplay{ArchivoBlack-Regular.ttf}[Path=../../../fonts/]
\newfontfamily\poplogo{SpaceGrotesk-Bold.ttf}[Path=../../../fonts/]
\newfontfamily\popsemi{PlusJakartaSans-SemiBold.ttf}[Path=../../../fonts/]
\newfontfamily\popxbold{PlusJakartaSans-ExtraBold.ttf}[Path=../../../fonts/]
\newfontfamily\popxboldls{PlusJakartaSans-ExtraBold.ttf}[Path=../../../fonts/,LetterSpace=3.0]

\setlist[enumerate]{leftmargin=1.7em,itemsep=5pt,topsep=4pt}
\setlist[itemize]{leftmargin=1.7em,itemsep=4pt,topsep=4pt,label={\color{poporange}\textbullet}}
\setlength{\parindent}{0pt}
\setlength{\parskip}{5pt}
\renewcommand{\arraystretch}{1.25}

% pgfplots : style Pop par défaut
\pgfplotsset{
  every axis/.append style={axis line style={popink,line width=1pt},tick style={popink},
    grid style={popline,line width=0.5pt},label style={font=\small},tick label style={font=\footnotesize}},
  popcurve/.style={poporange,line width=1.8pt,no markers,smooth},
}
% Même style disponible dans un \draw TikZ simple (hors pgfplots)
\tikzset{popcurve/.style={poporange,line width=1.8pt}}
\tikzset{popgrid/.style={popline,very thin}}
\colorlet{popcurve}{poporange} % le nom du style est parfois utilisé comme couleur

% ============================================================
% Pill / squiggle
% ============================================================
\newcommand{\poppill}[3]{%
  \tikz[baseline=-0.55ex]\node[draw=popink,line width=1.2pt,rounded corners=2.6pt,%
    fill=#2,inner xsep=6.5pt,inner ysep=3pt]{%
    \popxboldls\fontsize{7}{7}\selectfont\color{#3}\MakeUppercase{#1}};%
}
\newcommand{\popsquiggle}[1]{%
  \tikz[baseline=-0.4ex]{
    \draw[#1,line width=2.6pt,line cap=round]
      plot [smooth] coordinates {(0,0.09) (0.34,0.01) (0.68,0.21) (1.02,0.01) (1.36,0.10)};
  }%
}
% Mot-clé surligné (marqueur orange sous le texte)
\newcommand{\cle}[1]{{\popxbold\color{popdark}#1}}

% ============================================================
% En-tête / pied
% ============================================================
\pagestyle{fancy}
\fancyhf{}
\renewcommand{\headrulewidth}{0pt}
\renewcommand{\footrulewidth}{0pt}
\lhead{\raisebox{-0.15ex}{\poplogo\fontsize{13}{13}\selectfont\color{popink}OnBuch\color{poporange}+}}
\rhead{\poppill{\DOCMATIERE\ \textbullet\ \DOCNIVEAU\ \textbullet\ Leçon \DOCLECON}{white}{popink}}
\lfoot{\popxbold\fontsize{7.6}{7.6}\selectfont\color{popmuted}\MakeUppercase{Cours signé OnBuch+}\ \textbullet\ {\popsemi\DOCTITRE}}
\rfoot{\tikz[baseline=-0.6ex]\node[rounded corners=8.5pt,fill=popink,minimum height=17pt,minimum width=17pt,inner xsep=5pt,inner sep=0]{\popxbold\fontsize{8}{8}\selectfont\color{popcream}\thepage\,/\,\pageref*{LastPage}};}

% ============================================================
% Titres
% ============================================================
\newcommand{\chaptitre}[2]{%
  \begin{tcolorbox}[enhanced,colback=poporange,colframe=popink,boxrule=2.6pt,arc=11pt,
      drop shadow=popink,left=15pt,right=15pt,top=12pt,bottom=11pt,before skip=3pt,after skip=18pt]
    \poppill{#2}{popink}{popcream}\par\vspace{9pt}
    {\raggedright\hyphenpenalty=10000\exhyphenpenalty=10000\popdisplay\fontsize{22}{24}\selectfont\color{popcream}\MakeUppercase{#1}\par}
  \end{tcolorbox}
}
% Section numérotée : pastille orange avec le numéro + titre Archivo
\newcounter{popsec}
\newcommand{\coursec}[1]{%
  \stepcounter{popsec}\setcounter{popsub}{0}%
  \par\vspace{16pt}\noindent
  \begin{minipage}{\linewidth}
  \tikz[baseline=-0.7ex]\node[draw=popink,line width=1.6pt,fill=poporange,rounded corners=4pt,minimum size=22pt,inner sep=0]{\popdisplay\fontsize{12}{12}\selectfont\color{popcream}\thepopsec};%
  \hspace{8pt}{\popdisplay\fontsize{15}{17}\selectfont\color{popink}\MakeUppercase{#1}}\par
  \vspace{4pt}\noindent{\color{poporange}\rule{36pt}{3.6pt}}
  \end{minipage}\par\nopagebreak\vspace{8pt}
}
\newcounter{popsub}
\newcommand{\courssub}[1]{%
  \stepcounter{popsub}%
  \par\vspace{9pt}\noindent{\popxbold\fontsize{11.5}{13}\selectfont\color{popink}{\color{poporange}\thepopsec.\thepopsub}\hspace{6pt}#1}\par\nopagebreak\vspace{4pt}
}

% ============================================================
% Boîtes pédagogiques — même recette : fond blanc, bordure épaisse colorée,
% ombre dure encre, badge plein. Titre optionnel [#1].
% ============================================================
\tcbset{popbase/.style={breakable,enhanced,colback=white,boxrule=2.3pt,arc=9pt,
  shadow={3pt}{-3pt}{0pt}{popink},left=14pt,right=14pt,top=11pt,bottom=11pt,before skip=11pt,after skip=15pt,
  fontupper=\popsemi\fontsize{10.6}{14.5}\selectfont\color{popink},
  fontlower=\popsemi\fontsize{10.6}{14.5}\selectfont\color{popink}}}
% Coche dessinée (le glyphe ✓ n'existe pas dans les polices du document)
\newcommand{\popcheck}[1]{\tikz[baseline=-0.5ex]\draw[#1,line width=1.6pt,line cap=round,line join=round](-0.13,0.0)--(-0.04,-0.09)--(0.13,0.1);}
\providecommand{\coche}{\popcheck{popgreen}}
\providecommand{\resultat}[1]{\par\smallskip\noindent\fcolorbox{popgreen}{popgreenL}{\parbox{\dimexpr\linewidth-2\fboxsep-2\fboxrule\relax}{\textbf{Résultat.} #1}}\par\smallskip}
\newcommand{\popboxhead}[3]{\poppill{#1}{#2}{white}\ifx\relax#3\relax\else\hspace{6pt}{\popxbold\fontsize{10.6}{12}\selectfont\color{popink}#3}\fi\par\vspace{7pt}}

\newtcolorbox{aretenir}[1][]{popbase,colframe=popgreen,before upper={\popboxhead{À retenir}{popgreen}{#1}}}
% Texte en chinois (document de lecture, dialogue, extrait) : cadre
% d'étude. Un titre optionnel (source, auteur) s'affiche dans l'en-tête.
\newtcolorbox{textebox}[1][]{popbase,colframe=popblue,before upper={\popboxhead{文本 · Texte}{popblue}{#1}},after upper={}}
% Marques ✓ / ✗ (forme correcte / fautive) souvent utilisées par les modèles
\providecommand{\cross}{\textcolor{poppink}{\ensuremath{\times}}}
\providecommand{\xmark}{\cross}\providecommand{\crossmark}{\cross}\providecommand{\cmark}{\ensuremath{\checkmark}}
% Ligne de réponse / justification à compléter (inventée par les modèles)
\providecommand{\justification}[1]{\par\noindent\textit{Justification~:} \dotfill\par}
% \textipa (package tipa non chargé) : la police ne couvre pas l'API -> simple passage
\providecommand{\textipa}[1]{#1}
% Soulignement (ulem non chargé) : \uline, \ul
\providecommand{\uline}[1]{\underline{#1}}\providecommand{\ul}[1]{\underline{#1}}
% \glossaire{\item ...} inventée par les modèles : liste de glossaire
\providecommand{\glossaire}[1]{\begin{itemize}#1\end{itemize}}
% \alert (beamer) inventée par les modèles : texte mis en évidence
\providecommand{\alert}[1]{\textbf{\textcolor{poppink}{#1}}}
% variantes ASCII d'ordinaux écrites en macro
\providecommand{\iere}{\textsuperscript{re}}\providecommand{\ieme}{\textsuperscript{e}}\providecommand{\ere}{\textsuperscript{re}}\providecommand{\eme}{\textsuperscript{e}}\providecommand{\er}{\textsuperscript{er}}\providecommand{\ie}{\textsuperscript{e}}
% Ordinaux français écrits en macro par les modèles : 1\ière, 3\ième
\newcommand{\ière}{\textsuperscript{re}}\newcommand{\ième}{\textsuperscript{e}}
% \exemple (inventée par les modèles) : étiquette « Exemple : »
\providecommand{\exemple}{\textbf{Exemple~:}\ }
% \square utilisable aussi hors mode math (cases à cocher)
\AtBeginDocument{\let\squareMath\square\renewcommand{\square}{\ensuremath{\squareMath}}}
% Mot ou phrase en chinois à l'intérieur d'un paragraphe français
\newcommand{\zh}[1]{\ifmmode\text{#1}\else{#1}\fi}
\let\eng\zh \let\esp\zh \let\all\zh % alias tolérés
% flèches d'intonation inventées par les modèles
\providecommand{\textsearrow}{}\renewcommand{\textsearrow}{\ensuremath{\searrow}}\providecommand{\textnearrow}{}\renewcommand{\textnearrow}{\ensuremath{\nearrow}}
% \fr{..} : mot français (inventé par les modèles)
\providecommand{\fr}[1]{\textit{#1}}
% zhbox : encadré de texte chinois inventé par les modèles
\newenvironment{zhbox}{\begin{quote}}{\end{quote}}
% \courssubsub : titre de niveau 3 inventé par les modèles
\providecommand{\courssubsub}[1]{\par\medskip\noindent\textbf{#1}\par\smallskip}
% \zhbf : texte chinois en gras inventé par les modèles
\providecommand{\zhbf}[1]{\textbf{#1}}
% \vs : « versus » inventé par les modèles
\providecommand{\vs}{\textit{vs}\ }
% \blank : case à compléter inventée par les modèles
\providecommand{\blank}[1][]{\underline{\hspace{1.6em}}}
\newtcolorbox{exemplebox}[1][]{popbase,colframe=poporange,before upper={\popboxhead{Exemple}{poporange}{#1}}}
\newtcolorbox{definition}[1][]{popbase,colframe=popblue,before upper={\popboxhead{Définition}{popblue}{#1}}}
\newtcolorbox{propriete}[1][]{popbase,colframe=poppurple,before upper={\popboxhead{Propriété}{poppurple}{#1}}}
\newtcolorbox{methode}[1][]{popbase,colframe=popgold,before upper={\popboxhead{Méthode}{popgold}{#1}}}
\newtcolorbox{attention}[1][]{popbase,colframe=poppink,before upper={\popboxhead{Attention}{poppink}{#1}}}
\newtcolorbox{experience}[1][]{popbase,colframe=popblue,colback=popblueL,before upper={\popboxhead{Expérience}{popblue}{#1}}}
% « Le savais-tu ? » : carte violette pleine (contexte, culture, Cameroun)
\newtcolorbox{savaistu}[1][]{popbase,colback=poppurple,colframe=popink,
  fontupper=\popsemi\fontsize{10.4}{14.2}\selectfont\color{white},
  before upper={\poppill{Le savais-tu\,?}{white}{poppurple}\ifx\relax#1\relax\else\hspace{6pt}{\popxbold\color{white}#1}\fi\par\vspace{7pt}}}
% Exercice résolu : énoncé en haut, \tcblower puis la solution sur fond crème
\newcounter{popexo}\newcounter{popexr}
\newtcolorbox{exoresolu}[1][]{popbase,colframe=poporange,colbacklower=popcream,
  segmentation style={popink,line width=1.2pt,dashed},
  before upper={\stepcounter{popexr}\popboxhead{Exercice résolu \thepopexr}{poporange}{#1}},
  before lower={\poppill{Solution}{popink}{popcream}\par\vspace{6pt}}}
% Exercice (non résolu en place) — la correction est dans la partie Corrigés
\newtcolorbox{exercice}[1][]{popbase,colframe=popink,
  before upper={\stepcounter{popexo}\popboxhead{Exercice \thepopexo}{popink}{#1}}}
% Corrigé
\newtcolorbox{corrige}[1][]{popbase,colframe=popgreen,colback=popgreenL,
  before upper={\popboxhead{Corrigé}{popgreen}{#1}}}
% Objectifs / prérequis / bilan
\newtcolorbox{objectifs}{popbase,colframe=popink,colback=poporangeL,
  before upper={\popboxhead{Objectifs de la leçon}{popink}{}}}
\newtcolorbox{prerequis}{popbase,colframe=popink,colback=popblueL,
  before upper={\popboxhead{Prérequis}{popblue}{}}}
\newtcolorbox{bilan}{popbase,colframe=popink,colback=white,boxrule=2.8pt,
  before upper={\popboxhead{Fiche bilan}{poporange}{L'essentiel à connaître pour l'examen}}}
% Figure encadrée avec légende
\newtcolorbox{popfigure}[1][]{enhanced,breakable=false,colback=white,colframe=popline,boxrule=1.4pt,arc=8pt,
  left=10pt,right=10pt,top=10pt,bottom=8pt,before skip=10pt,after skip=12pt,halign=center,
  lower separated=false,halign lower=center,
  fontlower=\popsemi\fontsize{9}{11.5}\selectfont\color{popmuted},
  #1}
\newcommand{\legende}[1]{\par\vspace{6pt}{\popsemi\fontsize{9}{11.5}\selectfont\color{popmuted}#1}}

% ============================================================
% Signature OnBuch+
% ============================================================
\newcommand{\poplogoinline}[1]{{\poplogo\fontsize{#1}{#1}\selectfont\color{popink}OnBuch\color{poporange}+}}
\newcommand{\signatureonbuch}{%
  \begin{tcolorbox}[enhanced,colback=popink,colframe=popink,boxrule=2.2pt,arc=10pt,
      drop shadow=poporange,left=14pt,right=14pt,top=10pt,bottom=10pt,before skip=0pt,after skip=0pt]
    \tikz[baseline=-0.7ex]\node[circle,fill=popgreen,draw=popcream,line width=1.3pt,minimum size=22pt,inner sep=0]{\popcheck{white}};%
    \hspace{9pt}\begin{minipage}[c]{0.62\linewidth}
      {\popxbold\fontsize{12}{13}\selectfont\color{popcream}Signé\ \textbullet\ {\poplogo OnBuch\color{poporange}+}}\par\vspace{2pt}
      {\raggedright\popsemi\fontsize{8}{10.5}\selectfont\color{popline}\MakeUppercase{Équipe pédagogique OnBuch+}\\\MakeUppercase{Conforme au programme officiel MINESEC}\par}
    \end{minipage}\hfill
    {\poplogo\fontsize{22}{22}\selectfont\color{popcream}OnBuch\color{poporange}+}
  \end{tcolorbox}%
}

% ============================================================
% Page de garde
% #1 titre · #2 matière · #3 niveau · #4 module · #5 n° leçon · #6 durée ·
% #7 description / objectifs (texte ou itemize)
% ============================================================
\newcommand{\pagegarde}[7]{%
  \thispagestyle{empty}
  \newgeometry{a4paper,margin=1.7cm,top=1.6cm,bottom=1.4cm}%
  \begin{tikzpicture}[remember picture,overlay]
    \fill[poporange!18] ([xshift=-3.2cm,yshift=-3.6cm]current page.north east) circle (4.6cm);
    \fill[poppurple!14] ([xshift=2.4cm,yshift=7.6cm]current page.south west) circle (3.8cm);
  \end{tikzpicture}%
  {\poplogo\fontsize{30}{30}\selectfont\color{popink}OnBuch\color{poporange}+}\hfill
  \raisebox{0.6ex}{\poppill{Cours signé \textbullet\ Premium}{popink}{popcream}}\par
  \vspace{1pt}\popsquiggle{poppurple}\par
  \vspace{8pt}
  \begin{tcolorbox}[enhanced,colback=poporange,colframe=popink,boxrule=2.8pt,arc=13pt,
      drop shadow=popink,left=18pt,right=18pt,top=12pt,bottom=13pt,after skip=10pt]
    \poppill{#2 \textbullet\ #3}{popink}{popcream}\hspace{5pt}\poppill{#4}{white}{popink}\par\vspace{8pt}
    {\popdisplay\fontsize{14}{14}\selectfont\color{popink}LEÇON #5}\par\vspace{4pt}
    {\raggedright\hyphenpenalty=10000\exhyphenpenalty=10000\popdisplay\fontsize{25}{27}\selectfont\color{popcream}\MakeUppercase{#1}\par}
    \vspace{8pt}
    {\popxbold\fontsize{9.5}{9.5}\selectfont\color{popink}\MakeUppercase{Durée conseillée : #6}}
  \end{tcolorbox}
  \begin{tcolorbox}[enhanced,colback=white,colframe=popink,boxrule=2.2pt,arc=10pt,
      drop shadow=popink,left=16pt,right=16pt,top=10pt,bottom=10pt,after skip=8pt]
    \poppill{Dans cette leçon}{poppurple}{white}\par\vspace{5pt}
    \popsemi\fontsize{9.3}{12.6}\selectfont\color{popink}#7
  \end{tcolorbox}
  \noindent\begin{minipage}[t]{0.487\linewidth}
    \begin{tcolorbox}[enhanced,colback=popgreen,colframe=popink,boxrule=2.2pt,arc=10pt,
        drop shadow=popink,left=11pt,right=11pt,top=10pt,bottom=10pt,height=3.1cm,valign=center]
      \tcbox[colback=white,colframe=popink,boxrule=1.3pt,arc=5pt,boxsep=0pt,nobeforeafter,
        left=3pt,right=3pt,top=3pt,bottom=3pt,valign=center]{\qrcode[height=1.9cm]{https://play.google.com/store/apps/details?id=com.luvvix.onbuch&hl=fr}}%
      \hspace{8pt}\begin{minipage}[c]{0.5\linewidth}
        {\popdisplay\fontsize{11}{12.5}\selectfont\color{white}\MakeUppercase{Télécharge OnBuch}}\par\vspace{4pt}
        {\raggedright\popsemi\fontsize{8.4}{11}\selectfont\color{white}Gratuit \textbullet\ Google Play\\Tuteur IA, annales, concours, résultats\par}
      \end{minipage}
    \end{tcolorbox}
  \end{minipage}\hfill
  \begin{minipage}[t]{0.487\linewidth}
    \begin{tcolorbox}[enhanced,colback=poppurple,colframe=popink,boxrule=2.2pt,arc=10pt,
        drop shadow=popink,left=11pt,right=11pt,top=10pt,bottom=10pt,height=3.1cm,valign=center]
      \tikz[baseline=-0.7ex]\node[circle,fill=white,draw=popink,line width=1.3pt,minimum size=1.6cm,inner sep=0]{\popdisplay\fontsize{24}{24}\selectfont\color{poppurple}+};%
      \hspace{8pt}\begin{minipage}[c]{0.6\linewidth}
        {\popdisplay\fontsize{11}{12.5}\selectfont\color{white}\MakeUppercase{Passe à OnBuch+}}\par\vspace{4pt}
        {\raggedright\popsemi\fontsize{8.4}{11}\selectfont\color{white}Tous les cours signés, Tuteur IA illimité, fiches PDF — onglet OnBuch+ de l'app\par}
      \end{minipage}
    \end{tcolorbox}
  \end{minipage}
  \vspace{8pt}
  \popcontenu
  \vfill
  \signatureonbuch
  \restoregeometry
  \newpage
}

% Rangée « Ce cours contient » (page de garde)
\newcommand{\poptile}[3]{% #1 couleur #2 gros texte #3 légende
  \begin{tcolorbox}[enhanced,colback=#1,colframe=popink,boxrule=2pt,arc=9pt,shadow={2.5pt}{-2.5pt}{0pt}{popink},
      left=6pt,right=6pt,top=9pt,bottom=9pt,halign=center,valign=center,height=1.75cm,width=\linewidth,nobeforeafter]
    {\popdisplay\fontsize{12.5}{13}\selectfont\color{white}\MakeUppercase{#2}}\par\vspace{3pt}
    {\centering\popsemi\fontsize{7.8}{9.5}\selectfont\color{white}#3\par}
  \end{tcolorbox}}
\newcommand{\popcontenu}{%
  \par\noindent{\popxboldls\fontsize{7.5}{7.5}\selectfont\color{popmuted}CE COURS SIGNÉ CONTIENT}\par\vspace{7pt}
  \noindent\begin{minipage}[t]{0.235\linewidth}\poptile{popblue}{Cours}{complet, illustré, pas à pas}\end{minipage}\hfill
  \begin{minipage}[t]{0.235\linewidth}\poptile{poporange}{Méthodes}{et exercices résolus}\end{minipage}\hfill
  \begin{minipage}[t]{0.235\linewidth}\poptile{poppink}{Exercices}{type examen + corrigés}\end{minipage}\hfill
  \begin{minipage}[t]{0.235\linewidth}\poptile{popgreen}{Fiche bilan}{l'essentiel à retenir}\end{minipage}\par}

% Sommaire
\newtcolorbox{sommaire}{enhanced,colback=white,colframe=popink,boxrule=2.2pt,arc=10pt,
  drop shadow=popink,left=18pt,right=18pt,top=13pt,bottom=13pt,before skip=6pt,after skip=20pt,
  before upper={\poppill{Sommaire}{poppurple}{white}\par\vspace{9pt}\popsemi\fontsize{10.6}{14.5}\selectfont\color{popink}}}

% Signature de fin de cours — poussée tout en bas de la dernière page
\newcommand{\findecours}{%
  \par\vfill\nopagebreak
  \begin{center}\popsquiggle{poporange}\end{center}\vspace{4pt}
  \signatureonbuch
}
\AtBeginDocument{\def\llbracket{\mathopen{[\mkern-3.5mu[}}\def\rrbracket{\mathclose{]\mkern-3.5mu]}}} % intervalles d'entiers (glyphe ⟦ absent de la police)
% Glyphes absents de Fira Math : redéfinis à partir de glyphes disponibles
\AtBeginDocument{%
  \def\setminus{\mathbin{\backslash}}\let\smallsetminus\setminus
  \def\vdots{\mathord{\vbox{\baselineskip3.2pt\lineskiplimit0pt\kern4pt\hbox{.}\hbox{.}\hbox{.}}}}%
  \def\ddots{\mathinner{\mkern1mu\raise6.5pt\hbox{.}\mkern2mu\raise3.6pt\hbox{.}\mkern2mu\raise0.7pt\hbox{.}\mkern1mu}}%
  \def\triangle{\mathord{\scalebox{0.9}{$\symup{\Delta}$}}}%
  \def\nabla{\mathord{\raisebox{\depth}{\scalebox{1}[-1]{$\symup{\Delta}$}}}}%
  \def\checkmark{\popcheck{popdark}}%
  \def\star{\mathbin{\ast}}\def\bigstar{\mathord{\ast}}%
  \def\diamond{\mathbin{\scalebox{0.6}{$\Diamond$}}}%
  \def\bigcup{\mathop{\mathchoice{\raisebox{-0.3em}{\scalebox{1.7}{$\cup$}}}{\scalebox{1.25}{$\cup$}}{\cup}{\cup}}\displaylimits}%
  \def\bigcap{\mathop{\mathchoice{\raisebox{-0.3em}{\scalebox{1.7}{$\cap$}}}{\scalebox{1.25}{$\cap$}}{\cap}{\cap}}\displaylimits}%
  \def\lesseqqgtr{\mathrel{\lessgtr}}\def\bigoplus{\mathop{\oplus}\displaylimits}%
}
\providecommand{\ssi}{si et seulement si} % abréviation fréquente
\AtBeginDocument{\providecommand{\wideparen}{\overparen}} % arc de cercle

\begin{document}

\pagegarde{Leçon 31 — Grammaire : Les connecteurs 首先… 然后……其次……最后… ; 既…..又 ; 只有….才…. ; 只要……就…..}{Chinois}{Tle A}{Module 4 : Activités économiques et monde du travail}{ZH31}{2 h}{Cette leçon de grammaire présente quatre structures de connecteurs essentielles du chinois : l'enchaînement chronologique 首先…然后…其次…最后…, la cumulation 既…又…, la condition nécessaire 只有…才… et la condition suffisante 只要…就…. À travers l'observation d'exemples, des schémas de structure et des exercices de transformation, l'élève apprend à construire des phrases complexes claires et à éviter les pièges typiques des francophones.
\vspace{4pt}
\begin{multicols}{2}\small
\begin{itemize}
\item À la fin de cette leçon, je sais identifier les quatre structures de connecteurs dans un texte
\item À la fin de cette leçon, je sais construire un récit d'étapes avec 首先…然后…其次…最后…
\item À la fin de cette leçon, je sais cumuler deux qualités ou actions avec 既…又…
\item À la fin de cette leçon, je sais exprimer une condition nécessaire avec 只有…才…
\item À la fin de cette leçon, je sais exprimer une condition suffisante avec 只要…就…
\item À la fin de cette leçon, je sais distinguer 只有…才… et 只要…就… selon le sens voulu
\end{itemize}\end{multicols}}

\begin{sommaire}
\begin{multicols}{2}
\begin{enumerate}
\item Observation et découverte guidée
\item 首先…然后…其次…最后… : organiser une suite d'étapes
\item 既…又… : cumuler deux caractéristiques
\item 只有…才… : la condition nécessaire
\item 只要…就… : la condition suffisante et comparaison des deux conditions
\item Synthèse, entraînement et production
\item Activité d'intégration
\item Méthodes et exercices résolus
\item Exercices
\item Corrigés des exercices
\item Fiche bilan
\end{enumerate}
\end{multicols}
\end{sommaire}

\begin{objectifs}
\begin{itemize}
\item À la fin de cette leçon, je sais identifier les quatre structures de connecteurs dans un texte
\item À la fin de cette leçon, je sais construire un récit d'étapes avec 首先…然后…其次…最后…
\item À la fin de cette leçon, je sais cumuler deux qualités ou actions avec 既…又…
\item À la fin de cette leçon, je sais exprimer une condition nécessaire avec 只有…才…
\item À la fin de cette leçon, je sais exprimer une condition suffisante avec 只要…就…
\item À la fin de cette leçon, je sais distinguer 只有…才… et 只要…就… selon le sens voulu
\item À la fin de cette leçon, je sais placer correctement les connecteurs par rapport au sujet
\item À la fin de cette leçon, je sais rédiger un court paragraphe organisé utilisant au moins trois de ces structures
\end{itemize}
\end{objectifs}

\begin{prerequis}
\begin{itemize}
\item L'ordre de base de la phrase chinoise Sujet-Verbe-Complément (1re)
\item Les adverbes de liaison simples 先, 再, 然后 (1re)
\item La structure 因为…所以… et 虽然…但是… (début de Tle)
\item La place des adverbes (都, 也, 就) avant le verbe
\item Le vocabulaire du monde du travail et des activités économiques (Module 4)
\end{itemize}
\end{prerequis}

\newpage

\coursec{Observation et découverte guidée}

\courssub{Situation d'accroche et problématique}

Imagine que tu dois expliquer à un correspondant chinois comment se déroule la fête de la jeunesse à Yaoundé : d'abord le défilé, ensuite les discours, puis les matchs, enfin la remise des prix. Comment enchaîner ces étapes en chinois ? Et si tu veux dire que le Cameroun est à la fois riche en ressources et en cultures, ou que seule la pratique régulière permet de réussir au Bac ? Il te faut des connecteurs précis. C'est exactement ce que cette leçon t'offre : quatre outils pour organiser, cumuler et conditionner tes idées comme un locuteur authentique.

\textbf{Problématique :} comment le chinois relie-t-il les idées entre elles ? Pourquoi certains mots de liaison doivent-ils toujours aller par deux ?

\courssub{Lecture à voix haute : quatre phrases modèles}

Lis chaque phrase à voix haute, en respectant les tons. Observe les mots en gras : ce sont les connecteurs que nous allons étudier.

\begin{textebox}[Phrase modèle 1 : organiser des étapes]
首先，我们参观工厂，然后和工人谈话，其次了解生产过程，最后写报告。\\
\emph{Shǒuxiān, wǒmen cānguān gōngchǎng, ránhòu hé gōngrén tánhuà, qícì liǎojiě shēngchǎn guòchéng, zuìhòu xiě bàogào.}\\
« D'abord, nous visitons l'usine, ensuite nous parlons avec les ouvriers, puis nous découvrons le processus de production, enfin nous rédigeons un rapport. »
\end{textebox}

\begin{exemplebox}[Phrase modèle 2 : cumuler deux qualités]
\zh{这家公司既大又现代化。}\\
\emph{Zhè jiā gōngsī jì dà yòu xiàndàihuà.}\\
« Cette entreprise est à la fois grande et moderne. »
\end{exemplebox}

\begin{exemplebox}[Phrase modèle 3 : la condition nécessaire]
\zh{只有努力学习，才能通过考试。}\\
\emph{Zhǐyǒu nǔlì xuéxí, cái néng tōngguò kǎoshì.}\\
« Ce n'est qu'en travaillant dur qu'on peut réussir l'examen. »
\end{exemplebox}

\begin{exemplebox}[Phrase modèle 4 : la condition suffisante]
\zh{只要你努力，就会成功。}\\
\emph{Zhǐyào nǐ nǔlì, jiù huì chénggōng.}\\
« Tant que tu fais des efforts, tu réussiras. »
\end{exemplebox}

\courssub{Questions de découverte guidée}

Observe les quatre phrases et réponds mentalement aux questions suivantes avant de lire les réponses.

\begin{enumerate}
\item Dans la phrase 1, où se placent les mots \zh{首先}, \zh{然后}, \zh{其次}, \zh{最后} ? (Réponse : en tête de chaque proposition, souvent suivis d'une virgule chinoise \zh{，}.)
\item Dans la phrase 2, quels mots entourent les deux adjectifs \zh{大} et \zh{现代化} ? (Réponse : \zh{既} avant le premier, \zh{又} avant le second.)
\item Dans la phrase 3, quel mot répond à \zh{只有} dans la deuxième proposition ? (Réponse : \zh{才}, placé avant le verbe \zh{能}.)
\item Dans la phrase 4, quel mot répond à \zh{只要} ? (Réponse : \zh{就}, placé avant le verbe \zh{会}.)
\item Peut-on supprimer le deuxième mot de chaque paire ? (Réponse : non, la phrase devient incomplète ou incorrecte.)
\end{enumerate}

\begin{methode}[Découvrir une structure grammaticale chinoise]
\begin{enumerate}
\item Lis la phrase entière et repère les mots qui se \cle{répètent} ou qui se \cle{répondent} d'une proposition à l'autre.
\item Souligne ces mots et note leur position : début de phrase, avant le verbe, avant l'adjectif.
\item Compare avec le français : quel mot français correspond à chaque élément de la paire ?
\item Vérifie si le deuxième mot de la paire est obligatoire en le supprimant mentalement : si la phrase sonne inachevée, c'est un \cle{binôme obligatoire}.
\item Reformule une phrase de ta propre invention avec la même structure pour confirmer ta compréhension.
\end{enumerate}
\end{methode}

\courssub{Mise en évidence : des connecteurs qui fonctionnent en binômes}

La grande découverte de cette leçon est la suivante : en chinois, beaucoup de connecteurs logiques fonctionnent \cle{par paires inséparables}. Là où le français utilise un seul mot (« ensuite », « seulement »), le chinois exige souvent deux mots qui se répondent, comme deux parenthèses autour de l'idée.

\begin{aretenir}[Les quatre structures de la leçon]
\begin{itemize}
\item \zh{首先…然后…其次…最后…} : présenter des \cle{étapes chronologiques} (« d'abord… ensuite… puis… enfin… »).
\item \zh{既…又…} : \cle{cumuler} deux caractéristiques (« à la fois… et… »).
\item \zh{只有…才…} : exprimer la \cle{condition nécessaire} (« ce n'est que… que… »).
\item \zh{只要…就…} : exprimer la \cle{condition suffisante} (« tant que…, … »).
\end{itemize}
\end{aretenir}

\begin{center}
\begin{tabularx}{\linewidth}{|X|X|X|X|}
\hline
\rowcolor{popblueL} Structure & Fonction logique & Équivalent français & Exemple-clé \\
\hline
首先…然后…其次…最后… & chronologie, étapes & d'abord… ensuite… puis… enfin… & 首先参观，然后谈话 \\
\hline
既…又… & cumul de qualités & à la fois… et… & 既大又现代化 \\
\hline
只有…才… & condition nécessaire & seulement si… alors… & 只有努力，才能成功 \\
\hline
只要…就… & condition suffisante & tant que…, … & 只要努力，就会成功 \\
\hline
\end{tabularx}
\end{center}

\begin{popfigure}
\begin{tikzpicture}[mindmap, grow cyclic, every node/.style={concept, align=center, font=\small},
root concept/.append style={concept color=popblue, font=\bfseries},
level 1/.append style={level distance=3.4cm, sibling angle=90}]
\node[root concept, align=center]{Connecteurs\\ 连接词}
  child[concept color=poporange]{ node[align=center]{Chronologie\\ 首先…然后…\\ 其次…最后…} }
  child[concept color=popgreen]{ node[align=center]{Cumul\\ 既…又…} }
  child[concept color=poppurple]{ node[align=center]{Condition\\ nécessaire\\ 只有…才…} }
  child[concept color=poppink]{ node[align=center]{Condition\\ suffisante\\ 只要…就…} };
\end{tikzpicture}
\legende{Carte mentale des quatre connecteurs étudiés dans la leçon 31.}
\end{popfigure}

\begin{definition}[Précision sur l'ordre des étapes]
Dans une énumération soignée (à l'écrit, dans un exposé), l'ordre le plus courant est \zh{首先…其次…再次…最后…} (« d'abord… ensuite… de plus… enfin… »). Le mot \zh{然后} (« ensuite ») est plus oral et peut s'intercaler librement : \zh{首先…然后…最后…}. À l'examen, les deux schémas sont acceptés, mais veille à toujours terminer par \zh{最后}.
\end{definition}

\courssub{Premier entraînement : repérer les binômes}

\begin{exoresolu}[Repérage des connecteurs]
Énoncé : souligne les connecteurs dans la phrase suivante et indique leur fonction : \zh{我们只有学好汉语，才能去中国工作。}\\
\emph{Wǒmen zhǐyǒu xuéhǎo Hànyǔ, cái néng qù Zhōngguó gōngzuò.}
\tcblower
Solution détaillée : les connecteurs sont \zh{只有} (début de la première proposition) et \zh{才} (début de la deuxième proposition, avant le verbe \zh{能}). Il s'agit du binôme de la \cle{condition nécessaire}. Traduction : « Ce n'est qu'en apprenant bien le chinois que nous pourrons aller travailler en Chine. » On vérifie : si l'on supprime \zh{才}, la phrase \zh{我们只有学好汉语，能去中国工作} devient maladroite et ambiguë ; le binôme est donc obligatoire.
\end{exoresolu}

\begin{attention}[Erreurs fréquentes des francophones]
\begin{itemize}
\item Ne traduis jamais « seulement » par \zh{只有} seul : en français le mot suffit, en chinois il appelle obligatoirement \zh{才} dans la proposition suivante.
\item N'utilise pas \zh{和} (et) pour cumuler deux adjectifs : on ne dit pas \zh{公司大和现代化}, mais \zh{公司既大又现代化}. Rappel : \zh{和} ne relie que des noms ou des groupes nominaux, jamais des adjectifs ni des propositions.
\item Respecte l'ordre : \zh{才} et \zh{就} se placent \textbf{avant le verbe} de la deuxième proposition, jamais après.
\item Ne confonds pas les tons : \zh{既} (jì, 4\ieme{} ton) dans \zh{既…又…} n'a rien à voir avec \zh{即} (jí, 2\ieme{} ton, « c'est-à-dire »). À l'écrit, vérifie bien le caractère.
\end{itemize}
\end{attention}

\begin{savaistu}[Des connecteurs inséparables]
En chinois, contrairement au français, beaucoup de connecteurs logiques sont \cle{inséparables} : dire \zh{只有} sans \zh{才} laisse la phrase inachevée, comme une parenthèse ouverte jamais refermée. C'est une marque de la pensée chinoise : la relation logique est portée par \textbf{deux} mots qui encadrent l'idée, et non par un seul. Le locuteur expérimenté « entend » la paire entière dès le premier mot prononcé. Tu retrouveras ce principe avec d'autres binômes déjà vus ou à venir : \zh{因为…所以…} (parce que… donc…), \zh{虽然…但是…} (bien que… mais…).
\end{savaistu}

\begin{exercice}[À toi de jouer]
\begin{enumerate}
\item Souligne les binômes de connecteurs dans : \zh{这个工厂既大又干净。} et \zh{只要天气好，我们就去参观。}
\item Traduis en chinois : « D'abord nous écoutons, ensuite nous répétons. »
\item Complète : \zh{只有每天练习，} \_\_\_\_ \zh{说好汉语。}
\end{enumerate}
\end{exercice}

\begin{corrige}[Exercice « À toi de jouer »]
\begin{enumerate}
\item \zh{既…又…} (cumul : « à la fois grand et propre ») ; \zh{只要…就…} (condition suffisante : « tant qu'il fait beau, nous irons visiter »).
\item \zh{首先，我们听，然后我们重复。} \emph{Shǒuxiān, wǒmen tīng, ránhòu wǒmen chóngfù.}
\item \zh{只有每天练习，才能说好汉语。} \emph{Zhǐyǒu měitiān liànxí, cái néng shuōhǎo Hànyǔ.} « Ce n'est qu'en pratiquant chaque jour qu'on peut bien parler chinois. »
\end{enumerate}
\end{corrige}

Tu as maintenant observé les quatre structures en action. Dans les sections suivantes, nous étudierons chacune en détail : sa forme exacte, ses emplois, ses exceptions et ses pièges.

\coursec{首先…然后…其次…最后… : organiser une suite d'étapes}

Dans la section précédente, nous avons observé un jeune diplômé de Yaoundé raconter sa recherche d'emploi : il décrivait ses démarches \emph{dans un ordre précis}, étape après étape. En français, nous disons « d'abord… ensuite… puis… enfin… ». Le chinois possède un système de connecteurs parfaitement parallèle, mais avec des règles de placement et des nuances de registre qu'il faut maîtriser : c'est l'objet de cette section.

\courssub{Observation : repérer les connecteurs dans un texte}

Lisons d'abord ce court texte, extrait du discours d'un directeur d'entreprise à Douala qui explique comment recruter un employé.

\begin{textebox}[Comment recruter un employé ?]
我们公司招聘员工的时候，首先要发布招聘信息，然后看申请人的简历，其次安排面试，最后决定录用谁。\\
Wǒmen gōngsī zhāopìn yuángōng de shíhou, shǒuxiān yào fābù zhāopìn xìnxī, ránhòu kàn shēnqǐngrén de jiǎnlì, qícì ānpái miànshì, zuìhòu juédìng lùyòng shéi.\\
« Quand notre entreprise recrute du personnel, il faut d'abord publier l'annonce, ensuite examiner les CV des candidats, puis organiser les entretiens, enfin décider qui sera embauché. »
\end{textebox}

\textbf{Questions d'observation :}
\begin{enumerate}
\item Combien d'étapes le directeur décrit-il ? Soulignez les quatre mots qui les introduisent.
\item Où se placent ces mots : avant ou après le sujet de chaque proposition ?
\item Chaque connecteur est-il suivi d'une phrase complète (sujet + verbe) ?
\end{enumerate}

On constate que le texte comporte \textbf{quatre étapes}, introduites par \zh{首先} (d'abord), \zh{然后} (ensuite), \zh{其次} (en second lieu) et \zh{最后} (enfin). Chacun de ces connecteurs ouvre une \textbf{proposition complète}, séparée des autres par une virgule chinoise \zh{，}.

\courssub{La règle : schéma, forme et emploi}

\begin{propriete}[La chaîne 首先…然后…其次…最后…]
Pour organiser une suite d'étapes chronologiques ou d'arguments ordonnés, le chinois utilise la chaîne de connecteurs :
\begin{center}
\zh{首先} + phrase 1，\zh{然后} + phrase 2，\zh{其次} + phrase 3，\zh{最后} + phrase 4。
\end{center}
\begin{itemize}
\item Chaque connecteur introduit une \textbf{proposition complète} (sujet + verbe), jamais un simple mot isolé.
\item \zh{首先} se place \textbf{en tête de phrase} ou \textbf{juste après le sujet} : \zh{首先我要感谢老师} ou \zh{我首先要感谢老师} sont tous deux corrects.
\item \zh{然后}, \zh{其次} et \zh{最后} se placent \textbf{en tête de proposition}, après la virgule.
\item Les propositions sont séparées par la virgule chinoise \zh{，} ; seule la dernière prend le point final.
\end{itemize}
\end{propriete}

\begin{definition}[Valeur de chaque connecteur]
\begin{itemize}
\item \zh{首先} (shǒuxiān) : « d'abord, en premier lieu » — ouvre la série.
\item \zh{然后} (ránhòu) : « ensuite, puis » — connecteur neutre, le plus courant à l'oral.
\item \zh{其次} (qícì) : « ensuite, en second lieu » — plus \textbf{soutenu} que \zh{然后} ; il sert surtout à ordonner des \textbf{arguments} (dans un exposé, un discours, une dissertation).
\item \zh{最后} (zuìhòu) : « enfin, finalement » — ferme la série.
\end{itemize}
Remarque : à l'oral, on rencontre aussi \zh{接着} (jiēzhe, « ensuite, aussitôt après »), qui insiste sur l'enchaînement immédiat de deux actions ; à ce niveau, retenez simplement qu'il existe et que \zh{然后} reste le connecteur neutre de référence.
\end{definition}

\textbf{Emplois principaux :}
\begin{enumerate}
\item \textbf{Récit de processus} : décrire les étapes d'une démarche (chercher un emploi, préparer le ndolé, s'inscrire à l'université).
\item \textbf{Exposé et discours} : ordonner des arguments (d'abord… en second lieu… enfin…).
\item \textbf{Mode d'emploi} : expliquer comment faire quelque chose, étape par étape.
\end{enumerate}

\courssub{Frise des étapes : l'exemple de l'entretien d'embauche}

\begin{popfigure}
\begin{center}
\begin{tikzpicture}[font=\small]
\node[draw=popblue, fill=popblueL, rounded corners, minimum width=2.6cm, minimum height=1.1cm, align=center] (e1) at (0,0) {\zh{首先}\\ préparer le CV};
\node[draw=poporange, fill=poporangeL, rounded corners, minimum width=2.6cm, minimum height=1.1cm, align=center, right=1.1cm of e1] (e2) {\zh{然后}\\ envoyer la\\ candidature};
\node[draw=popgreen, fill=popgreenL, rounded corners, minimum width=2.6cm, minimum height=1.1cm, align=center, right=1.1cm of e2] (e3) {\zh{其次}\\ passer\\ l'entretien};
\node[draw=poppurple, fill=poppurpleL, rounded corners, minimum width=2.6cm, minimum height=1.1cm, align=center, right=1.1cm of e3] (e4) {\zh{最后}\\ attendre la\\ réponse};
\draw[-{Stealth}, very thick, popdark] (e1) -- (e2);
\draw[-{Stealth}, very thick, popdark] (e2) -- (e3);
\draw[-{Stealth}, very thick, popdark] (e3) -- (e4);
\end{tikzpicture}
\end{center}
\legende{Frise des quatre étapes d'une candidature : 首先 → 然后 → 其次 → 最后}
\end{popfigure}

Cette frise correspond à la phrase modèle de la leçon :

\zh{找工作的时候，首先要准备好简历，然后参加面试，最后等公司的通知。}

\emph{Zhǎo gōngzuò de shíhou, shǒuxiān yào zhǔnbèi hǎo jiǎnlì, ránhòu cānjiā miànshì, zuìhòu děng gōngsī de tōngzhī.}

« Pour chercher un emploi, il faut d'abord préparer son CV, ensuite passer l'entretien, enfin attendre la réponse de l'entreprise. »

Remarquez que cette phrase n'utilise que trois connecteurs : \zh{其次} est \textbf{facultatif}. C'est la première souplesse de la structure.

\courssub{Souplesses et exceptions de la structure}

\begin{propriete}[Variantes autorisées]
\begin{enumerate}
\item \textbf{Omission de} \zh{其次} : dans un récit de trois étapes, on enchaîne simplement \zh{首先…然后…最后…}。
\item \textbf{Multiplication de} \zh{然后} : dans un récit long (plus de quatre étapes), on peut répéter \zh{然后} autant de fois que nécessaire : \zh{首先…然后…然后…然后…最后…}。
\item \textbf{Variante courte à deux étapes} : \zh{先…然后…} (xiān… ránhòu…) signifie « d'abord… ensuite… ». C'est la forme la plus fréquente à l'oral pour deux actions seulement :
\zh{我们先吃饭，然后去看电影。} \emph{Wǒmen xiān chīfàn, ránhòu qù kàn diànyǐng.} « On mange d'abord, ensuite on va au cinéma. »
\item Attention : \zh{先} se place \textbf{après le sujet, devant le verbe} (c'est un adverbe), tandis que \zh{首先} peut ouvrir la phrase entière.
\end{enumerate}
\end{propriete}

\courssub{Tableau récapitulatif et comparaison avec le français}

\begin{center}
\begin{tabularx}{\linewidth}{|X|X|X|X|}
\hline
\rowcolor{popblueL} Caractères & Pinyin & Fonction & Équivalent français \\
\hline
首先 & shǒuxiān & ouvre la série (étape 1) & d'abord, en premier lieu \\
\hline
然后 & ránhòu & étape suivante (neutre, oral) & ensuite, puis \\
\hline
其次 & qícì & argument suivant (soutenu) & en second lieu, puis \\
\hline
最后 & zuìhòu & ferme la série & enfin, finalement \\
\hline
先…然后… & xiān… ránhòu… & deux étapes seulement (oral) & d'abord… ensuite… \\
\hline
\end{tabularx}
\end{center}

\begin{center}
\begin{tabularx}{\linewidth}{|X|X|X|}
\hline
\rowcolor{popgreenL} Structure & Exemple en caractères + pinyin & Traduction \\
\hline
首先…然后…其次…最后… & \zh{首先我要感谢老师，其次要感谢父母，最后祝大家成功。} \emph{Shǒuxiān wǒ yào gǎnxiè lǎoshī, qícì yào gǎnxiè fùmǔ, zuìhòu zhù dàjiā chénggōng.} & D'abord je remercie les professeurs, ensuite mes parents, enfin je souhaite à tous la réussite. \\
\hline
首先…然后…最后… & \zh{首先要准备好简历，然后参加面试，最后等公司的通知。} \emph{Shǒuxiān yào zhǔnbèi hǎo jiǎnlì, ránhòu cānjiā miànshì, zuìhòu děng gōngsī de tōngzhī.} & Il faut d'abord préparer son CV, ensuite passer l'entretien, enfin attendre la réponse de l'entreprise. \\
\hline
先…然后… & \zh{我先做作业，然后踢足球。} \emph{Wǒ xiān zuò zuòyè, ránhòu tī zúqiú.} & Je fais d'abord mes devoirs, ensuite je joue au football. \\
\hline
\end{tabularx}
\end{center}

\textbf{Comparaison avec le français.} Le français varie librement ses connecteurs (« d'abord, puis, ensuite, après, enfin ») sans ordre obligatoire. Le chinois, lui, préfère une \textbf{chaîne conventionnelle et ordonnée} : \zh{首先} ne peut pas venir après \zh{然后}, et \zh{最后} marque toujours la clôture. Autre différence : le français dit « premièrement, deuxièmement » avec des adverbes numérotés ; le chinois réserve \zh{第一}、\zh{第二} à la numérotation pure (voir le piège ci-dessous).

\courssub{Exemples variés en contexte camerounais}

\begin{exemplebox}[Quatre situations de la vie quotidienne]
\begin{enumerate}
\item \zh{做恩多莱的时候，首先洗好叶子，然后煮花生酱，其次加鱼，最后尝一尝味道。}\\
\emph{Zuò ēnduōlái de shíhou, shǒuxiān xǐ hǎo yèzi, ránhòu zhǔ huāshēngjiàng, qícì jiā yú, zuìhòu cháng yi cháng wèidao.}\\
« Pour préparer le ndolé, on lave d'abord les feuilles, ensuite on cuit la pâte d'arachide, puis on ajoute le poisson, enfin on goûte. »
\item \zh{开学的时候，首先交学费，然后领课本，最后认识新同学。}\\
\emph{Kāixué de shíhou, shǒuxiān jiāo xuéfèi, ránhòu lǐng kèběn, zuìhòu rènshi xīn tóngxué.}\\
« À la rentrée, on paie d'abord les frais de scolarité, ensuite on retire les manuels, enfin on fait connaissance avec les nouveaux camarades. »
\item \zh{我为什么喜欢杜阿拉？首先工作机会多，其次交通方便，最后靠海，风景很美。}\\
\emph{Wǒ wèishénme xǐhuan Dù'ālā? Shǒuxiān gōngzuò jīhuì duō, qícì jiāotōng fāngbiàn, zuìhòu kào hǎi, fēngjǐng hěn měi.}\\
« Pourquoi j'aime Douala ? D'abord il y a beaucoup d'emplois, ensuite les transports sont pratiques, enfin c'est au bord de la mer et le paysage est magnifique. » (Ici, \zh{其次} ordonne des \textbf{arguments}, pas des actions.)
\item \zh{你先告诉我你的名字，然后填这张表。}\\
\emph{Nǐ xiān gàosu wǒ nǐ de míngzi, ránhòu tián zhè zhāng biǎo.}\\
« Donne-moi d'abord ton nom, ensuite remplis ce formulaire. » (Variante courte à deux étapes ; notez le classificateur \zh{张} devant \zh{表}.)
\end{enumerate}
\end{exemplebox}

\courssub{Exercice résolu}

\begin{exoresolu}[Remettre un processus dans l'ordre]
Énoncé : voici les quatre étapes d'une inscription à l'université de Yaoundé, dans le désordre. Replacez chaque connecteur (\zh{首先}, \zh{然后}, \zh{其次}, \zh{最后}) au bon endroit, puis traduisez la phrase obtenue.

\zh{＿＿交注册费，＿＿选择专业，＿＿领学生证，＿＿在网上报名。}
\tcblower
Solution détaillée :
\begin{enumerate}
\item On identifie l'ordre logique des démarches : choisir sa filière → s'inscrire en ligne → payer les frais → retirer sa carte d'étudiant.
\item On place \zh{首先} devant la première étape (\zh{选择专业}), puis \zh{然后} (\zh{在网上报名}), puis \zh{其次} (\zh{交注册费}), enfin \zh{最后} (\zh{领学生证}).
\end{enumerate}
Phrase correcte :

\zh{首先选择专业，然后在网上报名，其次交注册费，最后领学生证。}

\emph{Shǒuxiān xuǎnzé zhuānyè, ránhòu zài wǎngshang bàomíng, qícì jiāo zhùcèfèi, zuìhòu lǐng xuéshēngzhèng.}

« D'abord on choisit sa filière, ensuite on s'inscrit en ligne, puis on paie les frais d'inscription, enfin on retire sa carte d'étudiant. »

Vérification : chaque connecteur introduit bien une proposition verbale complète, et l'ordre \zh{首先 → 然后 → 其次 → 最后} est respecté.
\end{exoresolu}

\courssub{Pièges à éviter : erreurs fréquentes des francophones}

\begin{attention}[Ne traduisez pas « d'abord » par 第一 !]
L'erreur la plus fréquente des élèves francophones est de traduire « d'abord / premièrement » par \zh{第一}. Or :
\begin{itemize}
\item \zh{第一} (dì yī) \textbf{numérote} : c'est « premier / numéro un » dans un classement ou une liste figée : \zh{他是第一名。} \emph{Tā shì dì yī míng.} « Il est premier (de la classe). »
\item \zh{首先} (shǒuxiān) \textbf{introduit une étape} d'un processus ou d'un raisonnement : c'est lui qu'il faut pour « d'abord » dans un récit.
\end{itemize}
Phrase fautive : \zh{第一我要感谢老师…} (à éviter dans un discours d'étapes). Phrase correcte : \zh{首先我要感谢老师…}。\\
Autres pièges :
\begin{itemize}
\item N'ouvrez pas un récit par \zh{然后} : \zh{然后} sert à \textbf{enchaîner} une étape sur une étape déjà mentionnée ; la première étape est introduite par \zh{首先} ou \zh{先}. On dit \zh{我们先吃饭，然后去看电影}, et non \zh{然后我们先吃饭…} en tout début de récit.
\item N'oubliez pas la virgule chinoise \zh{，} entre les propositions : la chaîne de connecteurs relie des phrases, pas des mots.
\item \zh{然后} ne se répète pas à l'infini dans un texte écrit soigné : au-delà de trois étapes, préférez \zh{其次} pour varier.
\end{itemize}
\end{attention}

\courssub{Méthode : rédiger un processus en quatre étapes}

\begin{methode}[Construire un texte avec 首先…然后…其次…最后…]
\begin{enumerate}
\item \textbf{Lister les étapes en français} : écrivez sur le brouillon les 3 ou 4 actions dans l'ordre logique (ex. : préparer le CV → envoyer la candidature → passer l'entretien → attendre la réponse).
\item \textbf{Vérifier le nombre d'étapes} : s'il n'y en a que deux, utilisez la variante courte \zh{先…然后…} ; s'il y en a trois ou quatre, utilisez la chaîne complète.
\item \textbf{Traduire chaque étape} en une proposition chinoise complète (sujet + verbe), avec le vocabulaire connu.
\item \textbf{Greffer les connecteurs} : \zh{首先} devant la première proposition, \zh{然后} devant la deuxième, \zh{其次} devant la troisième (facultatif), \zh{最后} devant la dernière.
\item \textbf{Relire} : vérifier la virgule \zh{，} entre chaque proposition, le point final \zh{。}, et l'ordre immuable des connecteurs.
\end{enumerate}
\end{methode}

\begin{exercice}[À vous de jouer]
Rédigez une phrase de quatre propositions avec la chaîne complète pour décrire l'une des situations suivantes : (a) comment tu prépares l'examen de chinois ; (b) comment on achète des légumes au marché de Mokolo ; (c) comment on organise un match de football au lycée. Le corrigé sera discuté en classe avant de passer à la section suivante sur \zh{既…又…}。
\end{exercice}

\coursec{既…又… : cumuler deux caractéristiques}

\courssub{Transition : de la séquence à la simultanéité}
Dans la section précédente, nous avons vu comment \textbf{首先…然后…其次…最后…} permet d'ordonner des actions dans le temps. Nous passons maintenant à une tout autre logique : non plus la succession, mais la \textbf{coexistence}. Le connecteur \textbf{既…又…} sert à affirmer qu'un même sujet possède \emph{deux} qualités ou accomplit \emph{deux} actions \emph{en même temps} — « à la fois… et… ». C'est un outil indispensable pour décrire, évaluer, argumenter (ex. : « Ce téléphone est à la fois léger \textbf{et} performant »).

\courssub{Observation guidée}
Lisez les deux phrases suivantes, issues d'une fiche technique pour un téléphone portable vendu à Yaoundé. Repérez la structure qui relie les deux adjectifs, puis les deux verbes.

\begin{textebox}[Fiche produit — Extrait]
这款手机\cle{既}轻\cle{又}薄，\cle{既}支持5G\cle{又}支持双卡双待。
\end{textebox}
\begin{center}
\begin{tabularx}{\linewidth}{|X|X|X|X|}
\hline \rowcolor{popblueL} Caractères & Pinyin & Nature & Traduction \\ \hline
这款 & zhè kuǎn & dém. + classif. & ce modèle (de téléphone) \\
手机 & shǒujī & n. & téléphone portable \\
既 & jì & adv. & à la fois (premier terme) \\
轻 & qīng & adj. & léger \\
又 & yòu & adv. & et (deuxième terme, « aussi ») \\
薄 & bó & adj. & fin, mince \\
支持 & zhīchí & v. & supporter, être compatible avec \\
5G & wǔ jī & n. & 5G \\
双卡双待 & shuāng kǎ shuāng dài & loc. & double SIM double veille \\ \hline
\end{tabularx}
\end{center}
\textbf{Traduction :} « Ce modèle de téléphone est à la fois léger et fin ; il supporte à la fois la 5G et la double SIM double veille. »

\noindent\textbf{Questions d'observation :}
\begin{enumerate}
\item Quel mot ouvre la première série ? Quel mot ouvre la deuxième ?
\item Les deux éléments reliés par \textbf{既…又…} sont-ils de même nature grammaticale ?
\item Le sujet est-il le même pour les deux qualités / actions ?
\end{enumerate}

\courssub{Règle générale : forme, emploi, restrictions}

\begin{propriete}[Structure 既…又…]
\textbf{Schéma :} \quad Sujet + \textbf{既} + Adjectif / Verbe 1 + \textbf{又} + Adjectif / Verbe 2

\textbf{Emploi :} Exprimer la \textbf{cumulativité} : le sujet possède simultanément deux qualités (deux adjectifs) ou réalise simultanément deux actions / états (deux verbes). Équivalent français : « à la fois… et… », « à la fois… et aussi… ».

\textbf{Conditions de correction :}
\begin{enumerate}
\item \textbf{Même sujet} pour les deux termes (pas de changement de sujet).
\item \textbf{Parallélisme strict} : les deux éléments doivent être de \textbf{même catégorie syntaxique} (deux adjectifs, deux verbes, deux groupes verbaux). On ne mélange pas adjectif et verbe.
\item \textbf{Registre :} \textbf{既…又…} appartient au style \emph{soutenu / écrit} (articles, discours, descriptions techniques, rédaction formelle). À l'oral familier, on préfère souvent la structure \textbf{又…又…} (ex. : 又便宜又好).
\end{enumerate}
\end{propriete}

\begin{methode}[Construire une phrase avec 既…又…]
\begin{enumerate}
\item Identifier le \textbf{sujet unique}.
\item Choisir deux prédicats de \textbf{même nature} (deux adjectifs \textbf{ou} deux verbes / groupes verbaux).
\item Placer \textbf{既} devant le premier prédicat et \textbf{又} devant le second.
\item Vérifier l'absence de \textbf{和} (hé) ou de virgule entre les deux termes (la structure elle-même assure la liaison).
\end{enumerate}
\end{methode}

\begin{exemplebox}[Exemples variés — Adjectifs / Verbes d'état]
\begin{itemize}
\item 杜阿拉\cle{既}是一个港口，\cle{又}是一个工业中心。\\
Dùālā \textbf{jì} shì yí ge gǎngkǒu, \textbf{yòu} shì yí ge gōngyè zhōngxīn.\\
Douala est à la fois un port et un centre industriel. \quad(\textit{deux groupes « 是 + GN » parallèles})
\item 这道菜\cle{既}鲜嫩\cle{又}入味。\\
Zhè dào cài \textbf{jì} xiānnèn \textbf{yòu} rùwèi.\\
Ce plat est à la fois tendre et savoureux. \quad(\textit{deux adjectifs})
\item 学中文\cle{既}有用\cle{又}有趣。\\
Xué Zhōngwén \textbf{jì} yǒuyòng \textbf{yòu} yǒuqù.\\
Apprendre le chinois est à la fois utile et intéressant. \quad(\textit{deux adjectifs à valeur verbale})
\end{itemize}
\end{exemplebox}

\begin{exemplebox}[Exemples variés — Verbes d'action / Groupes verbaux modaux]
\begin{itemize}
\item 王老师\cle{既}教语法\cle{又}管实习。\\
Wáng lǎoshī \textbf{jì} jiào yǔfǎ \textbf{yòu} guǎn shíxí.\\
M. Wang enseigne à la fois la grammaire et supervise les stages. \quad(\textit{deux VO})
\item 我们\cle{既}要保护环境\cle{又}要发展经济。\\
Wǒmen \textbf{jì} yào bǎohù huánjìng \textbf{yòu} yào fāzhǎn jīngjì.\\
Nous devons à la fois protéger l'environnement et développer l'économie. \quad(\textit{deux groupes verbaux modaux « 要 + VO »})
\item 这款软件\cle{既}能翻译\cle{又}能朗读。\\
Zhè kuǎn ruǎnjiàn \textbf{jì} néng fānyì \textbf{yòu} néng lǎngdú.\\
Ce logiciel peut à la fois traduire et lire à haute voix. \quad(\textit{deux capacités « 能 + V »})
\end{itemize}
\end{exemplebox}

\courssub{Visualisation de la structure}

\begin{popfigure}
\begin{tikzpicture}[
  node distance=1.2cm and 0.8cm,
  main/.style={draw=popblue, thick, rounded corners, fill=popblueL, align=center, minimum height=1.1cm},
  brace/.style={decorate, decoration={brace, amplitude=6pt}, thick, popdark},
  arrow/.style={-Stealth, thick, poporange},
  labelbox/.style={font=\small\bfseries, text=popdark}
]
\node[main] (sujet) {SUJET};
\node[main, right=of sujet] (ji) {既 + Élément 1};
\node[main, below=of ji] (you) {又 + Élément 2};
\draw[brace] (ji.north west) -- (you.south west) node[midway, left=4pt, labelbox, align=center]{même nature\\grammaticale};
\draw[arrow] (sujet) -- (ji);
\draw[arrow] (sujet) -- (you);
\node[labelbox, above=0.2cm of sujet] {Ex. : 这款手机};
\node[labelbox, right=0.2cm of ji] {轻 (adj.)};
\node[labelbox, right=0.2cm of you] {薄 (adj.)};
\end{tikzpicture}
\legende{Schéma de la structure 既…又… : un seul sujet gouverne deux éléments parallèles reliés par 既 et 又.}
\end{popfigure}

\courssub{Comparaison avec le français et pièges fréquents}

\begin{attention}[Erreurs typiques des francophones]
\begin{enumerate}
\item \textbf{Utiliser \emph{和} (hé) entre deux adjectifs :} \zh{*这个产品便宜和实用} ✗. En chinois, \textbf{和} relie des \emph{noms}, pas des adjectifs. \textbf{Correction :} 这个产品\cle{既}便宜\cle{又}实用。 ou 这个产品\cle{又}便宜\cle{又}实用。 (oral) ou 这个产品\cle{很}便宜，\cle{也}很实用。
\item \textbf{Mélanger les catégories syntaxiques :} \zh{*他既高大又跑步} ✗ (adjectif + verbe). \textbf{Correction :} 他既高大\cle{又}强壮 (deux adj.) \textbf{ou} 他既\cle{会}打篮球\cle{又}会跑步 (deux verbes).
\item \textbf{Changer de sujet :} \zh{*既我喜欢，又他不喜欢} ✗. \textbf{既…又…} ne relie qu'un seul sujet. Pour deux sujets, utilisez \textbf{也} : 我喜欢，他\cle{也}喜欢。
\item \textbf{Confondre 既 (jì) et 即 (jí) :} \textbf{既} (déjà, à la fois) vs \textbf{即} (immédiatement, c'est-à-dire). Ex. : \zh{既然} (jìrán, « puisque ») contient 既 ; \zh{即刻} (jíkè, « immédiatement ») contient 即.
\end{enumerate}
\end{attention}

\begin{savaistu}[既 seul dans des expressions figées]
Le caractère \textbf{既} (jì) apparaît seul dans plusieurs mots courants où il garde le sens de « déjà / depuis que » :
\begin{itemize}
\item \zh{既然} (jìrán) — « puisque, dès lors que » : \zh{既然来了，就安心学习。} (Puisque tu es venu, étudie tranquillement.)
\item \zh{既定} (jìdìng) — « préétabli, fixé » : \zh{既定方针} (ligne directrice fixée).
\item \zh{既得利益} (jìdé lìyì) — « intérêts acquis / établis ».
\end{itemize}
Ne les confondez pas avec le connecteur de cumul \textbf{既…又…} étudié ici.
\end{savaistu}

\courssub{Tableau récapitulatif : 既…又… vs 又…又… vs 很…也很… vs 不但…而且…}

\begin{center}
\begin{tabularx}{\linewidth}{|X|X|X|X|}
\hline \rowcolor{popblueL} Structure & Registre & Nature des éléments & Exemple (car. + pinyin + trad.) \\ \hline
\textbf{既…又…} & Soutenu, écrit & Adj. / V. parallèles & 既便宜又实用 (à la fois bon marché et pratique) \\ \hline
\textbf{又…又…} & Oral, familier & Adj. (souvent monosyll.) & 又快又好 (à la fois rapide et bien fait) \\ \hline
\textbf{很…也很…} & Neutre, universel & Adj. / SV & 很便宜，也很实用 (très bon marché et aussi très pratique) \\ \hline
\textbf{不但…而且…} & Neutre, argumentatif & V. / phrases complètes & 不但便宜，而且耐用 (non seulement bon marché, mais encore durable) \\ \hline
\end{tabularx}
\end{center}

\begin{aretenir}[Points clés à retenir]
\begin{itemize}
\item \textbf{既…又…} = cumul de deux propriétés/actions sur \textbf{un même sujet}.
\item \textbf{Parallélisme obligatoire} : Adj + Adj \textbf{ou} V + V (même structure interne).
\item \textbf{Registre soutenu} : privilégié à l'écrit, aux exposés, aux descriptions formelles.
\item \textbf{Pas de « 和 »} entre les deux termes.
\item Distinction graphique : \textbf{既} (jì, cumul) $\neq$ \textbf{即} (jí, immédiat / c'est-à-dire).
\end{itemize}
\end{aretenir}

\courssub{Exercice résolu : transformation et correction}

\begin{exoresolu}[Transformer en 既…又… (style soutenu) / Corriger les phrases]
\textbf{Énoncé :}
\begin{enumerate}
\item Réécrivez chaque phrase en utilisant \textbf{既…又…}.
\item Corrigez les phrases fautives en justifiant.
\end{enumerate}
\textbf{Phrases :}
\begin{enumerate}
\item 这个城市很现代，也很古老。 (adj. + adj.)
\item 她会弹钢琴，也会拉小提琴。 (V. + V.)
\item *这本书既有趣和有用。 (erreur : 和)
\item *既我去，又他去。 (erreur : deux sujets)
\item 这道菜又辣又烫。 (passer au style soutenu)
\end{enumerate}

\tcblower
\textbf{Solution détaillée :}
\begin{enumerate}
\item 这个城市\cle{既}现代\cle{又}古老。 (Zhège chéngshì jì xiàndài yòu gǔlǎo.) — Cette ville est à la fois moderne et ancienne. \quad(\textit{Parallélisme adjectif + adjectif respecté})
\item 她\cle{既}会弹钢琴\cle{又}会拉小提琴。 (Tā jì huì tán gāngqín yòu huì lā xiǎotíqín.) — Elle sait à la fois jouer du piano et du violon. \quad(\textit{Parallélisme « 会 + VO » respecté})
\item \textbf{Correction :} 这本书\cle{既}有趣\cle{又}有用。 (\textbf{既…又…} ne s'utilise jamais avec \textbf{和}).
\item \textbf{Correction :} 我\cle{既}去\cle{又}去 (un seul sujet, action répétée) \textbf{ou mieux} : 我去，他\cle{也}去 (deux sujets différents $\rightarrow$ \textbf{也}).
\item 这道菜\cle{既}辣\cle{又}烫。 (Zhè dào cài jì là yòu tàng.) — Ce plat est à la fois épicé et brûlant. \quad(\textit{Passage oral 又…又… $\rightarrow$ écrit 既…又…})
\end{enumerate}
\end{exoresolu}

\courssub{Entraînement guidé : production écrite}

\begin{exercice}[À vous de jouer — Description d'un lieu camerounais]
Rédigez un court paragraphe (4–5 phrases) pour présenter \textbf{le marché Mokolo à Yaoundé} ou \textbf{le mont Cameroun} en utilisant \textbf{au moins trois fois} la structure \textbf{既…又…} (adjectifs et/ou verbes). Respectez le parallélisme et le registre soutenu.

\textbf{Mots utiles :} 繁华 (fánhuá, animé/prospère), 历史悠久 (lìshǐ yōujiǔ, longue histoire), 热闹 (rènào, animé), 壮观 (zhuàngguān, spectaculaire), 吸引 (xīyǐn, attirer), 游客 (yóukè, touristes), 商品 (shāngpǐn, marchandises), 价格实惠 (jiàgé shíhuì, prix abordables), 生物多样性 (shēngwù duōyàngxìng, biodiversité), 登山者 (dēngshānzhě, alpinistes).
\end{exercice}

\begin{corrige}[Exercice — Proposition de correction]
\textbf{Le marché Mokolo :}
莫科洛市场\cle{既}繁华\cle{又}历史悠久，\cle{既}商品齐全\cle{又}价格实惠，\cle{既}吸引本地居民\cle{又}吸引外地游客。
(Mòkēluò shìchǎng jì fánhuá yòu lìshǐ yōujiǔ, jì shāngpǐn qíquán yòu jiàgé shíhuì, jì xīyǐn běndì jūmín yòu xīyǐn wàidì yóukè.)
Le marché Mokolo est à la fois prospère et chargé d'histoire ; il offre à la fois une gamme complète de marchandises et des prix abordables ; il attire à la fois les habitants locaux et les touristes d'ailleurs.

\textbf{Le mont Cameroun (variante) :}
喀麦隆山\cle{既}壮观\cle{又}生物多样性丰富，\cle{既}吸引科学家\cle{又}吸引登山者，\cle{既}是国家象征\cle{又}是自然遗产。
(Kāmàilóng shān jì zhuàngguān yòu shēngwù duōyàngxìng fēngfù, jì xīyǐn kēxuéjiā yòu xīyǐn dēngshānzhě, jì shì guójiā xiàngzhēng yòu shì zìrán yíchǎn.)
Le mont Cameroun est à la fois spectaculaire et riche en biodiversité ; il attire à la fois les scientifiques et les alpinistes ; il est à la fois un symbole national et un patrimoine naturel.
\end{corrige}

\begin{experience}[Activité orale : « Le produit idéal »]
En binôme, imaginez un objet innovant (ex. : un sac à dos solaire, une application pour les taxis moto « bend skin »). Présentez-le à la classe en utilisant \textbf{au moins deux fois} \textbf{既…又…} pour vanter ses mérites (qualités + fonctions). Les autres élèves notent les structures entendues et votent pour le meilleur « pitch ».
\end{experience}

\coursec{只有…才… : la condition nécessaire}

Dans la section précédente, nous avons appris à \cle{cumuler} deux caractéristiques avec \zh{既…又…} (jì…yòu…) : une personne peut être à la fois sérieuse \emph{et} gentille, un produit à la fois bon \emph{et} bon marché. Nous passions ainsi d'une idée à une autre idée \emph{additionnée}. Nous changeons maintenant de logique : avec \zh{只有…才…}, il ne s'agit plus d'additionner, mais de poser une \cle{condition indispensable}. Entre les deux propositions, il existe un lien de nécessité : \emph{sans la première, la seconde est impossible}. C'est exactement ce qu'exprime le français « ce n'est que… que… » ou « seulement si… ».

\courssub{Observation : découvrir la structure dans un texte}

Lisons d'abord ce court texte sur la recherche d'emploi, thème central de notre module « Activités économiques et monde du travail ».

\begin{textebox}[À l'agence pour l'emploi de Douala]
在杜阿拉，很多年轻人想找好工作。\\
\emph{Zài Dù'ālā, hěn duō niánqīngrén xiǎng zhǎo hǎo gōngzuò.}\\
« À Douala, beaucoup de jeunes veulent trouver un bon emploi. »\\[4pt]
只有努力学习，才能考上好大学。\\
\emph{Zhǐyǒu nǔlì xuéxí, cái néng kǎoshang hǎo dàxué.}\\
« Ce n'est qu'en travaillant dur qu'on peut entrer dans une bonne université. »\\[4pt]
只有会说汉语，才能在中国公司工作。\\
\emph{Zhǐyǒu huì shuō Hànyǔ, cái néng zài Zhōngguó gōngsī gōngzuò.}\\
« Ce n'est qu'en sachant parler chinois qu'on peut travailler dans une entreprise chinoise. »\\[4pt]
只有你去，他才来。\\
\emph{Zhǐyǒu nǐ qù, tā cái lái.}\\
« Il ne viendra que si tu y vas. »
\end{textebox}

\textbf{Questions d'observation}
\begin{enumerate}
\item Soulignez dans chaque phrase les deux mots qui se répondent : \zh{只有} et \zh{才}. Où se placent-ils ?
\item Dans la phrase \zh{只有你去，他才来}, supprimez mentalement \zh{只有} : le sens change-t-il ?
\item Traduisez chaque phrase en français. Quelle tournure française revient à chaque fois ?
\end{enumerate}

On observe que : \zh{只有} ouvre la première proposition (la \cle{condition}), et \zh{才} se place juste \emph{avant le verbe} de la seconde proposition (le \cle{résultat}). Le français rend cette relation par « ce n'est que… que… », « seulement si… », « il faut… pour… ».

\courssub{La règle : schéma, emploi, place des mots}

\begin{propriete}[只有…才… : la condition nécessaire]
\zh{只有…才…} exprime une \cle{condition nécessaire} : le résultat de la seconde proposition n'est possible \textbf{que si} la condition de la première proposition est remplie. Sans la condition, le résultat est impossible.\par
\textbf{Schéma :} 只有 + condition ，（sujet）+ 才 + verbe / résultat\par
\begin{itemize}
\item \zh{只有} (zhǐyǒu, « seulement ») introduit la condition, souvent suivie d'un verbe ou d'un groupe verbal : \zh{只有努力}, \zh{只有会说汉语}, \zh{只有你去}.
\item \zh{才} (cái) se place \textbf{avant le verbe} de la seconde proposition, après le sujet s'il y en a un : \zh{才能工作}, \zh{他才来}.
\item Les deux éléments sont \textbf{inséparables} : \zh{只有} appelle obligatoirement \zh{才}.
\end{itemize}
\end{propriete}

\begin{exemplebox}[Exemples fondateurs]
\zh{只有会说汉语，才能在中国公司工作。}\\
\emph{Zhǐyǒu huì shuō Hànyǔ, cái néng zài Zhōngguó gōngsī gōngzuò.}\\
« Ce n'est qu'en sachant parler chinois qu'on peut travailler dans une entreprise chinoise. »\\[4pt]
\zh{只有你去，他才来。}\\
\emph{Zhǐyǒu nǐ qù, tā cái lái.}\\
« Il ne viendra que si tu y vas. » (deux sujets différents : \zh{你} pour la condition, \zh{他} pour le résultat)\\[4pt]
\zh{只有每天练习，才能写好汉字。}\\
\emph{Zhǐyǒu měitiān liànxí, cái néng xiě hǎo Hànzì.}\\
« Ce n'est qu'en s'entraînant chaque jour qu'on peut bien écrire les caractères. »
\end{exemplebox}

\begin{popfigure}
\begin{tikzpicture}[
  box/.style={draw=popblue, fill=popblueL, rounded corners=4pt, align=center, minimum height=1.2cm, text width=4.2cm, font=\small},
  rbox/.style={draw=popgreen, fill=popgreenL, rounded corners=4pt, align=center, minimum height=1.2cm, text width=4.2cm, font=\small},
  arr/.style={-{Stealth[length=3mm]}, very thick, poporange}
]
\node[box, align=center] (cond) at (0,0){\textbf{Condition}\\ 只有 + verbe\\ \zh{只有努力学习}};
\node[rbox, align=center] (res) at (7.2,0){\textbf{Résultat}\\ 才 + verbe\\ \zh{才能成功}};
\draw[arr] (cond) -- node[above, font=\small\itshape, text=popdark]{indispensable} (res);
\node[align=center, font=\small, text=popmuted] at (3.6,-1.4) {Sans la condition $\rightarrow$ résultat impossible};
\end{tikzpicture}
\legende{Organigramme de la condition nécessaire : la flèche n'existe que si la condition est remplie.}
\end{popfigure}

\courssub{Comparaison avec le français}

Le français dispose de plusieurs tournures pour la condition nécessaire ; toutes se traduisent par \zh{只有…才…} :

\begin{center}
\begin{tabularx}{\linewidth}{|X|X|X|}
\hline
\rowcolor{popblueL} \textbf{Français} & \textbf{Chinois (caractères + pinyin)} & \textbf{Remarque} \\
\hline
« Ce n'est qu'en travaillant qu'on réussit. » & \zh{只有努力，才能成功。} \emph{Zhǐyǒu nǔlì, cái néng chénggōng.} & tournure la plus proche \\
\hline
« Seulement si tu viens, je reste. » & \zh{只有你来，我才留下。} \emph{Zhǐyǒu nǐ lái, wǒ cái liúxià.} & deux sujets explicites \\
\hline
« Il faut avoir un visa pour entrer en Chine. » & \zh{只有有签证，才能去中国。} \emph{Zhǐyǒu yǒu qiānzhèng, cái néng qù Zhōngguó.} & « il faut… pour… » \\
\hline
« On ne peut entrer qu'avec un billet. » & \zh{只有有票，才能进去。} \emph{Zhǐyǒu yǒu piào, cái néng jìnqù.} & restriction \\
\hline
\end{tabularx}
\end{center}

Remarquez que le français varie les tournures (« ce n'est que », « seulement si », « il faut », « ne… que »), alors que le chinois utilise \textbf{toujours la même paire} \zh{只有…才…}. C'est une bonne nouvelle : une seule structure à mémoriser pour quatre tournures françaises !

\courssub{La nuance de 才 : la restriction}

Dans cette structure, \zh{才} n'est pas un simple connecteur : il marque la \cle{restriction}. Il signifie littéralement « seulement alors », « dans ce cas seulement ». Le résultat est \emph{limité} à cette seule condition : aucune autre voie ne mène au résultat.

\begin{exemplebox}[Sentir la restriction]
\zh{只有下雨，比赛才取消。}\\
\emph{Zhǐyǒu xià yǔ, bǐsài cái qǔxiāo.}\\
« Le match ne sera annulé que s'il pleut. » (s'il ne pleut pas, le match aura lieu : la pluie est la \emph{seule} cause possible d'annulation)\\[4pt]
\zh{只有老师同意，我们才能去参观工厂。}\\
\emph{Zhǐyǒu lǎoshī tóngyì, wǒmen cái néng qù cānguān gōngchǎng.}\\
« Nous ne pourrons visiter l'usine que si le professeur accepte. »
\end{exemplebox}

\begin{attention}[Ne pas confondre les emplois de 才]
\zh{才} (cái) possède un autre emploi fréquent : il marque la \textbf{tardiveté}, « ne… que » au sens temporel, et il est alors \emph{seul}, sans \zh{只有} :
\zh{他八点才来。} \emph{Tā bā diǎn cái lái.} « Il n'est arrivé qu'à huit heures. » (il est arrivé tard)\par
Dans \zh{只有…才…}, \zh{才} marque la \textbf{condition} ; seul, après un complément de temps, il marque le \textbf{retard}. Même caractère, deux emplois : regardez toujours si \zh{只有} est présent dans la phrase.
\end{attention}

\courssub{Erreurs fréquentes des francophones}

\begin{attention}[Le partenaire obligatoire de 只有 est 才]
Erreur n°1 : oublier \zh{才}. *\zh{只有努力，能成功。} est \textbf{fautif} : \zh{只有} ne peut pas rester sans réponse.\par
Erreur n°2 : remplacer \zh{才} par \zh{就}. *\zh{只有努力，就能成功。} est \textbf{fautif}. \zh{就} est le partenaire de \zh{只要} (condition suffisante, section suivante), jamais de \zh{只有}.\par
Retenez le couple comme une expression figée : \textbf{只有…才…}, et non \zh{只有…就…}.\par
Erreur n°3 : mal placer \zh{才}. Il se place \emph{avant le verbe}, après le sujet : \zh{他才来} et non *\zh{才他来}.
\end{attention}

\begin{methode}[Comment choisir 只有…才… ?]
\begin{enumerate}
\item Traduisez mentalement la phrase française : peut-on la formuler avec « \textbf{ce n'est que si…} » ou « \textbf{seulement si…} » ? Si oui, c'est \zh{只有…才…}.
\item Identifiez la condition (ce qui est exigé) : placez-la après \zh{只有}.
\item Identifiez le résultat : placez \zh{才} juste avant son verbe.
\item Vérifiez le couple : \zh{只有} au début, \zh{才} avant le verbe du résultat — jamais \zh{就}.
\end{enumerate}
\end{methode}

\begin{exoresolu}[Traduction guidée]
Énoncé : traduisez en chinois « Ce n'est qu'en parlant chinois tous les jours qu'on progresse vite. »
\tcblower
\textbf{Étape 1} : la phrase contient « ce n'est qu'en… que… » $\Rightarrow$ structure \zh{只有…才…}.\par
\textbf{Étape 2} : condition = « parler chinois tous les jours » $\Rightarrow$ \zh{每天说汉语} (měitiān shuō Hànyǔ). On la place après \zh{只有} : \zh{只有每天说汉语}.\par
\textbf{Étape 3} : résultat = « progresser vite » $\Rightarrow$ \zh{进步很快} (jìnbù hěn kuài). On place \zh{才} avant : \zh{才进步很快}.\par
\textbf{Solution} : \zh{只有每天说汉语，才进步很快。} \emph{Zhǐyǒu měitiān shuō Hànyǔ, cái jìnbù hěn kuài.}
\end{exoresolu}

\begin{exercice}[À vous de jouer]
Traduisez en chinois avec \zh{只有…才…} :
\begin{enumerate}
\item Ce n'est qu'en ayant un passeport qu'on peut voyager à l'étranger.
\item Il ne viendra au marché que si tu l'appelles.
\item Ce n'est qu'en écoutant le professeur qu'on comprend la leçon.
\end{enumerate}
\end{exercice}

\begin{corrige}[Exercice « À vous de jouer »]
\begin{enumerate}
\item \zh{只有有护照，才能去国外旅行。} \emph{Zhǐyǒu yǒu hùzhào, cái néng qù guówài lǚxíng.}
\item \zh{只有你给他打电话，他才来市场。} \emph{Zhǐyǒu nǐ gěi tā dǎ diànhuà, tā cái lái shìchǎng.}
\item \zh{只有听老师的话，才能听懂课文。} \emph{Zhǐyǒu tīng lǎoshī de huà, cái néng tīng dǒng kèwén.}
\end{enumerate}
\end{corrige}

\courssub{Production : à l'oral et à l'écrit}

\begin{experience}[Conditions de réussite : débat guidé]
\textbf{Consigne :} En binôme, préparez un court exposé (1 minute) sur le thème : « Réussir ses études / Trouver un bon emploi au Cameroun / Travailler dans une entreprise chinoise ». Utilisez \textbf{au moins trois phrases} avec \zh{只有…才…}.\par
\textbf{Exemple de structure :}\\
\zh{只有认真复习，才能考上好大学。}\\
\zh{只有会说英语和汉语，才能在大公司工作。}\\
\zh{只有积累经验，才能成为专家。}\\
\emph{Zhǐyǒu rènzhēn fùxí, cái néng kǎoshang hǎo dàxué.}\\
\emph{Zhǐyǒu huì shuō Yīngyǔ hé Hànyǔ, cái néng zài dà gōngsī gōngzuò.}\\
\emph{Zhǐyǒu jīlěi jīngyàn, cái néng chéngwéi zhuānjiā.}\\
\textbf{Déroulement :} 5 min de préparation \textrightarrow{} 2 min de présentation par binôme \textrightarrow{} vote de la classe pour l'exposé le plus convaincant.
\end{experience}

\begin{aretenir}[L'essentiel]
\zh{只有 + condition，才 + résultat} exprime la \cle{condition nécessaire} : « ce n'est que… que… », « seulement si… ». \zh{才} se place avant le verbe du résultat et marque la restriction. Le partenaire de \zh{只有} est \textbf{toujours} \zh{才}, jamais \zh{就}. Attention à ne pas confondre ce \zh{才} avec le \zh{才} de tardiveté (\zh{他八点才来}).
\end{aretenir}

\coursec{只要…就… : la condition suffisante et comparaison des deux conditions}

\courssub{Transition et observation guidée}

Dans la section précédente, nous avons vu que \textbf{只有…才…} exprime une \emph{condition nécessaire} : sans l'effort, pas de réussite. Nous découvrons maintenant \textbf{只要…就…}, qui exprime une \emph{condition suffisante} : la réalisation de la condition \emph{garantit} le résultat. C'est toute la différence entre « il faut au moins cela » et « cela suffit amplement ».

Observons une situation professionnelle camerounaise :

\begin{textebox}[Dialogue : Conditions de paiement et de livraison]
\textbf{Contexte :} Un entrepreneur de Douala (M. Ngono) négocie avec une fournisseuse de Yaoundé (Mme Fouda).

M. Ngono：\zh{只要你们按时发货，我们就马上付款。} \\
(Zhǐyào nǐmen ànshí fāhuò, wǒmen jiù mǎshàng fùkuǎn.) \\
— « Tant que vous expédiez à temps, nous payons immédiatement. »

Mme Fouda：\zh{只有收到定金，我才能安排发货。} \\
(Zhǐyǒu shōudào dìngjīn, wǒ cái néng ānpái fāhuò.) \\
— « Ce n'est qu'après avoir reçu l'acompte que je peux organiser l'expédition. »

M. Ngono：\zh{只要合同签好，定金今天就转账。} \\
(Zhǐyào hétong qiān hǎo, dìngjīn jīntiān jiù zhuǎnzhàng.) \\
— « Du moment que le contrat est signé, l'acompte est viré aujourd'hui. »

Mme Fouda：\zh{好，只要钱到账，货明天就发。} \\
(Hǎo, zhǐyào qián dào zhàng, huò míngtiān jiù fā.) \\
— « D'accord, dès que l'argent arrive sur le compte, la marchandise part demain. »
\end{textebox}

\begin{aretenir}[Analyse rapide]
\begin{itemize}
    \item \zh{只要你们按时发货，我们就马上付款} : la condition (expédier à temps) \emph{suffit} à déclencher le paiement. C'est \textbf{只要…就…}.
    \item \zh{只有收到定金，我才能安排发货} : l'acompte est \emph{indispensable} (condition nécessaire). C'est \textbf{只有…才…}.
    \item Remarquez la position de \zh{就} (juste avant le verbe : \zh{付款}, \zh{转账}, \zh{发}) et de \zh{才} (juste avant \zh{能安排}).
\end{itemize}
\end{aretenir}

\courssub{Structure et emploi de 只要…就…}

\begin{propriete}[Règle : 只要…就… = condition suffisante]
\textbf{Schéma :} \zh{只要} + \textbf{Condition} ，(+ \textbf{Sujet}) + \zh{就} + \textbf{Résultat}

\textbf{Emploi :}
\begin{itemize}
    \item Indique que \emph{dès que} la condition est remplie, le résultat se produit inévitablement.
    \item Équivalents français : « tant que… », « du moment que… », « pourvu que… », « dès que… ».
    \item La condition est \textbf{suffisante} : elle seule assure la réalisation du résultat (aucune autre condition n'est exigée).
\end{itemize}

\textbf{Position de \zh{就} :}
\begin{itemize}
    \item \zh{就} se place \textbf{avant le verbe (ou le verbe modal)} de la proposition principale.
    \item Si le sujet est répété dans la seconde proposition, \zh{就} suit ce sujet : \zh{只要 + 条件，主语 + 就 + 动词}.
    \item Si le sujet est le même dans les deux propositions, on l'omet souvent dans la seconde et \zh{就} précède directement le verbe : \zh{只要 + 条件，就 + 动词}.
\end{itemize}

\textbf{Négation :} on nie en général le résultat dans la proposition principale (\zh{就不} / \zh{就不会}). Pour nier la condition, on préfère la tournure \zh{只要不…就…} (« tant que… ne… pas »).
\end{propriete}

\begin{exemplebox}[Exemples variés avec 只要…就…]
\begin{itemize}
    \item \zh{只要你努力学习，就一定能通过HSK4。} \\
    (Zhǐyào nǐ nǔlì xuéxí, jiù yídìng néng tōngguò HSK4.) \\
    — « Du moment que tu étudies sérieusement, tu réussiras certainement le HSK4. »

    \item \zh{只要天不下雨，我们就去克里比海边野餐。} \\
    (Zhǐyào tiān bù xiàyǔ, wǒmen jiù qù Kèlǐbǐ hǎibiān yěcān.) \\
    — « Tant qu'il ne pleut pas, nous allons pique-niquer au bord de la mer à Kribi. »

    \item \zh{这个软件只要安装好，就能自动翻译中文。} \\
    (Zhège ruǎnjiàn zhǐyào ānzhuāng hǎo, jiù néng zìdòng fānyì Zhōngwén.) \\
    — « Ce logiciel, dès qu'il est bien installé, peut traduire le chinois automatiquement. »

    \item \zh{只要大家配合，工作就能按时完成。} \\
    (Zhǐyào dàjiā pèihé, gōngzuò jiù néng ànshí wánchéng.) \\
    — « Pourvu que tout le monde coopère, le travail pourra être terminé à temps. »
\end{itemize}
\end{exemplebox}

\courssub{Comparaison contrastive : 只有…才… contre 只要…就…}

C'est le point crucial pour le baccalauréat et pour une communication précise. Ne les confondez plus jamais.

\begin{center}
\begin{tabularx}{\linewidth}{|X|X|X|}
\hline
\rowcolor{popblueL}
\textbf{Critère} & \textbf{只有 … 才 …} & \textbf{只要 … 就 …} \\ \hline
\textbf{Type de condition} & Condition \textbf{nécessaire} (indispensable) & Condition \textbf{suffisante} (garantie) \\ \hline
\textbf{Logique} & Sans A $\rightarrow$ pas de B. \newline (A est le \emph{seul} chemin vers B) & Avec A $\rightarrow$ B assuré. \newline (A \emph{suffit} pour obtenir B) \\ \hline
\textbf{Équivalents français} & « ce n'est que si… que… », « seulement si… » & « tant que… », « du moment que… », « pourvu que… » \\ \hline
\textbf{Degré d'exigence} & \textbf{Élevé, restrictif} : A est obligatoire, mais peut ne pas suffire. & \textbf{Ouvert} : A suffit, même si d'autres conditions pourraient aussi convenir. \\ \hline
\textbf{Exemple type} & \zh{只有办好签证，才能去中国。} \newline (Zhǐyǒu bàn hǎo qiānzhèng, cái néng qù Zhōngguó.) \newline « Ce n'est qu'avec un visa qu'on peut aller en Chine. » & \zh{只要有车票，就能上车。} \newline (Zhǐyào yǒu chēpiào, jiù néng shàngchē.) \newline « Du moment qu'on a un billet, on peut monter à bord. » \\ \hline
\textbf{Piège} & Ne jamais écrire *只有…就… & Ne jamais écrire *只要…才… \\ \hline
\end{tabularx}
\end{center}

\begin{popfigure}
\begin{tikzpicture}[
    node distance=1.5cm,
    box/.style={draw=popblue, thick, rounded corners, fill=popblueL, text width=6.8cm, align=center, inner sep=8pt}
]
\node[box, align=center] (nec){
    {\large\bfseries\color{popdark} 只有 … 才 …}\\[2pt]
    {\small (condition nécessaire)}\\[4pt]
    {\small \textbf{Logique :} sans A $\rightarrow$ pas de B}\\[2pt]
    {\small \textbf{Français :} « ce n'est que si A que B »}\\[4pt]
    {\small \zh{只有努力，才会成功。}}\\
    {\small (Zhǐyǒu nǔlì, cái huì chénggōng.)}\\
    {\small « Ce n'est qu'en travaillant dur qu'on réussit. »}
};
\node[box, right=of nec, align=center] (suf){
    {\large\bfseries\color{popdark} 只要 … 就 …}\\[2pt]
    {\small (condition suffisante)}\\[4pt]
    {\small \textbf{Logique :} avec A $\rightarrow$ B assuré}\\[2pt]
    {\small \textbf{Français :} « tant que A, B »}\\[4pt]
    {\small \zh{只要努力，就会成功。}}\\
    {\small (Zhǐyào nǔlì, jiù huì chénggōng.)}\\
    {\small « Du moment qu'on travaille dur, on réussit. »}
};
\draw[<->, very thick, poporange] (nec.east) -- node[above, font=\small\bfseries, text=poporange, align=center] {degré\\d'exigence} (suf.west);
\end{tikzpicture}
\legende{Comparaison visuelle : condition nécessaire (restrictive) contre condition suffisante (garantie).}
\end{popfigure}

\begin{attention}[Erreurs fatales à bannir — pièges du baccalauréat]
\begin{enumerate}
    \item \textbf{Croisement interdit 1 :} *\zh{只有你努力，就能成功。} $\rightarrow$ \textbf{FAUX}. \zh{只有} appelle obligatoirement \zh{才}.
    \item \textbf{Croisement interdit 2 :} *\zh{只要你努力，才能成功。} $\rightarrow$ \textbf{FAUX}. \zh{只要} appelle obligatoirement \zh{就}.
    \item \textbf{Oubli de \zh{就} ou \zh{才} :} *\zh{只要你来，我高兴。} $\rightarrow$ \textbf{INCOMPLET}. Il faut \zh{只要你来，我就高兴。} (Zhǐyào nǐ lái, wǒ jiù gāoxìng.) — « Tant que tu viens, je suis content. »
    \item \textbf{Ordre des mots :} \zh{就} et \zh{才} se placent \textbf{avant} le verbe ou le modal, et \textbf{après} le sujet s'il est exprimé.
    \begin{itemize}
        \item Correct : \zh{他只要来，我们就走。} (Tā zhǐyào lái, wǒmen jiù zǒu.) — « Dès qu'il vient, nous partons. »
        \item Incorrect : *\zh{他只要来，我们就走就。}
    \end{itemize}
\end{enumerate}

\textbf{Moyen mnémotechnique :} \cle{有才，要就} (yǒu cái, yào jiù).
\begin{itemize}
    \item \zh{有} (dans \zh{只有}) va avec \zh{才}.
    \item \zh{要} (dans \zh{只要}) va avec \zh{就}.
\end{itemize}
Récitez-le : « \emph{Seul celui qui a du talent (有才) obtient ce qu'il veut (要就)} ».
\end{attention}

\courssub{Méthode : transformer une phrase pour tester la logique}

\begin{methode}[Passer de 只要…就… à 只有…才… (et inversement)]
\textbf{Objectif :} vérifier comment le sens de la phrase change quand on remplace une structure par l'autre.

\begin{enumerate}
    \item \textbf{Isolez la condition (A) et le résultat (B).} \\
    Exemple : \zh{只要下雨，比赛就取消。} (A = \zh{下雨} « pleuvoir », B = \zh{比赛取消} « le match est annulé »).
    \item \textbf{Appliquez la transformation :} remplacez \zh{只要} par \zh{只有} \emph{et} \zh{就} par \zh{才} (les deux ensemble !). \\
    $\rightarrow$ \zh{只有下雨，比赛才取消。}
    \item \textbf{Interprétez les deux versions :}
    \begin{itemize}
        \item Version 1 (\zh{只要…就…}) : « S'il pleut, le match est annulé. » La pluie \textbf{suffit} à annuler (mais la neige pourrait aussi annuler le match).
        \item Version 2 (\zh{只有…才…}) : « \textbf{Seulement} s'il pleut, le match est annulé. » La pluie est la \textbf{seule} cause d'annulation (s'il neige, le match a lieu !).
    \end{itemize}
    \item \textbf{Choisissez selon la réalité :} dans la vie réelle, un match peut être annulé pour la pluie \textbf{ou} la neige. La version \zh{只要} (la pluie suffit) est donc logique ; la version \zh{只有} (seule la pluie annule) est factuellement fausse.
\end{enumerate}
\end{methode}

\begin{exoresolu}[Transformation et choix du connecteur]
\textbf{Énoncé :} Transformez les phrases suivantes en utilisant l'autre structure (\zh{只有…才…} $\leftrightarrow$ \zh{只要…就…}) et dites si le sens devient plus restrictif ou plus large.

1. \zh{只有通过面试，才能得到这份工作。} (Zhǐyǒu tōngguò miànshì, cái néng dédào zhè fèn gōngzuò.) — « Ce n'est qu'en réussissant l'entretien qu'on peut obtenir ce poste. »

2. \zh{只要你会说中文，就能申请这个岗位。} (Zhǐyào nǐ huì shuō Zhōngwén, jiù néng shēnqǐng zhège gǎngwèi.) — « Du moment que tu sais parler chinois, tu peux postuler à ce poste. »

3. \zh{只有缴纳会费，才能参加协会活动。} (Zhǐyǒu jiǎonà huìfèi, cái néng cānjiā xiéhuì huódòng.) — « Ce n'est qu'en payant la cotisation qu'on peut participer aux activités de l'association. »

\tcblower
\textbf{Solution détaillée :}

\textbf{1.} Transformation : \zh{只要通过面试，就能得到这份工作。} \\
\textbf{Analyse :} l'originale (\zh{只有}) fait de l'entretien une étape \textbf{nécessaire} (mais il peut y avoir aussi un test écrit). La transformée (\zh{只要}) affirme que l'entretien \textbf{suffit} (aucun autre test). \textbf{Sens : plus large.} Dans la réalité, l'entretien est souvent nécessaire mais pas suffisant : la version \zh{只要} est donc souvent trop optimiste.

\textbf{2.} Transformation : \zh{只有会说中文，才能申请这个岗位。} \\
\textbf{Analyse :} l'originale (\zh{只要}) affirme que le chinois \textbf{suffit} (aucun diplôme exigé). La transformée (\zh{只有}) fait du chinois une condition \textbf{indispensable} (mais un diplôme peut être exigé en plus). \textbf{Sens : plus restrictif.}

\textbf{3.} Transformation : \zh{只要缴纳会费，就能参加协会活动。} \\
\textbf{Analyse :} l'originale (\zh{只有}) rend le paiement \textbf{obligatoire} (mais un parrainage peut être exigé en plus). La transformée (\zh{只要}) affirme que le paiement \textbf{suffit} à lui seul. \textbf{Sens : plus large.}
\end{exoresolu}

\courssub{Culture de l'examen : le piège classique du baccalauréat}

\begin{savaistu}[QCM de grammaire : repérer 才 et 就]
Au baccalauréat (épreuve écrite comme à l'oral), les items sur \zh{只有} / \zh{只要} sont quasi systématiques. Voici la technique infaillible :

\begin{enumerate}
    \item \textbf{Repérez le deuxième mot-outil :} cherchez \zh{才} ou \zh{就} dans la phrase (souvent juste avant le verbe principal).
    \item \textbf{Déduisez le premier mot :}
    \begin{itemize}
        \item Si vous voyez \zh{才} $\rightarrow$ le début est obligatoirement \zh{只有}.
        \item Si vous voyez \zh{就} $\rightarrow$ le début est obligatoirement \zh{只要}.
    \end{itemize}
    \item \textbf{Vérifiez la logique :} s'agit-il d'une condition \emph{unique et indispensable} (\zh{只有}) ou \emph{simple et suffisante} (\zh{只要}) ?
\end{enumerate}

\textbf{Exemple type baccalauréat :} \\
\zh{\_\_\_\_\_\_ 你按时交作业，老师 \_\_\_\_\_\_ 会表扬你。} \\
A. 只有 … 就 \qquad B. 只要 … 才 \qquad C. 只有 … 才 \qquad D. 只要 … 就

\textbf{Résolution rapide :} les paires A et B sont des croisements interdits (\textbf{FAUX}). Restent C et D. Dans le contexte scolaire, rendre ses devoirs à temps \textbf{suffit} en général pour être félicité (condition suffisante) $\rightarrow$ \textbf{réponse D}. \\
(Si l'énoncé signifiait « ce n'est que si tu as 20/20 que le professeur te félicite », ce serait une condition nécessaire $\rightarrow$ réponse C.)
\end{savaistu}

\courssub{Entraînement : exercices d'application}

\begin{exercice}[Complétez avec 只有…才… ou 只要…就…]
\textbf{Consigne :} choisissez la structure adaptée au sens logique et placez \zh{才} ou \zh{就} à la bonne place.

1. \zh{\_\_\_\_\_\_\_\_ 你认真复习，\_\_\_\_\_\_\_\_ 考得好。} (Contexte : « Si tu révises sérieusement, tu auras de bonnes notes » — condition suffisante.)

2. \zh{\_\_\_\_\_\_\_\_ 你考满分，\_\_\_\_\_\_\_\_ 能得奖学金。} (Contexte : « Seule une note parfaite donne droit à la bourse » — condition nécessaire stricte.)

3. \zh{\_\_\_\_\_\_\_\_ 天气好，我们 \_\_\_\_\_\_\_\_ 去登山。} (Contexte : « Beau temps = on part en randonnée » — condition suffisante.)

4. \zh{\_\_\_\_\_\_\_\_ 办好签证，\_\_\_\_\_\_\_\_ 能去中国旅游。} (Contexte : « Il faut un visa pour aller en Chine » — condition administrative indispensable.)

5. \zh{\_\_\_\_\_\_\_\_ 我们团结一致，\_\_\_\_\_\_\_\_ 能赢得比赛。} (Contexte : discours de l'entraîneur, « l'unité suffit pour gagner » — condition suffisante.)

\textbf{Vocabulaire utile :} \zh{复习} (fùxí, réviser), \zh{满分} (mǎnfēn, note maximale), \zh{奖学金} (jiǎngxuéjīn, bourse d'études), \zh{登山} (dēngshān, faire de la randonnée en montagne), \zh{办签证} (bàn qiānzhèng, faire une demande de visa), \zh{团结一致} (tuánjié yīzhì, être unis et solidaires).
\end{exercice}

\begin{corrige}[Exercice — Complétion]
\textbf{1.} \zh{只要你认真复习，就考得好。} \\
(Zhǐyào nǐ rènzhēn fùxí, jiù kǎo de hǎo.) — Condition suffisante (réviser suffit). \zh{就} se place avant \zh{考}.

\textbf{2.} \zh{只有你考满分，才能得奖学金。} \\
(Zhǐyǒu nǐ kǎo mǎnfēn, cái néng dé jiǎngxuéjīn.) — Condition nécessaire (seule la note maximale donne droit à la bourse). \zh{才} se place avant \zh{能}.

\textbf{3.} \zh{只要天气好，我们就去登山。} \\
(Zhǐyào tiānqì hǎo, wǒmen jiù qù dēngshān.) — Condition suffisante (le beau temps suffit). \zh{就} suit le sujet \zh{我们} et précède \zh{去}.

\textbf{4.} \zh{只有办好签证，才能去中国旅游。} \\
(Zhǐyǒu bàn hǎo qiānzhèng, cái néng qù Zhōngguó lǚyóu.) — Condition nécessaire (le visa est obligatoire). \zh{才} se place avant \zh{能}. Note : \zh{办好} (bàn hǎo) signifie « mener à bien, obtenir ».

\textbf{5.} \zh{只要我们团结一致，就能赢得比赛。} \\
(Zhǐyào wǒmen tuánjié yīzhì, jiù néng yíngdé bǐsài.) — Condition suffisante (l'unité garantit la victoire, discours de motivation). \zh{就} suit le sujet \zh{我们} et précède \zh{能}.
\end{corrige}

\begin{exercice}[Expression écrite guidée : situation professionnelle]
\textbf{Consigne :} rédigez un court paragraphe (5 à 6 phrases) en chinois pour décrire les conditions d'embauche dans une entreprise fictive, la « \zh{中非科技有限公司} » (Zhōngfēi Kējì Yǒuxiàn Gōngsī, « Société de technologie sino-africaine SARL ») à Douala. Utilisez \textbf{au moins deux fois 只有…才…} et \textbf{deux fois 只要…就…}.

\textbf{Pistes :}
\begin{itemize}
    \item Diplôme et expérience (souvent avec \zh{只有…才…}).
    \item Compétences techniques ou linguistiques (peut être \zh{只要…就…} si c'est le critère unique).
    \item Période d'essai et signature du contrat (enchaînement logique avec \zh{只要…就…}).
\end{itemize}
\textbf{Mots-clés :} \zh{招聘} (zhāopìn, recruter), \zh{应聘者} (yìngpìnzhě, candidat), \zh{试用期} (shìyòngqī, période d'essai), \zh{转正} (zhuǎnzhèng, être titularisé), \zh{签合同} (qiān hétong, signer le contrat).
\end{exercice}

\begin{corrige}[Exercice — Expression écrite (modèle de réponse)]
\textbf{Production modèle :}

\zh{中非科技公司招聘工程师。只有拥有硕士学历和两年经验的应聘者，才能进入面试环节。只要通过技术面试，就能获得录用通知。只有签署劳动合同并缴纳保险，才能正式入职。只要试用期表现合格，就可以转正并享受全部福利。}

\textbf{Pinyin :} \\
Zhōngfēi Kējì Gōngsī zhāopìn gōngchéngshī. Zhǐyǒu yōngyǒu shuòshì xuélì hé liǎng nián jīngyàn de yìngpìnzhě, cái néng jìnrù miànshì huánjié. Zhǐyào tōngguò jìshù miànshì, jiù néng huòdé lùyòng tōngzhī. Zhǐyǒu qiānshǔ láodòng hétong bìng jiǎonà bǎoxiǎn, cái néng zhèngshì rùzhí. Zhǐyào shìyòngqī biǎoxiàn hégé, jiù kěyǐ zhuǎnzhèng bìng xiǎngshòu quánbù fúlì.

\textbf{Traduction :} \\
« La société de technologie sino-africaine recrute des ingénieurs. Seuls les candidats possédant un master et deux ans d'expérience peuvent accéder à l'étape de l'entretien. Du moment qu'on réussit l'entretien technique, on reçoit la notification d'embauche. Ce n'est qu'après avoir signé le contrat de travail et payé les assurances qu'on peut entrer officiellement en fonction. Tant que les résultats de la période d'essai sont satisfaisants, on peut être titularisé et bénéficier de l'ensemble des avantages. »

\textbf{Analyse des structures utilisées :}
\begin{itemize}
    \item \zh{只有拥有…才能进入…} : condition nécessaire (diplôme et expérience obligatoires).
    \item \zh{只要通过…就能获得…} : condition suffisante (réussir l'entretien technique suffit pour recevoir l'offre).
    \item \zh{只有签署…并缴纳…才能正式…} : condition nécessaire administrative (contrat et assurances indispensables).
    \item \zh{只要试用期…合格，就可以转正…} : condition suffisante (validation de l'essai = titularisation).
\end{itemize}
\end{corrige}

\coursec{Synthèse, entraînement et production}

Dans la section précédente, nous avons opposé \zh{只有……才……} (condition nécessaire) et \zh{只要……就……} (condition suffisante), et nous avons vu que croiser les deux paires est l'erreur la plus fréquente des francophones. Il est maintenant temps de rassembler les quatre connecteurs de la leçon dans un tableau de synthèse, de nous entraîner systématiquement, puis de produire un court texte professionnel qui les utilise tous.

\courssub{Tableau de synthèse des quatre structures}

\begin{aretenir}[Les quatre connecteurs de la leçon]
\begin{center}
\begin{tabularx}{\linewidth}{|X|X|X|X|}\hline
\rowcolor{popblueL} Structure & Sens logique & Équivalent français & Exemple (caractères, pinyin, traduction) \\ \hline
首先……然后……其次……最后…… & étapes successives & d'abord… ensuite… puis… enfin… & \zh{首先报名，然后考试，最后录取。} \\
& & & \emph{Shǒuxiān bàomíng, ránhòu kǎoshì, zuìhòu lùqǔ.} \\
& & & « D'abord s'inscrire, ensuite passer l'examen, enfin être admis. » \\ \hline
既……又…… & cumul de deux qualités & à la fois… et… ; non seulement… mais aussi… & \zh{这个工作既稳定又有意思。} \\
& & & \emph{Zhège gōngzuò jì wěndìng yòu yǒu yìsi.} \\
& & & « Ce travail est à la fois stable et intéressant. » \\ \hline
只有……才…… & condition nécessaire & ce n'est que… que… ; seulement si… & \zh{只有努力，才能成功。} \\
& & & \emph{Zhǐyǒu nǔlì, cái néng chénggōng.} \\
& & & « Ce n'est qu'en travaillant qu'on réussit. » \\ \hline
只要……就…… & condition suffisante & pourvu que… ; dès que… & \zh{只要努力，就会进步。} \\
& & & \emph{Zhǐyào nǔlì, jiù huì jìnbù.} \\
& & & « Pourvu qu'on travaille, on progresse. » \\ \hline
\end{tabularx}
\end{center}
\end{aretenir}

\begin{popfigure}
\begin{tikzpicture}[mindmap, grow cyclic, every node/.style={concept, align=center, font=\small},
root concept/.append style={concept color=popblue, font=\bfseries},
level 1/.append style={level distance=3.6cm, sibling angle=90}]
\node[root concept, align=center]{Connecteurs\\ZH31}
  child[concept color=poporange] { node[align=center]{首先…然后…\\其次…最后…\\ \emph{d'abord… enfin}\\ \textbf{étapes}} }
  child[concept color=popgreen] { node[align=center]{既…又…\\ \emph{à la fois… et}\\ \textbf{cumul}} }
  child[concept color=poppurple] { node[align=center]{只有…才…\\ \emph{ce n'est que… que}\\ \textbf{nécessaire}} }
  child[concept color=poppink] { node[align=center]{只要…就…\\ \emph{pourvu que}\\ \textbf{suffisant}} };
\end{tikzpicture}
\legende{Carte mentale de synthèse : les quatre connecteurs, leur schéma, leur équivalent français et leur mot-clé logique.}
\end{popfigure}

\begin{methode}[Réussir les exercices de connecteurs en trois étapes]
\begin{enumerate}
\item \textbf{Identifier le sens logique} de la phrase : s'agit-il d'étapes successives (首先/然后/其次/最后), d'un cumul de qualités (既……又), d'une condition nécessaire (只有……才) ou suffisante (只要……就) ?
\item \textbf{Choisir la paire complète} correspondant à ce sens, sans croiser les moitiés de deux paires différentes.
\item \textbf{Vérifier que les deux moitiés de la paire sont présentes} : relis ta phrase et souligne chaque élément du connecteur ; s'il en manque un, la phrase est fautive.
\end{enumerate}
\end{methode}

\begin{attention}[Bilan : les trois erreurs fréquentes des francophones]
\begin{itemize}
\item \textbf{Oublier le second élément de la paire} : *只有练习，能进步。 est fautif ; il faut \zh{只有练习，才能进步。} (Zhǐyǒu liànxí, cái néng jìnbù.) « Ce n'est qu'en s'entraînant qu'on progresse. »
\item \textbf{Croiser 只有 et 就} : *只有努力，就能成功。 mélange les deux conditions. Choisis : \zh{只有努力，才能成功。} (nécessaire) ou \zh{只要努力，就能成功。} (suffisant).
\item \textbf{Traduire « et » par 和 entre deux adjectifs ou verbes} : *这个工作有意思和工资高 est impossible. 和 ne relie que des noms ; entre adjectifs ou verbes, utilise \zh{既……又……} : \zh{这个工作既有意思，工资又高。}
\end{itemize}
\end{attention}

\courssub{Exercices d'application}

\begin{exoresolu}[Exercice 1 — Compléter avec 首先 / 然后 / 其次 / 最后]
Énoncé : complète le texte suivant avec les quatre connecteurs d'étapes.\\
\zh{……我要学好汉语，……通过HSK四级考试，……申请去中国留学，……在一家中国公司找到工作。}\\
\tcblower
Solution : le texte décrit quatre étapes successives d'un projet, dans l'ordre chronologique. On place donc les connecteurs dans l'ordre :\\
\zh{首先我要学好汉语，然后通过HSK四级考试，其次申请去中国留学，最后在一家中国公司找到工作。}\\
\emph{Shǒuxiān wǒ yào xuéhǎo Hànyǔ, ránhòu tōngguò HSK sì jí kǎoshì, qícì shēnqǐng qù Zhōngguó liúxué, zuìhòu zài yì jiā Zhōngguó gōngsī zhǎodào gōngzuò.}\\
« D'abord je veux bien apprendre le chinois, ensuite passer le HSK 4, puis demander à partir étudier en Chine, enfin trouver un emploi dans une entreprise chinoise. »
\end{exoresolu}

\begin{exercice}[Exercice 2 — Relier deux adjectifs avec 既……又……]
Transforme chaque paire en une phrase avec 既……又…… :
\begin{enumerate}
\item 这个工作 + 有意思 + 工资高
\item 雅温得 + 大 + 漂亮
\item 我的老师 + 严格 + 耐心
\item 汉语 + 难 + 有用
\end{enumerate}
\end{exercice}

\begin{exercice}[Exercice 3 — Transformation : 只要……就…… en 只有……才……]
Transforme les phrases suivantes et commente le changement de sens :
\begin{enumerate}
\item 只要多练习，你就会进步。
\item 只要我们一起努力，这个项目就能成功。
\end{enumerate}
\end{exercice}

\begin{exercice}[Exercice 4 — Correction de phrases fautives]
Corrige les phrases suivantes :
\begin{enumerate}
\item *只有练习，就能进步。
\item *这个工作有意思和工资高。
\item *首先我去学校，最后然后回家。
\item *既便宜又质量好，所以和我很喜欢。
\end{enumerate}
\end{exercice}

\begin{exercice}[Exercice 5 — Traduction (thème)]
Traduis en chinois :
\begin{enumerate}
\item « Ce n'est qu'en travaillant ensemble que nous réussirons. »
\item « Le Cameroun est à la fois beau et riche. »
\end{enumerate}
\end{exercice}

\begin{corrige}[Exercices 2 à 5]
\textbf{Exercice 2} :
\begin{enumerate}
\item \zh{这个工作既有意思，工资又高。} \emph{(Zhège gōngzuò jì yǒu yìsi, gōngzī yòu gāo.)} « Ce travail est à la fois intéressant et bien payé. »
\item \zh{雅温得既大又漂亮。} \emph{(Yǎwēndé jì dà yòu piàoliang.)} « Yaoundé est à la fois grande et belle. »
\item \zh{我的老师既严格又耐心。} \emph{(Wǒ de lǎoshī jì yángé yòu nàixīn.)} « Mon professeur est à la fois strict et patient. »
\item \zh{汉语既难又有用。} \emph{(Hànyǔ jì nán yòu yǒuyòng.)} « Le chinois est à la fois difficile et utile. »
\end{enumerate}

\textbf{Exercice 3} :
\begin{enumerate}
\item \zh{只有多练习，你才能进步。} \emph{(Zhǐyǒu duō liànxí, nǐ cái néng jìnbù.)}
\item \zh{只有我们一起努力，这个项目才能成功。} \emph{(Zhǐyǒu wǒmen yìqǐ nǔlì, zhège xiàngmù cái néng chénggōng.)}
\end{enumerate}
\textit{Commentaire :} avec \zh{只要……就}, la condition \emph{suffit} (d'autres moyens existent peut-être) ; avec \zh{只有……才}, elle devient \emph{indispensable} : sans elle, pas de résultat. Le message est plus exigeant.

\textbf{Exercice 4} :
\begin{enumerate}
\item \zh{只有练习，才能进步。} (ajout de \zh{才}, suppression de \zh{就})
\item \zh{这个工作既有意思，工资又高。} (remplacement de \zh{和} par \zh{既……又……})
\item \zh{首先我去学校，然后回家。} (suppression de \zh{最后}, ordre logique rétabli : deux étapes $\rightarrow$ 首先…然后…)
\item \zh{既便宜，质量又好，所以我很喜欢。} (parallélisme rétabli : adjectif / sujet-verbe ; suppression de \zh{和} inutile après \zh{所以})
\end{enumerate}

\textbf{Exercice 5} :
\begin{enumerate}
\item \zh{只有一起努力，我们才能成功。} \emph{(Zhǐyǒu yìqǐ nǔlì, wǒmen cái néng chénggōng.)}
\item \zh{喀麦隆既美丽又富饶。} \emph{(Kāmàilóng jì měilì yòu fùráo.)}
\end{enumerate}
\end{corrige}

\courssub{Production guidée : mon projet professionnel}

\begin{exemplebox}[Modèle de production]
\zh{首先我要高中毕业，然后学习汉语，其次去中国实习，最后在一家中国公司工作。只要努力，我的梦想就会实现。}\\
\emph{Shǒuxiān wǒ yào gāozhōng bìyè, ránhòu xuéxí Hànyǔ, qícì qù Zhōngguó shíxí, zuìhòu zài yì jiā Zhōngguó gōngsī gōngzuò. Zhǐyào nǔlì, wǒ de mèngxiǎng jiù huì shíxiàn.}\\
« D'abord je veux obtenir mon baccalauréat, ensuite apprendre le chinois, puis partir en stage en Chine, enfin travailler dans une entreprise chinoise. Pourvu que je travaille dur, mon rêve se réalisera. »
\end{exemplebox}

\begin{experience}[À toi d'écrire : les étapes de ton projet professionnel]
Rédige 5 à 6 phrases décrivant les étapes de ton projet professionnel (études, stage, emploi à Douala, Yaoundé ou en Chine…). Contraintes :
\begin{enumerate}
\item utilise au moins \textbf{trois structures} de la leçon ;
\item commence par une suite d'étapes avec 首先……然后……其次……最后…… ;
\item ajoute une phrase avec 既……又…… pour décrire le métier visé ;
\item termine par une condition avec 只要……就…… ou 只有……才…… ;
\item relis-toi avec la méthode en trois étapes : sens logique, paire choisie, deux moitiés présentes.
\end{enumerate}
\end{experience}

\begin{savaistu}[Un propos structuré, un exposé valorisé]
Dans les discours officiels chinois — discours de dirigeants, présentations d'entreprise, soutenances universitaires — la suite 首先……其次……最后…… est la marque d'un propos bien structuré et maîtrisé. L'auditeur chinois y est très sensible : annoncer clairement le plan de ce qu'on va dire est considéré comme un signe de rigueur et de respect envers l'auditoire. Au Cameroun, à l'oral du Bac de chinois, utiliser ces connecteurs dans ton exposé montre au jury que tu organises tes idées : c'est un critère de réussite directement valorisé dans la grille d'évaluation de l'expression orale.
\end{savaistu}

\newpage

\coursec{Activité d'intégration}

\coursec{Leçon 31 — Grammaire : Les connecteurs \zh{首先\ldots 然后\ldots 其次\ldots 最后\ldots} ; \zh{既\ldots 又\ldots} ; \zh{只有\ldots 才\ldots} ; \zh{只要\ldots 就\ldots}}

\courssub{Observation et découverte guidée}

\begin{textebox}[Dialogue : Préparer l'orientation — 规划未来]
\textbf{Contexte :} Dans le bureau d'orientation du lycée de Yaoundé, M. \textsc{Biya} conseille \textsc{Ngo}, élève de Terminale A, pour le concours de la Chambre de commerce Cameroun--Chine.\\[6pt]
\textbf{老师}：Ngo，你想好职业规划了吗？比赛要求\zh{首先}说目标，\zh{然后}说步骤，\zh{其次}说困难，\zh{最后}说决心。\\ 
\textbf{Lǎoshī}：Ngo, nǐ xiǎng hǎo zhíyè guīhuà le ma? Bǐsài yāoqiú \zh{shǒuxiān} shuō mùbiāo, \zh{ránhòu} shuō bùzhòu, \zh{qícì} shuō kùnnán, \zh{zuìhòu} shuō juéxīn.\\
« Ngo, as-tu réfléchi à ton plan de carrière ? Le concours exige : \textbf{d'abord} l'objectif, \textbf{ensuite} les étapes, \textbf{puis} les difficultés, \textbf{enfin} la détermination. »\\[6pt]
\textbf{Ngo}：我\zh{既}想当翻译，\zh{又}想做经理。可是\zh{只有}中文\zh{才}能当翻译，\zh{只要}我努力，\zh{就}能行。\\
\textbf{Ngo}：Wǒ \zh{jì} xiǎng dāng fānyì, \zh{yòu} xiǎng zuò jīnglǐ. Kěshì \zh{zhǐyǒu} Zhōngwén \zh{cái} néng dāng fānyì, \zh{zhǐyào} wǒ nǔlì, \zh{jiù} néng xíng.\\
« Je veux \textbf{à la fois} être traducteur \textbf{et} manager. Mais \textbf{ce n'est qu'avec} le chinois qu'\textbf{on peut} être traducteur ; \textbf{tant que} je m'efforce, \textbf{je pourrai} y arriver. »\\[6pt]
\textbf{老师}：很好！\zh{只有}逻辑清晰，\zh{才}能打动评委。\zh{只要}你用好四个连接词，\zh{就}有机会去杜阿拉实习。\\
\textbf{Lǎoshī}：Hěn hǎo! \zh{Zhǐyǒu} luóji qīngchu, \zh{cái} néng dǒngwéi wěiyuán. \zh{Zhǐyào} nǐ yòng hǎo sì ge liánjiēcí, \zh{jiù} yǒu jīhuì qù Dùālā shíxí.\\
« Très bien ! \textbf{Seule} une logique claire \textbf{peut} convaincre le jury. \textbf{Du moment que} tu utilises bien les quatre connecteurs, \textbf{tu auras} une chance de stage à Douala. »
\end{textebox}

\begin{experience}[Activité de découverte : Repérer les connecteurs]
\begin{enumerate}
\item Surligne dans le dialogue les quatre groupes de connecteurs étudiés aujourd'hui.
\item Pour chaque groupe, note : (a) sa place dans la phrase (début ? milieu ?) ; (b) le type de relation logique qu'il exprime (chronologie, addition, condition nécessaire, condition suffisante).
\item Compare avec le français : comment traduit-on \zh{才} et \zh{juste} ? Pourquoi \zh{就} ne se traduit-il pas toujours par « alors » ?
\end{enumerate}
\end{experience}

\courssub{\zh{首先\ldots 然后\ldots 其次\ldots 最后\ldots} : Organiser une suite d'étapes}

\begin{propriete}[Structure \zh{首先\ldots 然后\ldots 其次\ldots 最后\ldots}]
\textbf{Schéma :} \zh{首先}，S + V\ldots \zh{，然后}，S + V\ldots \zh{，其次}，S + V\ldots \zh{，最后}，S + V\ldots\\[4pt]
\textbf{Emploi :} Énumérer des actions ou des étapes dans un \textbf{ordre chronologique} ou \textbf{logique strict} (procédure, récit, plan, argumentation).\\[4pt]
\textbf{Règles de position :}
\begin{itemize}
\item Ces connecteurs se placent \textbf{en tête de proposition} (sujet + connecteur ou connecteur + sujet), suivis d'une virgule \zh{，}.
\item L'ordre est \textbf{figé} : \zh{首先} (premier) $\rightarrow$ \zh{然后} (deuxième / suite immédiate) $\rightarrow$ \zh{其次} (troisième / point suivant, souvent moins urgent) $\rightarrow$ \zh{最后} (final).
\item On ne saute pas d'étape (on ne commence pas par \zh{其次}).
\end{itemize}
\textbf{Exemples :}
\begin{center}
\begin{tabularx}{\linewidth}{|X|X|X|}
\hline
\rowcolor{popblueL} Contexte & Phrase chinoise (caractères + pinyin) & Traduction \\
\hline
Projet d'études & \zh{首先，我要通过高考。然后，上大学学中文。} \\ Shǒuxiān, wǒ yào tōngguò gāokǎo. Ránhòu, shàng dàxué xué Zhōngwén. & D'abord, je dois réussir le bac. Ensuite, entrer à l'université pour étudier le chinois. \\
\hline
Recrutement & \zh{首先看简历，然后面试，其次考察技能，最后决定录用。} \\ Shǒuxiān kàn jiǎnlì, ránhòu miànshì, qícì kǎochá jìnéng, zuìhòu juédìng lùyòng. & D'abord lire le CV, puis entretien, ensuite tester les compétences, enfin décider de l'embauche. \\
\hline
\end{tabularx}
\end{center}
\end{propriete}

\begin{attention}[Pièges francophones — Ordre et ponctuation]
\begin{itemize}
\item \textbf{Ordre inversé :} En français, on dit « ensuite, puis, enfin ». En chinois, \zh{其次} (qícì) vient \textbf{après} \zh{然后} (ránhòu) et \textbf{avant} \zh{最后} (zuìhòu). Ne confondez pas \zh{其次} (troisième étape) et \zh{然后} (deuxième étape).
\item \textbf{Virgule obligatoire :} \zh{首先，} \zh{然后，} \zh{其次，} \zh{最后，} sont toujours suivis d'une virgule chinoise \zh{，} avant le sujet ou le verbe.
\item \textbf{Sujet répété ou non :} \zh{首先我去，然后你看} (sujets différents) ou \zh{首先去，然后看} (même sujet ellipté). Les deux sont corrects.
\end{itemize}
\end{attention}

\courssub{\zh{既\ldots 又\ldots} : Cumuler deux caractéristiques}

\begin{propriete}[Structure \zh{既\ldots 又\ldots}]
\textbf{Schéma :} S + \zh{既} + Adj/V + \zh{又} + Adj/V\\[4pt]
\textbf{Emploi :} Exprimer la \textbf{coexistence} de deux qualités, états ou actions pour \textbf{un même sujet}. Sens : « à la fois... et... », « aussi bien... que... ».\\[4pt]
\textbf{Contraintes morphologiques :}
\begin{itemize}
\item Les deux éléments (X et Y) doivent être de \textbf{même nature grammaticale} : deux adjectifs (\zh{既高又大}), deux verbes (\zh{既会说又会写}), deux compléments (\zh{既吃米又吃面} — rare, préférer \zh{又\ldots 又\ldots} pour compléments).
\item Le sujet est \textbf{unique} et placé \textbf{avant} \zh{既}.
\item Négation : \zh{不既\ldots 又\ldots} (peu usité) ; on préfère \zh{既不\ldots 也不\ldots} (« ni... ni... »).
\end{itemize}
\textbf{Exemples :}
\begin{center}
\begin{tabularx}{\linewidth}{|X|X|X|}
\hline
\rowcolor{popblueL} Type & Phrase chinoise & Traduction \\
\hline
Adj + Adj & \zh{这份工作既轻松又赚钱。} Zhè fèn gōngzuò jì qīngsòng yòu zhuànqián. & Ce travail est à la fois facile et bien payé. \\
\hline
Verbe + Verbe & \zh{他既会说法语，又会说中文。} Tā jì huì shuō Fǎyǔ, yòu huì shuō Zhōngwén. & Il sait parler à la fois français et chinois. \\
\hline
Négation & \zh{这个决定既不公平，也不合理。} Zhège juédìng jì bù gōngpíng, yě bù hélǐ. & Cette décision n'est ni juste, ni raisonnable. \\
\hline
\end{tabularx}
\end{center}
\end{propriete}

\begin{attention}[Erreurs fréquentes — \zh{既\ldots 又\ldots} vs \zh{又\ldots 又\ldots}]
\begin{itemize}
\item \zh{既\ldots 又\ldots} : Style \textbf{soutenu / écrit}, insiste sur la \textbf{simultanéité} de deux attributs stables. Ex: \zh{既红又专} (à la fois rouge et expert / politiquement correct et professionnel).
\item \zh{又\ldots 又\ldots} : Style \textbf{courant / oral}, souvent pour des actions successives ou des états temporaires. Ex: \zh{他又累又饿} (il est à la fois fatigué et affamé \emph{maintenant}).
\item \textbf{Interdit :} Relier deux phrases complètes avec deux sujets différents. \ding{55} \zh{我既喜欢中文，他又喜欢法文。} \ding{51} \zh{我既喜欢中文，又喜欢法文。}
\end{itemize}
\end{attention}

\courssub{\zh{只有\ldots 才\ldots} : La condition nécessaire}

\begin{propriete}[Structure \zh{只有\ldots 才\ldots}]
\textbf{Schéma :} \zh{只有} + Condition (Sujet + Verbe) \zh{，才} + Résultat (Sujet + Verbe)\\[4pt]
\textbf{Sémantique :} \textbf{Condition nécessaire} (necessary condition). Le résultat ne se produit \textbf{que si} la condition est réalisée. « Ce n'est que si [Condition] que [Résultat] » / « Seulement si [Condition], [Résultat] ».\\[4pt]
\textbf{Points clés :}
\begin{itemize}
\item \zh{只有} (zhǐyǒu) introduit la proposition conditionnelle (souvent en tête de phrase).
\item \zh{才} (cái) est \textbf{obligatoire} dans la proposition principale, placé \textbf{immédiatement avant le verbe} (ou l'auxiliaire modal) du résultat.
\item Le sujet de la condition et du résultat est souvent le même (ellipse possible) mais peut être différent.
\end{itemize}
\textbf{Exemples :}
\begin{center}
\begin{tabularx}{\linewidth}{|X|X|X|}
\hline
\rowcolor{popblueL} Situation & Phrase chinoise & Traduction \\
\hline
Même sujet & \zh{只有努力学习，才能通过HSK4。} Zhǐyǒu nǔlì xuéxí, cái néng tōngguò HSK4. & Ce n'est qu'en étudiant dur qu'on peut réussir le HSK4. \\
\hline
Sujets différents & \zh{只有公司同意，我才能去出差。} Zhǐyǒu gōngsī tóngyì, wǒ cái néng qù chūchāi. & Ce n'est que si l'entreprise accepte que je pourrai partir en mission. \\
\hline
Avec \zh{不} & \zh{只有不放弃，才不会失败。} Zhǐyǒu bù fàngqì, cái bù huì shībài. & Tant qu'on n'abandonne pas, on n'échoue pas. \\
\hline
\end{tabularx}
\end{center}
\end{propriete}

\begin{aretenir}[Français vs Chinois : \zh{才} $\neq$ « alors »]
En français, « ce n'est que... que » structure la phrase (ex: \emph{Ce n'est qu'en travaillant que tu réussiras}). En chinois, \zh{只有} marque le début de la condition, et \zh{才} verrouille le résultat. \textbf{Ne jamais oublier \zh{才}}.
\begin{itemize}
\item \ding{55} \zh{只有你来，我去。} (Incomplet, sonne comme une simple succession)
\item \ding{51} \zh{只有你来，我\zh{才}去。} (Je n'irai \textbf{que si} tu viens).
\end{itemize}
\end{aretenir}

\courssub{\zh{只要\ldots 就\ldots} : La condition suffisante et comparaison}

\begin{propriete}[Structure \zh{只要\ldots 就\ldots}]
\textbf{Schéma :} \zh{只要} + Condition (Sujet + Verbe) \zh{，就} + Résultat (Sujet + Verbe)\\[4pt]
\textbf{Sémantique :} \textbf{Condition suffisante} (sufficient condition). La réalisation de la condition \textbf{entraîne automatiquement} le résultat. « Il suffit que... pour que... » / « Dès que... / Tant que..., ... ».\\[4pt]
\textbf{Points clés :}
\begin{itemize}
\item \zh{只要} (zhǐyào) = « seulement besoin de / du moment que ».
\item \zh{就} (jiù) est \textbf{obligatoire} dans la principale, placé \textbf{avant le verbe} (marque la conséquence immédiate ou inévitable).
\item Nuance : \zh{只要} implique souvent une facilité ou une généralité (« à chaque fois que... »).
\end{itemize}
\textbf{Exemples :}
\begin{center}
\begin{tabularx}{\linewidth}{|X|X|X|}
\hline
\rowcolor{popblueL} Situation & Phrase chinoise & Traduction \\
\hline
Conseil & \zh{只要你坚持练习，口语就会进步。} Zhǐyào nǐ jiānchí liànxí, kǒuyǔ jiù huì jìnbù. & Du moment que tu t'exerces régulièrement, l'oral progressera. \\
\hline
Règle générale & \zh{只要有护照，就可以申请签证。} Zhǐyào yǒu hùzhào, jiù kěyǐ shēnqǐng qiānzhèng. & Il suffit d'avoir un passeport pour demander un visa. \\
\hline
Condition minimale & \zh{只要五十块钱。} Zhǐyào wǔshí kuài qián. & Il ne faut que cinquante yuans / Cinquante yuans suffisent. (Ellipse de la principale) \\
\hline
\end{tabularx}
\end{center}
\end{propriete}

\begin{propriete}[Comparaison décisive : \zh{只有\ldots 才\ldots} vs \zh{只要\ldots 就\ldots}]
\begin{center}
\begin{tabularx}{\linewidth}{|X|X|X|}
\hline
\rowcolor{popblueL} Critère & \zh{只有\ldots 才\ldots} (Nécessaire) & \zh{只要\ldots 就\ldots} (Suffisante) \\
\hline
Logique & Condition \textbf{indispensable} : Sans A, pas de B. & Condition \textbf{suffisante} : Si A, alors B (mais B possible par d'autres moyens). \\
\hline
Mots-clés FR & « Ce n'est que si... que... », « Seulement si... » & « Il suffit que... », « Du moment que... », « Dès que... » \\
\hline
Marqueurs & \zh{才} (cái) : retard, difficulté, exclusivité. & \zh{就} (jiù) : rapidité, certitude, conséquence naturelle. \\
\hline
Exemple type & \zh{只有有钱，才能买房。} (L'argent est \textbf{la} condition sine qua non) & \zh{只要有钱，就能买房。} (L'argent \textbf{permet} d'acheter, mais il faut aussi des papiers, un bien dispo...) \\
\hline
\end{tabularx}
\end{center}
\end{propriete}

\begin{methode}[Comment choisir le bon connecteur conditionnel ?]
\begin{enumerate}
\item \textbf{Identifiez la relation logique :} Est-ce que A est \emph{la seule porte} vers B (\zh{只有\ldots{}才\ldots}) ou \emph{une porte large ouverte} vers B (\zh{只要\ldots{}就\ldots}) ?
\item \textbf{Testez par la négation :}
\begin{itemize}
\item « Sans A, B est impossible » $\rightarrow$ \zh{只有 A，才 B}.
\item « Avec A, B est garanti » $\rightarrow$ \zh{只要 A，就 B}.
\end{itemize}
\item \textbf{Placez \zh{才} / \zh{juste} avant le verbe principal :} \zh{才能去}, \zh{就可以去}. Jamais en fin de phrase.
\end{enumerate}
\end{methode}

\courssub{Synthèse, entraînement et production}

\courssub{Lexique du module 4 : Monde du travail (Révision + Extension)}

\begin{center}
\begin{tabularx}{\linewidth}{|X|X|X|X|}
\hline
\rowcolor{popblueL} Caractères & Pinyin & Nature & Traduction \\
\hline
职业规划 & zhíyè guīhuà & n. & plan de carrière, orientation professionnelle \\
\hline
高考 & gāokǎo & n. & examen national d'entrée à l'université (bac) \\
\hline
专业 & zhuānyè & n. & spécialité, filière, major \\
\hline
实习 & shíxí & v./n. & stage ; faire un stage \\
\hline
面试 & miànshì & v./n. & entretien d'embauche ; passer un entretien \\
\hline
简历 & jiǎnlì & n. & CV, curriculum vitæ \\
\hline
录用 & lùyòng & v. & embaucher, recruter \\
\hline
技能 & jìnéng & n. & compétence, savoir-faire \\
\hline
薪水 & xīnshuǐ & n. & salaire, paye \\
\hline
晋升 & jìnshēng & v. & promotion, avancement \\
\hline
中喀合资 & Zhōng-Kā hézī & n. & coentreprise sino-camerounaise \\
\hline
双语 & shuāngyǔ & adj. & bilingue \\
\hline
\end{tabularx}
\end{center}

\begin{exercice}[Entraînement 1 : Transformation de phrases]
Réécrivez chaque phrase en utilisant le connecteur imposé. Respectez l'ordre des mots et la place de \zh{才}/\zh{就}.
\begin{enumerate}
\item 我想当医生。我想当老师。 (\zh{既\ldots 又\ldots})
\item 你不努力，不能成功。 (\zh{只有\ldots 才\ldots})
\item 你会中文，你能在中喀公司工作。 (\zh{只要\ldots 就\ldots})
\item 我先回家，再吃饭，然后睡觉。 (\zh{首先\ldots 然后\ldots 最后\ldots} — adapter le nombre d'étapes)
\item Ce travail est fatigant. Ce travail est bien payé. (\zh{既\ldots 又\ldots})
\item Si tu as le visa, tu peux aller en Chine. (\zh{只要\ldots 就\ldots})
\item Ce n'est qu'en passant l'entretien qu'on est embauché. (\zh{只有\ldots 才\ldots})
\end{enumerate}
\end{exercice}

\begin{corrige}[Exercice 1]
\begin{enumerate}
\item \zh{我既想当医生，又想当老师。} (Wǒ jì xiǎng dāng yīshēng, yòu xiǎng dāng lǎoshī.)
\item \zh{只有努力，才能成功。} (Zhǐyǒu nǔlì, cái néng chénggōng.) — \emph{Sujet ellipté.}
\item \zh{只要你会中文，就能在中喀公司工作。} (Zhǐyào nǐ huì Zhōngwén, jiù néng zài Zhōng-Kā hézī gōngzuò.)
\item \zh{首先回家，然后吃饭，最后睡觉。} (Shǒuxiān huíjiā, ránhòu chīfàn, zuìhòu shuìjiào.) — \emph{3 étapes : on utilise 首先, 然后, 最后.}
\item \zh{这份工作既累又赚钱。} (Zhè fèn gōngzuò jì lèi yòu zhuànqián.)
\item \zh{只要有签证，就可以去中国。} (Zhǐyào yǒu qiānzhèng, jiù kěyǐ qù Zhōngguó.)
\item \zh{只有通过面试，才能被录用。} (Zhǐyǒu tōngguò miànshì, cái néng bèi lùyòng.) — \emph{Passif \zh{被} pour « être embauché ».}
\end{enumerate}
\end{corrige}

\begin{exercice}[Entraînement 2 : Complétion et choix du connecteur]
Complétez avec \zh{首先/然后/其次/最后}, \zh{既/又}, \zh{只有/才}, \zh{只要/就}.
\begin{enumerate}
\item \_\_\_\_\_\_ 你认真复习，\_\_\_\_\_\_ 考试没问题。
\item 这个项目 \_\_\_\_\_\_ 有创意，\_\_\_\_\_\_ 可行，老板才会批准。 (Deux adjectifs pour le sujet « projet »)
\item \_\_\_\_\_\_ 通过面试，\_\_\_\_\_\_ 能拿到这份工作。
\item 我的计划是：\_\_\_\_\_\_ 考上大学，\_\_\_\_\_\_ 学好专业，\_\_\_\_\_\_ 找实习，\_\_\_\_\_\_ 正式入职。
\end{enumerate}
\end{exercice}

\begin{corrige}[Exercice 2]
\begin{enumerate}
\item \zh{只要} 你认真复习，\zh{就} 考试没问题。 (Condition suffisante : réviser suffit pour réussir).
\item 这个项目 \zh{既} 有创意，\zh{又} 可行，老板才会批准。 (Cumul de qualités + condition nécessaire pour la validation : \zh{既\ldots{}又\ldots} pour le sujet « projet », \zh{只有\ldots{}才\ldots} pour la validation. \textbf{Attention :} La phrase mélange deux structures. Correction attendue : « Ce projet est à la fois créatif et réalisable. \textbf{Seulement ainsi} le patron l'approuvera. » $\rightarrow$ \zh{这个项目既有创意又可行，只有如此，老板才会批准。} Ou plus simple : \zh{这个项目既有创意又可行，只要这样，老板就会批准。}) \textbf{Version simplifiée pour l'exercice :} \zh{这个项目既有创意又可行。} (Focus sur \zh{既\ldots{}又\ldots}).
\item \zh{只有} 通过面试，\zh{才} 能拿到这份工作。 (Condition nécessaire).
\item \zh{首先} 考上大学，\zh{然后} 学好专业，\zh{其次} 找实习，\zh{最后} 正式入职。 (Ordre chronologique complet).
\end{enumerate}
\end{corrige}

\courssub{Production intégrée : Concours « Mon plan de carrière en 5 étapes »}

\begin{methode}[Rédiger et présenter son plan de carrière (Écrit + Oral)]
\begin{enumerate}
\item \textbf{Planifier (Brouillon) :} Notez 5 étapes réelles (Bac $\rightarrow$ Université/École $\rightarrow$ Stage/Spécialisation $\rightarrow$ Premier emploi $\rightarrow$ Objectif 5 ans). Choisis le vocabulaire du tableau.
\item \textbf{Structurer (Écrit) :} Rédigez 8--12 phrases.
\begin{itemize}
\item Ouvrez avec \zh{首先} (Étape 1 : Bac / Décision).
\item Développez avec \zh{然后} (Étape 2 : Études) et \zh{其次} (Étape 3 : Stage / Chine / Spécialisation).
\item Concluez avec \zh{最后} (Étape 4/5 : Emploi visé).
\item \textbf{Insérez obligatoirement :} 1$\times$ \zh{既\ldots{}又\ldots} (qualités du poste / compétences), 1$\times$ \zh{只有\ldots{}才\ldots} (condition sine qua non : ex. bilingue, diplôme), 1$\times$ \zh{只要\ldots{}就\ldots} (votre détermination).
\end{itemize}
\item \textbf{Vérifier (Auto-correction) :} \zh{才}/\zh{就} avant le verbe ? Virgules après connecteurs de séquence ? Sujet unique pour \zh{既\ldots{}又\ldots} ? Tons notés sur le brouillon ?
\item \textbf{Présenter (Oral 2 min) :} Ne lisez pas. Regardez la classe. Articuez les connecteurs (marqueurs de structure). La classe note sur grille : Tons / Structure / Vocabulaire / Fluidité.
\end{enumerate}
\end{methode}

\begin{exoresolu}[Production modèle : Plan de carrière de \textsc{Moukoko}, Tle A, Douala]
Énoncé : Rédiger un plan en 5 étapes avec les 4 connecteurs + lexique module 4.
\tcblower
\textbf{Texte :}\\[4pt]
我叫Moukoko，在杜阿拉读终末年。\\ Wǒ jiào Moukoko, zài Dùālā dú zhōngmò nián.\\ « Je m'appelle Moukoko, je suis en Terminale à Douala. »\\[4pt]
我的职业目标\zh{既}明确\zh{又}现实。\\ Wǒ de zhíyè mùbiāo \zh{jì} míngquè \zh{yòu} xiànshí.\\ « Mon objectif professionnel est à la fois clair et réaliste. »\\[4pt]
\zh{首先}，我要全力备战高考，考上雅温得大学。\\ \zh{Shǒuxiān}, wǒ yào quánlì bèizhàn gāokǎo, kǎoshàng Yǎwēndé dàxué.\\ « \textbf{D'abord}, je dois me préparer à fond pour le bac et entrer à l'université de Yaoundé. »\\[4pt]
\zh{然后}，我选择国际贸易专业，重点练习商务中文。\\ \zh{Ránhòu}, wǒ xuǎnzé guójì màoyì zhuānyè, zhòngdiǎn liànxí shāngwù Zhōngwén.\\ « \textbf{Ensuite}, je choisis la spécialité Commerce international, en insistant sur le chinois des affaires. »\\[4pt]
\zh{其次}，我申请去北京实习半年，积累公司经验。\\ \zh{Qícì}, wǒ shēnqǐng qù Běijīng shíxí bàn nián, jīlěi gōngsī jīngyàn.\\ « \textbf{Puis}, je postule pour un stage de six mois à Pékin pour accumuler de l'expérience en entreprise. »\\[4pt]
\zh{最后}，我回喀麦隆，在中喀贸易公司做经理助理。\\ \zh{Zuìhòu}, wǒ huí Kāmàilóng, zài Zhōng-Kā màoyì gōngsī zuò jīnglǐ zhùlǐ.\\ « \textbf{Enfin}, je rentre au Cameroun pour être assistant manager dans une société de commerce sino-camerounaise. »\\[4pt]
我知道，\zh{只有}掌握双语技能\zh{才}能胜任这个岗位。\\ Wǒ zhīdào, \zh{zhǐyǒu} zhǎngwò shuāngyǔ jìnéng \zh{cái} néng shèngrèn zhège gǎngwèi.\\ « Je sais que \textbf{ce n'est qu'en maîtrisant} les compétences bilingues \textbf{que je pourrai} assumer ce poste. »\\[4pt]
\zh{只要}我坚持不懈，\_\_\_\_\_\_ 就一定能实现梦想！\\ \zh{Zhǐyào} wǒ jiānchí bùxiè, \zh{jiù} yídìng néng shíxiàn mèngxiǎng!\\ « \textbf{Tant que} je persévère, \textbf{je réaliserai} sûrement mon rêve ! »\\[6pt]
\textbf{Analyse linguistique :}
\begin{itemize}
\item \zh{既明确又现实} : Adjectifs coordonnés, sujet unique (\zh{目标}).
\item Séquence \zh{首先/然后/其次/最后} : Respect strict de l'ordre, virgules, sujets elliptés (même sujet \zh{我}).
\item \zh{只有\ldots{}才\ldots} : Condition nécessaire (\zh{掌握双语技能}) pour le résultat (\zh{胜任}). \zh{才} devant \zh{能}.
\item \zh{只要\ldots{}就\ldots} : Condition suffisante (\zh{坚持不懈}) pour le résultat final (\zh{实现梦想}). \zh{就} devant \zh{能}.
\item Lexique réinvesti : \zh{职业目标}, \zh{国际贸易}, \zh{商务中文}, \zh{实习}, \zh{经验}, \zh{双语}, \zh{岗位}, \zh{中喀}.
\end{itemize}
\end{exoresolu}

\begin{experience}[Interaction orale : Questions croisées \zh{只有/才} vs \zh{只要/就}]
En binôme, posez à votre camarade deux questions sur son plan en utilisant les structures conditionnelles.
\begin{itemize}
\item \textbf{Modèle A (Nécessaire) :} \zh{你只有\_\_\_\_\_\_，才能\_\_\_\_\_\_吗？} (Ex: \zh{你只有拿到HSK5，才能去北京实习吗？})
\item \textbf{Modèle B (Suffisante) :} \zh{只要你\_\_\_\_\_\_，就能\_\_\_\_\_\_吗？} (Ex: \zh{只要你面试成功，就能被录用吗？})
\end{itemize}
Le camarade répond par une phrase complète avec le même connecteur.
\end{experience}

\begin{savaistu}[Le marché de l'emploi : Cameroun vs Chine]
\begin{itemize}
\item \textbf{Cameroun :} Le bilinguisme \zh{法语/英语} (Français/Anglais) est un atout majeur pour la fonction publique et les grandes entreprises. L'apprentissage du \zh{中文} (chinois) devient un \textbf{avantage concurrentiel décisif} (\zh{核心竞争力}) pour les postes dans les \zh{中资企业} (entreprises chinoises : China Railway, Sinohydro, Huawei, usines de transformation bois/agricoles à Kribi, Douala, Yaoundé).
\item \textbf{Chine :} Le marché valorise les \zh{海归} (retournés de l'étranger). Un Camerounais diplômé en Chine (\zh{留学生}) avec une spécialité technique (ingénierie, médecine, commerce, IA) et parlant chinois + français/anglais est très recherché pour les projets \zh{一带一路} (Belt and Road Initiative) en Afrique.
\item \textbf{Stages :} Les \zh{实习} (stages) sont la porte d'entrée principale. En Chine, les stages d'été (\zh{暑期实习}) sont très codifiés ; au Cameroun, les \zh{中喀合资企业} recrutent souvent via le \zh{职业规划比赛} (concours de plan de carrière) ou les forums \zh{校招} (recrutement sur campus).
\end{itemize}
\end{savaistu}

\courssub{Bilan de la leçon 31}

\begin{aretenir}[Check-list de révision — Suis-je prêt pour l'examen ?]
\begin{itemize}
\item[\ding{51}] Je sais ordonner \zh{首先 $\rightarrow$ 然后 $\rightarrow$ 其次 $\rightarrow$ 最后} sans erreur.
\item[\ding{51}] J'utilise \zh{既\ldots{}又\ldots} pour deux adjectifs/verbes de \textbf{même sujet}.
\item[\ding{51}] Je distingue \textbf{condition nécessaire} (\zh{只有\ldots{}才\ldots} = « ce n'est que si... que ») et \textbf{condition suffisante} (\zh{只要\ldots{}就\ldots} = « il suffit que... »).
\item[\ding{51}] Je place \zh{才} et \zh{就} \textbf{immédiatement avant le verbe} de la principale.
\item[\ding{51}] Je mobilise le lexique \zh{职业规划, 专业, 实习, 面试, 简历, 录用, 双语, 中喀合资} dans un texte cohérent.
\item[\ding{51}] Je peux présenter mon projet oralement en 2 min avec une prononciation tonale correcte.
\end{itemize}
\end{aretenir}

\begin{exercice}[Devoir maison / Évaluation formative]
Rédigez votre propre plan de carrière (10 phrases minimum) en respectant scrupuleusement la structure imposée (5 étapes + 3 connecteurs logiques). Soulignez les connecteurs. Préparez l'oral pour la prochaine séance.
\end{exercice}

\newpage

\coursec{Méthodes et exercices résolus}

\coursec{Leçon 31 — Grammaire : Les connecteurs 首先… 然后……其次……最后… ; 既…..又 ; 只有….才…. ; 只要……就…..}

\courssub{Observation et découverte guidée}

\begin{textebox}[Dialogue : Préparer un exposé sur l'économie camerounaise]
\textbf{Contexte :} Li Wei (étudiant chinois) et Paul (lycéen camerounais) préparent un exposé commun pour le Module 4.
\begin{enumerate}
\item \zh{保罗：我们要怎么安排这个 exposé？} \emph{Bǎoluó: Wǒmen yào zěnme ānpái zhège exposé?} « Paul : Comment allons-nous organiser cet exposé ? »
\item \zh{李伟：首先，我们介绍喀麦隆的农业。然后，我们谈工业。其次，分析中喀合作。最后，提出建议。} \emph{Lǐ Wěi: Shǒuxiān, wǒmen jièshào Kāmàilóng de nóngyè. Ránhòu, wǒmen tán gōngyè. Qícì, fēnxī Zhōng-Kā hézuò. Zuìhòu, tíchū jiànyì.} « Li Wei : D'abord, nous présentons l'agriculture du Cameroun. Ensuite, nous parlons de l'industrie. Puis, nous analysons la coopération sino-camerounaise. Enfin, nous formulons des recommandations. »
\item \zh{保罗：好主意！喀麦隆既有丰富的资源，又有年轻的人口。} \emph{Bǎoluó: Hǎo zhǔyì! Kāmàilóng jì yǒu fēngfù de zīyuán, yòu yǒu niánqīng de rénkǒu.} « Paul : Bonne idée ! Le Cameroun a à la fois des ressources abondantes et une population jeune. »
\item \zh{李伟：只有团结合作，才能实现共同发展。只要我们努力，就一定能做好。} \emph{Lǐ Wěi: Zhǐyǒu tuánjié hézuò, cái néng shíxiàn gòngtóng fāzhǎn. Zhǐyào wǒmen nǔlì, jiù yídìng néng zuò hǎo.} « Li Wei : Ce n'est qu'en s'unissant qu'on peut réaliser un développement commun. Pourvu que nous fassions des efforts, nous y arriverons sûrement. »
\end{enumerate}
\end{textebox}

\begin{aretenir}[Repérage des structures cibles]
Dans ce dialogue, identifiez et soulignez :
\begin{itemize}
\item Les quatre mots qui ordonnent le plan de l'exposé (\zh{首先、然后、其次、最后}).
\item La structure qui cumule deux atouts du Cameroun (\zh{既…又…}).
\item La structure qui exprime une condition \emph{indispensable} (\zh{只有…才…}).
\item La structure qui exprime une condition \emph{suffisante} (\zh{只要…就…}).
\end{itemize}
\end{aretenir}

\courssub{Vocabulaire de la leçon (HSK 3-4 — Module 4)}

\begin{center}
\begin{tabularx}{\linewidth}{|X|X|X|X|}
\hline
\rowcolor{popblueL} \textbf{Caractères} & \textbf{Pinyin} & \textbf{Nature} & \textbf{Traduction} \\ \hline
首先 & shǒuxiān & adv. & d'abord, en premier lieu \\ \hline
然后 & ránhòu & adv. & ensuite, alors \\ \hline
其次 & qícì & adv. & en second lieu, de plus \\ \hline
最后 & zuìhòu & adv. & enfin, en dernier lieu \\ \hline
既 & jì & adv. & à la fois (correlatif) \\ \hline
又 & yòu & adv. & aussi, de nouveau (correlatif) \\ \hline
只有 & zhǐyǒu & conj. & seulement si, le seul moyen est \\ \hline
才 & cái & adv. & seulement alors (marque condition nécessaire) \\ \hline
只要 & zhǐyào & conj. & pourvu que, il suffit que \\ \hline
就 & jiù & adv. & alors, dès lors (marque conséquence/suffisance) \\ \hline
安排 & ānpái & v. & organiser, arranger \\ \hline
介绍 & jièshào & v. & présenter, introduire \\ \hline
分析 & fēnxī & v. & analyser \\ \hline
提出 & tíchū & v. & formuler, soumettre (une idée) \\ \hline
建议 & jiànyì & n. & suggestion, recommandation \\ \hline
丰富 & fēngfù & adj. & abondant, riche \\ \hline
资源 & zīyuán & n. & ressource \\ \hline
团结 & tuánjié & v./adj. & s'unir, uni \\ \hline
合作 & hézuò & v./n. & coopérer, coopération \\ \hline
实现 & shíxiàn & v. & réaliser, concrétiser \\ \hline
共同 & gòngtóng & adj. & commun, partagé \\ \hline
发展 & fāzhǎn & v./n. & développer, développement \\ \hline
坚持 & jiānchí & v. & persévérer, persister \\ \hline
进步 & jìnbù & v./n. & progresser, progrès \\ \hline
通过 & tōngguò & v. & passer (examen), par le biais de \\ \hline
毕业 & bìyè & v. & obtenir son diplôme, finir ses études \\ \hline
\end{tabularx}
\end{center}

\courssub{Écriture des caractères clés : ordre des traits et radicaux}

\begin{definition}[Caractères à maîtriser pour l'écrit (Bac)]
\begin{itemize}
\item \zh{首} (shǒu, tête, premier) : Radiciel \zh{首} (tête). Ordre : haut → bas (nez, œil, cheveux), puis corps. 9 traits.
\item \zh{既} (jì, déjà, à la fois) : Radiciel \zh{无} (sans) + \zh{田} (champ) + \zh{儿} (enfant). Phonétique \zh{即} (jí). 9 traits. \textbf{Attention :} \zh{既} (jì, 4\textsuperscript{e} ton) $\neq$ \zh{即} (jí, 2\textsuperscript{e} ton, « immédiatement »).
\item \zh{才} (cái, talent / seulement alors) : Radiciel \zh{扌} (main) + \zh{才} (phonétique). Ordre : trait horizontal, vertical, crochet (main), puis \zh{才} (haut → bas). 7 traits. \textbf{Attention :} \zh{才} (cái, adv.) $\neq$ \zh{材} (cái, n. matériau, bois).
\item \zh{要} (yào, vouloir / il faut) : Radiciel \zh{覀} (couvrir) + \zh{女} (femme). 9 traits. Ton 4 (yào) seul ; ton 1 (yāo) dans \zh{如果} (rúguǒ) mais \zh{只要} se lit \zh{zhǐyào} (3-4).
\item \zh{则} (zé, règle / alors) : Radiciel \zh{刂} (sabre). 6 traits. Apparaît dans \zh{就} (jiù) : \zh{京} (capitale) + \zh{刂}. 12 traits.
\end{itemize}
\end{definition}

\courssub{Structure 1 : 首先… 然后… 其次… 最后… — Organiser une suite d'étapes}

\begin{propriete}[Schéma et emplois de la série chronologique]
\begin{center}
\begin{tabularx}{\linewidth}{|X|X|X|}
\hline
\rowcolor{popblueL} \textbf{Connecteur} & \textbf{Pinyin} & \textbf{Équivalent français / Emploi} \\ \hline
首先 & shǒuxiān & D'abord, premièrement — ouvre la séquence. \\ \hline
然后 & ránhòu & Ensuite, puis — 2\textsuperscript{e} étape (action successive). \\ \hline
其次 & qícì & En second lieu, de plus — 3\textsuperscript{e} étape (argument ou action supplémentaire). \\ \hline
最后 & zuìhòu & Enfin, finalement — clôture la séquence. \\ \hline
\end{tabularx}
\end{center}
\textbf{Structure canonique :} \textbf{Connecteur + Sujet + Verbe + Compléments。} \\
Le connecteur se place \emph{en tête de proposition} (sujet optionnel s'il est identique au précédent). Pas de virgule obligatoire après le connecteur si le sujet suit immédiatement, mais recommandée pour la clarté.
\end{propriete}

\begin{exemplebox}[Exemples variés]
\begin{enumerate}
\item \zh{首先，我们要分析题目。} \emph{Shǒuxiān, wǒmen yào fēnxī tímù.} (D'abord, il faut analyser le sujet.)
\item \zh{然后，查阅资料。} \emph{Ránhòu, chàyuè zīliào.} (Ensuite, consulter la documentation.)
\item \zh{其次，写提纲。} \emph{Qícì, xiě tígāng.} (Puis, rédiger le plan.)
\item \zh{最后，认真书写。} \emph{Zuìhòu, rènzhēn shūxiě.} (Enfin, écrire soigneusement.)
\end{enumerate}
\end{exemplebox}

\begin{methode}[Construire un texte en étapes avec 首先…然后…其次…最后…]
\begin{enumerate}
\item Découper le processus ou l'argumentation en 2 à 4 étapes logiques.
\item Associer chaque étape à son connecteur : \cle{首先} (1), \cle{然后} (2), \cle{其次} (3), \cle{最后} (4).
\item Placer le connecteur \emph{avant le sujet} (ou le verbe si sujet omis) de chaque phrase.
\item Vérifier la cohérence temporelle ou logique : \zh{然后} implique une succession temporelle ; \zh{其次} ajoute un point (souvent argumentatif).
\item Ne jamais placer le connecteur en fin de phrase (calque du français « d'abord » en fin).
\end{enumerate}
\end{methode}

\begin{exoresolu}[Rédiger un paragraphe : Préparer le ndolé]
\textbf{Énoncé.} Votre correspondant chinois vous demande la recette du ndolé. Rédigez 4 phrases avec les 4 connecteurs.\\[2pt]
\tcblower
\textbf{Raisonnement.} Étapes : 1. Laver feuilles. 2. Cuire avec arachides. 3. Ajouter viande/poisson. 4. Servir avec accompagnement.\\[4pt]
\textbf{Réponse.}\\
做ndolé很容易。首先，我们把ndolé的叶子洗干净。然后，把叶子和花生一起煮。其次，加上肉或者鱼。最后，配米饭或者miondo吃，味道好极了！\\
\emph{Zuò ndolé hěn róngyì. Shǒuxiān, wǒmen bǎ ndolé de yèzi xǐ gānjìng. Ránhòu, bǎ yèzi hé huāshēng yìqǐ zhǔ. Qícì, jiā shang ròu huòzhě yú. Zuìhòu, pèi mǐfàn huòzhě miondo chī, wèidao hǎo jí le!}\\
« Préparer le ndolé est facile. D'abord, nous lavons les feuilles de ndolé. Ensuite, on cuit les feuilles avec les arachides. Puis, on ajoute viande ou poisson. Enfin, on mange avec riz ou miondo, c'est délicieux ! »\\[4pt]
\textbf{Justifications.} Structure \zh{把} pour l'objet antéposé (laver/cuire) ; \zh{好极了} exclamation. Connecteurs en tête de phrase.
\end{exoresolu}

\begin{exercice}[Entraînement : Ordonner un récit]
Remettez les phrases dans l'ordre logique en insérant les connecteurs appropriés (首/然/次/最) :\\
1. \_\_\_\_，我们到达了北京。 2. \_\_\_\_，我们参观了故宫。 3. \_\_\_\_，我们坐飞机出发。 4. \_\_\_\_，我们品尝了烤鸭。
\end{exercice}

\begin{corrige}[Exercice : Ordonner un récit]
1. \textbf{首先}，我们坐飞机出发。 (Départ)
2. \textbf{然后}，我们到达了北京。 (Arrivée)
3. \textbf{其次}，我们参观了故宫。 (Activité 1)
4. \textbf{最后}，我们品尝了烤鸭。 (Activité finale / repas)
\emph{Ordre chronologique strict respecté.}
\end{corrige}

\courssub{Structure 2 : 既…又… — Cumuler deux caractéristiques}

\begin{propriete}[Schéma, contraintes et emplois de 既…又…]
\textbf{Structure :} \textbf{Sujet + 既 + Adj/Verbe 1 + 又 + Adj/Verbe 2。} \\
\textbf{Contraintes strictes :}
\begin{enumerate}
\item \textbf{Même sujet} pour les deux prédicats.
\item \textbf{Même nature grammaticale} : Adj + Adj, Verbe + Verbe, Verbe+Objet + Verbe+Objet.
\item \textbf{Pas de mot de liaison} (pas de \zh{和}, \zh{也}, \zh{还} entre les deux).
\item \textbf{Négation :} \zh{不既…又…} (peu usité) ; préférer \zh{既不…也不…} (ni… ni…) au niveau supérieur. Ici on retient l'affirmatif.
\end{enumerate}
\end{propriete}

\begin{exemplebox}[Exemples par catégorie grammaticale]
\begin{enumerate}
\item \textbf{Adjectifs :} \zh{这道菜既好吃又不贵。} \emph{Zhè dào cài jì hǎochī yòu bù guì.} (Ce plat est à la fois bon et pas cher.)
\item \textbf{Verbes d'état :} \zh{他既聪明又勤奋。} \emph{Tā jì cōngmíng yòu qínfèn.} (Il est à la fois intelligent et travailleur.)
\item \textbf{Verbes d'action (même verbe, objets différents) :} \zh{他既会说法语，又会说英语。} \emph{Tā jì huì shuō Fǎyǔ, yòu huì shuō Yīngyǔ.} (Il sait parler à la fois français et anglais.) $\rightarrow$ Répéter le verbe si objets différents.
\item \textbf{Verbes d'action (verbes différents) :} \zh{周末我既要复习，又要打球。} \emph{Zhōumò wǒ jì yào fùxí, yòu yào dǎqiú.} (Le week-end je dois à la fois réviser et jouer au ballon.)
\end{enumerate}
\end{exemplebox}

\begin{methode}[Transformer deux phrases en une avec 既…又…]
\begin{enumerate}
\item Vérifier : Sujet identique ? Nature identique (Adj/Adj ou V/V) ?
\item Écrire : Sujet + \zh{既} + Prédicat 1 + \zh{又} + Prédicat 2.
\item Si Prédicat = Verbe + Objet différent $\rightarrow$ Répéter le Verbe après \zh{又}.
\item Supprimer l'adverbe de degré \zh{很} (souvent implicite dans 既…又…).
\end{enumerate}
\end{methode}

\begin{exoresolu}[Transformation et pièges]
\textbf{Énoncé.} Fusionnez :\\
1. 杜阿拉很大。杜阿拉很热闹。\\
2. 我喜欢吃ndolé。我喜欢吃miondo。\\
3. 这个工作很累。这个工作很危险。 (Attention : nature ?)\\[2pt]
\tcblower
\textbf{Réponses.}\\
1. 杜阿拉既大又热闹。 \emph{Dù'ālā jì dà yòu rènao.} (Adj + Adj, pas de \zh{很}.)
2. 我既喜欢吃ndolé，又喜欢吃miondo。 \emph{Wǒ jì xǐhuan chī ndolé, yòu xǐhuan chī miondo.} (V+O différents $\rightarrow$ répéter \zh{喜欢吃}.)
3. 这个工作既累又危险。 \emph{Zhège gōngzuò jì lèi yòu wēixiǎn.} (Adj + Adj : \zh{累} et \zh{危险} sont tous deux adjectifs prédicatifs.)
\end{exoresolu}

\begin{exercice}[Entraînement : 既…又…]
Transformez en une phrase avec 既…又… :\\
1. 我的汉语老师很严格。我的汉语老师很幽默。\\
2. 我们要学好语法。我们要练习口语。\\
3. 这件衣服很贵。这件衣服很旧。 (Piège : sens péjoratif possible)
\end{exercice}

\begin{corrige}[Exercice : 既…又…]
1. 我的汉语老师既严格又幽默。\\
2. 我们既要学好语法，又要练习口语。 (Répétition de \zh{要} car compléments différents.)\\
3. 这件衣服既贵又旧。 (Structure correcte, sens : « à la fois chère et vieille ».)
\end{corrige}

\courssub{Structure 3 : 只有…才… — Condition nécessaire (exclusive)}

\begin{propriete}[Schéma, place de 才 et nuance]
\textbf{Structure :} \textbf{只有 + Condition，+ (Sujet) + 才 + Résultat。} \\
\begin{center}
\begin{tikzpicture}[node distance=1.2cm, >=Stealth]
\node[draw, rounded corners, fill=popblueL, text width=3cm, align=center] (cond) {只有 + Condition};
\node[draw, rounded corners, fill=poporangeL, text width=4cm, align=center, right of=cond] (res) {Sujet + 才 + Verbe + Complément};
\draw[->, thick, popblue] (cond) -- (res) node[midway, above, font=\small] {Condition INDISPENSABLE};
\end{tikzpicture}
\end{center}
\textbf{Règles d'or :}
\begin{itemize}
\item \zh{才} se place \textbf{immédiatement avant le verbe} (ou l'auxiliaire modal \zh{能、会、可以、要}).
\item Sujet : peut être avant \zh{才} (\zh{你才能去}) ou après \zh{只有} condition si différent (\zh{只有你努力，才能成功} — sujet générique).
\item \textbf{Interdiction formelle :} \zh{只有…就…} (faute grave).
\item Sens : « Ce n'est que si [Condition] que [Résultat] » $\rightarrow$ Sans condition, résultat impossible.
\end{itemize}
\end{propriete}

\begin{exemplebox}[Variantes de position du sujet]
\begin{enumerate}
\item \zh{只有努力，才能成功。} \emph{Zhǐyǒu nǔlì, cái néng chénggōng.} (Sujet générique omis.)
\item \zh{只有你努力，你才能成功。} \emph{Zhǐyǒu nǐ nǔlì, nǐ cái néng chénggōng.} (Sujet répété.)
\item \zh{只有你努力，才能成功。} \emph{Zhǐyǒu nǐ nǔlì, cái néng chénggōng.} (Sujet unique après condition, sous-entendu avant \zh{才}.)
\end{enumerate}
\end{exemplebox}

\begin{methode}[Traduire « Ce n'est que… que… » / « Seulement si… »]
\begin{enumerate}
\item Repérer la structure française « Ce n'est que [X] que [Y] » ou « Seulement si [X], [Y] ».
\item Traduire [X] après \zh{只有}.
\item Traduire [Y] en plaçant \zh{才} devant le verbe.
\item Vérifier : pas de \zh{就}, pas de \zh{则}.
\end{enumerate}
\end{methode}

\begin{exoresolu}[Complétion et traduction]
\textbf{Énoncé.} 1. Complétez : 只有每天练习汉字，你\_\_\_记得住。\\
2. Traduisez : « Ce n'est qu'en passant le Bac que vous entrerez à l'université. » (Bac = 高考 / 会考 ; entrer à l'université = 上大学).\\[2pt]
\tcblower
\textbf{Réponses.}\\
1. 只有每天练习汉字，你\textbf{才}记得住。 \emph{Zhǐyǒu měitiān liànxí Hànzì, nǐ cái jìde zhù.}\\
2. 只有通过会考，你们才能上大学。 \emph{Zhǐyǒu tōngguò huìkǎo, nǐmen cái néng shàng dàxué.}
\textbf{Justification.} \zh{才} précède \zh{记得住} / \zh{能上}. \zh{只有} marque l'unicité de la voie.
\end{exoresolu}

\begin{exercice}[Entraînement : 只有…才…]
Complétez par 才 ou 就 (un seul choix correct) et traduisez :\\
1. 只有认真听课，你\_\_\_能懂语法。\\
2. 只有明天不下雨，我们\_\_\_去郊游。\\
3. Traduisez : « Ce n'est qu'en respectant les règles qu'on gagne le match. » (respecter = 遵守 ; règles = 规则 ; gagner = 赢 ; match = 比赛).
\end{exercice}

\begin{corrige}[Exercice : 只有…才…]
1. \textbf{才} (condition nécessaire). « Ce n'est qu'en écoutant attentivement que tu pourras comprendre la grammaire. »
2. \textbf{才} (condition nécessaire). « Ce n'est que s'il ne pleut pas demain que nous irons en excursion. »
3. 只有遵守规则，才能赢比赛。 \emph{Zhǐyǒu zūnshǒu guīzé, cái néng yíng bǐsài.}
\end{corrige}

\courssub{Structure 4 : 只要…就… — Condition suffisante (ouverte)}

\begin{propriete}[Schéma, place de 就 et opposition à 只有…才…]
\textbf{Structure :} \textbf{只要 + Condition，+ (Sujet) + 就 + Résultat。} \\
\begin{center}
\begin{tabularx}{\linewidth}{|X|X|X|}
\hline
\rowcolor{popblueL} \textbf{Critère} & \textbf{只有…才…} & \textbf{只要…就…} \\ \hline
\textbf{Type de condition} & Nécessaire (exclusive) & Suffisante (ouverte) \\ \hline
\textbf{Marqueur français} & « Ce n'est que… que… », « Seulement si… » & « Pourvu que… », « Il suffit que… », « Si… (alors) » \\ \hline
\textbf{Corrélatif} & 才 (cái) & 就 (jiù) \\ \hline
\textbf{Autres moyens ?} & Non, impossible autrement & Oui, possibles \\ \hline
\end{tabularx}
\end{center}
\textbf{Règle :} \zh{就} se place avant le verbe/auxiliaire du résultat. \textbf{Interdiction :} \zh{只要…才…}.
\end{propriete}

\begin{exemplebox}[Exemples concrets]
\begin{enumerate}
\item \zh{只要你想学，我就教你。} \emph{Zhǐyào nǐ xiǎng xué, wǒ jiù jiāo nǐ.} (Pourvu que tu veuilles apprendre, je t'enseigne.)
\item \zh{只要有时间，我就去图书馆。} \emph{Zhǐyào yǒu shíjiān, wǒ jiù qù túshūguǎn.} (Dès que j'ai du temps, j vais à la biblio — habitude.)
\item \zh{只要坚持不懈，就没有做不成的事。} \emph{Zhǐyào jiānchí bùxiè, jiù méiyǒu zuò bù chéng de shì.} (Pourvu qu'on persévère, rien n'est impossible.)
\end{enumerate}
\end{exemplebox}

\begin{methode}[Choisir entre 只有…才… et 只要…就… : Test du « Seule voie »]
\begin{enumerate}
\item Posez la question : « Est-ce le \emph{seul} moyen d'obtenir le résultat ? »
\item OUI $\rightarrow$ \zh{只有…才…} (Ex: passer l'examen pour le diplôme).
\item NON (c'est \emph{un} moyen parmi d'autres, ou une condition facile) $\rightarrow$ \zh{只要…就…} (Ex: avoir le temps pour aller au marché).
\item Attention aux phrases hypothétiques : \zh{如果…就…} (si… alors…) est neutre ; \zh{只要…就…} insiste sur la \emph{facilité} ou la \emph{suffisance} de la condition.
\end{enumerate}
\end{methode}

\begin{exoresolu}[Choix discriminé et justification]
\textbf{Énoncé.} Choisissez la structure correcte et justifiez :\\
1. \_\_\_\_你努力，你\_\_\_考上大学。 (Seule voie = Bac/Concours obligatoire)\\
2. \_\_\_\_你有空，我们\_\_\_喝茶。 (Condition simple, suffisant pour le plaisir)\\
3. \_\_\_\_通过面试，你\_\_\_得到这份工作。 (Seule voie administrative)\\[2pt]
\tcblower
\textbf{Réponses.}\\
1. \textbf{只有}你努力，你\textbf{才}考上大学。 (Nécessaire : sans effort, impossible.)
2. \textbf{只要}你有空，我们\textbf{就}喝茶。 (Suffisant : avoir du temps suffit à déclencher l'invitation.)
3. \textbf{只有}通过面试，你\textbf{才}能得到这份工作。 (Nécessaire : l'entretien est étape obligatoire.)
\end{exoresolu}

\begin{exercice}[Entraînement : 只要…就… vs 只有…才…]
Complétez avec la bonne paire (只要…就 / 只有…才) :\\
1. \_\_\_\_你按时交作业，老师\_\_\_不批评你。\\
2. \_\_\_\_考上大学，你\_\_\_能拿学位证。\\
3. \_\_\_\_下雨，地\_\_\_湿。 (Loi naturelle / condition suffisante universelle)
\end{exercice}

\begin{corrige}[Exercice : Choix conditionnel]
1. \textbf{只要}你按时交作业，老师\textbf{就}不批评你。 (Condition suffisante pour éviter la critique.)
2. \textbf{只有}考上大学，你\textbf{才}能拿学位证。 (Condition nécessaire : pas de diplôme sans entrée à l'univ.)
3. \textbf{只要}下雨，地\textbf{就}湿. (Condition suffisante universelle. \zh{如果}下雨也 possible mais \zh{只要} insiste sur l'automaticité.)
\end{corrige}

\courssub{Synthèse comparative, pièges et production}

\begin{propriete}[Tableau récapitulatif des 4 structures (Niveau Bac)]
\begin{center}
\begin{tabularx}{\linewidth}{|X|X|X|X|}
\hline
\rowcolor{popblueL} \textbf{Structure} & \textbf{Fonction} & \textbf{Corrélatifs} & \textbf{Ordre mots-clé} \\ \hline
首先 / 然后 / 其次 / 最后 & Énumération chronologique / logique & — & Connecteur + Proposition \\ \hline
既…又… & Cumul de qualités/actions (même sujet, même nature) & 既 — 又 & Suj + 既 + Préd1 + 又 + Préd2 \\ \hline
只有…才… & Condition \textbf{nécessaire} (exclusive) & 只有 — 才 & 只有 + Cond, (Suj) + 才 + Rés \\ \hline
只要…就… & Condition \textbf{suffisante} (ouverte) & 只要 — 就 & 只要 + Cond, (Suj) + 就 + Rés \\ \hline
\end{tabularx}
\end{center}
\end{propriete}

\begin{attention}[Les 5 erreurs fatales au Bac (Rappel et approfondissement)]
\begin{enumerate}
\item \textbf{Croisement des corrélatifs :} \zh{只有…就…} $\times$ ; \zh{只要…才…} $\times$. \textbf{Mnémotechnique :} \zh{只} (seul) appelle \zh{才} (seulement) ; \zh{要} (demande) appelle \zh{就} (alors/immédiat).
\item \textbf{Mauvaise place de 才 / 就 :} Ils sont des \emph{adverbes préverbaux}. \zh{他才会来} (Juste) ; \zh{才他会来} $\times$ ; \zh{他会才来} $\times$.
\item \textbf{Intrus dans 既…又… :} \zh{既高又和壮} $\times$ (pas de \zh{和}) ; \zh{既高也壮} $\times$ (pas de \zh{也}).
\item \textbf{Connecteur d'étape en position finale :} \zh{我们洗叶子首先。} $\times$. Toujours en tête : \zh{首先，我们洗叶子。}
\item \textbf{Confusion graphique / tonale :} \zh{既} (jì, 4\textsuperscript{e}) vs \zh{即} (jí, 2\textsuperscript{e}) ; \zh{才} (cái) vs \zh{材} (cái) ; \zh{只} (zhǐ, 3\textsuperscript{e} dans \zh{只要}) vs \zh{只} (zhī, 1\textsuperscript{e} classifieur). \textbf{Sanction double :} Grammaire + Écriture.
\item \textbf{Bonus :} Oublier le sujet après \zh{只有} si changement de sujet. \zh{只有你去，我才去。} (Correct) vs \zh{只有去，我才去。} (Ambigu : qui va ?).
\end{enumerate}
\end{attention}

\courssub{Production écrite guidée : Paragraphe argumenté (Type Bac)}

\begin{methode}[Rédiger un paragraphe cohérent HSK 3-4 — 5 à 6 phrases]
\begin{enumerate}
\item \textbf{Phrase d'amorce :} Thème + \zh{既…又…} (deux qualités du sujet).
\item \textbf{Question transition :} \zh{怎么才能…呢？} / \zh{如何…？}
\item \textbf{Développement ordonné :} \zh{首先…}，\zh{其次…}，\zh{最后…} (3 arguments/actions).
\item \textbf{Conclusion conditionnelle :} \zh{只要…就…} (encouragement) OU \zh{只有…才…} (exigence).
\item \textbf{Relecture :} 1 connecteur/phrase, 才/就 préverbaux, pas de \zh{和} dans 既…又…, caractères exacts.
\end{enumerate}
\end{methode}

\begin{exoresolu}[Composition : « Comment réussir son orientation professionnelle ? »]
\textbf{Énoncé.} Rédigez 6 phrases : 既…又… (qualités du métier), 首先/其次/最后 (3 étapes), 只有…才… (condition nécessaire).\\[2pt]
\tcblower
\textbf{Plan :} Métier idéal = 既稳定又有前途. Étapes : 1. S'informer (首先). 2. Choisir filière (其次). 3. Travailler dur (最后). Condition : 只有努力才能成功.\\[4pt]
\textbf{Réponse.}\\
一个理想的职业既稳定又有前途。怎么选择适合自己的路呢？首先，我们要多了解各种专业的情况。其次，要根据自己的兴趣和特长决定。最后，必须在学校里刻苦学习。只有脚踏实地，才能实现职业梦想。\\
\emph{Yí ge lǐxiǎng de zhíyè jì wěndìng yòu yǒu qiántú. Zěnme xuǎnzé shìhé zìjǐ de lù ne? Shǒuxiān, wǒmen yào duō liǎojiě gè zhǒng zhuānyè de qíngkuàng. Qícì, yào gēnjù zìjǐ de xìngqù hé tècháng juédìng. Zuìhòu, bìxū zài xuéxiào lǐ kěkǔ xuéxí. Zhǐyǒu jiǎotāshídì, cái néng shíxiàn zhíyè mèngxiǎng.}\\
« Un métier idéal est à la fois stable et prometteur. Comment choisir sa voie ? D'abord, s'informer sur les filières. Ensuite, décider selon ses intérêts et talents. Enfin, étudier sérieusement à l'école. Ce n'est qu'en étant pragmatique (pieds sur terre) qu'on réalise son rêve professionnel. »\\[4pt]
\textbf{Justifications.} \zh{既稳定又有前途} (Adj + V-O, nature proche acceptée) ; \zh{脚踏实地} (chengyu, « les pieds sur terre ») ; \zh{才能} ordre strict.
\end{exoresolu}

\begin{experience}[Activité orale : Débat « Étudier en Chine ou au Cameroun ? »]
\textbf{Consigne :} Par binômes. L'un défend la Chine, l'autre le Cameroun pour les études supérieures.
\textbf{Contraintes :} Utiliser au moins 1 fois : 既…又… (atouts du pays), 首先/其次/最后 (arguments), 只有…才… (condition réussite).
\textbf{Exemple amorce :} \zh{中国既有名校又有奖学金。首先… 其次… 最后… 只有努力，才能毕业。}
\textbf{Évaluation :} Fluidité, justesse structures, vocabulaire Module 4 (专业、奖学金、学费、文化、适应).
\end{experience}

\begin{exercice}[Production écrite individuelle (Devoir / Bac blanc)]
\textbf{Sujet :} « Pour développer l'économie du Cameroun, que faut-il faire ? »\\
Rédigez un paragraphe de 80-100 caractères chinois. Obligatoire : 既…又… (richesses du pays), 首先/然后/最后 (3 mesures), 只要…就… (optimisme).
\textbf{Mots-clés fournis :} 资源、农业、工业、投资、基础设施、发展、繁荣.
\end{exercice}

\begin{corrige}[Exercice : Développement économie]
\textbf{Proposition de corrigé :}\\
喀麦隆既有丰富的自然资源，又有年轻的劳动力。首先，要发展现代农业。然后，要建设工业基础设施。最后，要吸引外国投资。只要政府和人民团结一心，国家就一定会繁荣昌盛。\\
\emph{Kāmàilóng jì yǒu fēngfù de zìrán zīyuán, yòu yǒu niánqīng de láodònglì. Shǒuxiān, yào fāzhǎn xiàndài nóngyè. Ránhòu, yào jiànshè gōngyè jīchǔ shèshī. Zuìhòu, yào xīyǐn wàiguó tóuzī. Zhǐyào zhèngfǔ hé rénmín tuánjié yīxīn, guójiā jiù yídìng huì fánróng chāngshèng.}
\end{corrige}

\courssub{Élément socioculturel : Logique du discours et monde du travail}

\begin{savaistu}[Structurer sa pensée : Approche chinoise vs Occidentale / Camerounaise]
\begin{itemize}
\item \textbf{Chine (写作/汇报) :} Structure \emph{déductive} et \emph{hiérarchique} très codifiée. L'usage de \zh{首先/其次/最后} est la norme attendue dans les rapports administratifs, les exposés scolaires (高考作文) et les discours officiels. Cela montre la \zh{条理} (logique, clarté d'esprit).
\item \textbf{Cameroun / Francophonie :} Structure souvent \emph{inductive} (arguments puis thèse) ou \emph{dialectique} (thèse/antithèse/synthèse). Les connecteurs « d'abord, ensuite, enfin » existent mais sont moins ritualisés qu'en chinois écrit formel.
\item \textbf{Conseil pour l'examen :} En production écrite (Bac), l'absence de ces 4 connecteurs pour lister 3 arguments ou plus est perçue comme un manque de \zh{条理} et coûte des points sur la note de « Structure/Coherence ». Utilisez-les \textbf{systématiquement}.
\item \textbf{Monde du travail :} Dans une entreprise chinoise (ou sino-camerounaise), un rapport d'activité (\zh{工作汇报}) sans \zh{首先/其次/最后} paraîtra brouillon. Maîtriser ces connecteurs, c'est acquérir une \textbf{compétence professionnelle transversale}.
\end{itemize}
\end{savaistu}

\newpage

\coursec{Exercices}

\courssub{Je vérifie mes connaissances}

\begin{exercice}[QCM : choisir le bon connecteur]
Pour chaque phrase, choisis la réponse correcte (A, B, C ou D).

\begin{enumerate}
\item 只有努力学习，\zh{＿＿} 能考上好大学。\\
\emph{Zhǐyǒu nǔlì xuéxí, \_\_\_ néng kǎoshang hǎo dàxué.}\\
A. 就 \quad B. 才 \quad C. 也 \quad D. 都
\item 只要天气好，我们 \zh{＿＿} 去踢足球。\\
\emph{Zhǐyào tiānqì hǎo, wǒmen \_\_\_ qù tī zúqiú.}\\
A. 才 \quad B. 就 \quad C. 只 \quad D. 又
\item 这个工作 \zh{＿＿} 工资高，\zh{＿＿} 离家近。\\
\emph{Zhège gōngzuò \_\_\_ gōngzī gāo, \_\_\_ lí jiā jìn.}\\
A. 只有……才…… \quad B. 只要……就…… \quad C. 既……又…… \quad D. 首先……然后……
\item 面试的时候，\zh{＿＿} 要介绍自己，\zh{＿＿} 要回答问题。\\
\emph{Miànshì de shíhou, \_\_\_ yào jièshào zìjǐ, \_\_\_ yào huídá wèntí.}\\
A. 既……又…… \quad B. 首先……然后…… \quad C. 只有……才…… \quad D. 只要……就……
\end{enumerate}
\end{exercice}

\begin{exercice}[Vrai ou faux ? Justifie ta réponse]
Dis si chaque phrase est correcte (V) ou fautive (F), puis justifie en une phrase en français.

\begin{enumerate}
\item \zh{只要多练习，才会进步。} \emph{Zhǐyào duō liànxí, cái huì jìnbù.}
\item \zh{这家公司既大又有名。} \emph{Zhè jiā gōngsī jì dà yòu yǒumíng.}
\item \zh{首先吃饭，然后我起床。} \emph{Shǒuxiān chīfàn, ránhòu wǒ qǐchuáng.}
\item \zh{只有会说中文，才能在这家公司工作。} \emph{Zhǐyǒu huì shuō Zhōngwén, cái néng zài zhè jiā gōngsī gōngzuò.}
\end{enumerate}
\end{exercice}

\begin{exercice}[Vocabulaire : complète le tableau]
Complète le tableau suivant en donnant le pinyin (avec les tons) et la traduction française de chaque connecteur.

\begin{center}
\begin{tabularx}{\linewidth}{|X|X|X|}
\hline
\rowcolor{popblueL} Caractères & Pinyin & Traduction \\
\hline
首先 & & \\
\hline
然后 & & \\
\hline
其次 & & \\
\hline
最后 & & \\
\hline
既……又…… & & \\
\hline
只有……才…… & & \\
\hline
只要……就…… & & \\
\hline
\end{tabularx}
\end{center}
\end{exercice}

\begin{exercice}[Associe les caractères au pinyin]
Relie chaque mot (1 à 6) à son pinyin (a à f).

\begin{enumerate}
\item \zh{努力} \quad a. \emph{jìnbù}
\item \zh{成功} \quad b. \emph{nǔlì}
\item \zh{进步} \quad c. \emph{miànshì}
\item \zh{面试} \quad d. \emph{chénggōng}
\item \zh{工资} \quad e. \emph{gōngzī}
\item \zh{经验} \quad f. \emph{jīngyàn}
\end{enumerate}
\end{exercice}

\courssub{Je m'entraîne}

\begin{exercice}[Texte à trous : le recrutement]
Complète le texte suivant avec les connecteurs : \zh{首先}、\zh{然后}、\zh{其次}、\zh{最后} (chaque connecteur n'est utilisé qu'une fois).

\begin{textebox}[Le processus de recrutement]
\zh{＿＿，公司在网上发布招聘信息。＿＿，他们看申请人的简历。＿＿，他们请最好的人来面试。＿＿，公司选择最合适的人。}\\
\emph{\_\_\_\_, gōngsī zài wǎngshang fābù zhāopìn xìnxī. \_\_\_\_, tāmen kàn shēnqǐngrén de jiǎnlì. \_\_\_\_, tāmen qǐng zuì hǎo de rén lái miànshì. \_\_\_\_, gōngsī xuǎnzé zuì héshì de rén.}\\
« D'abord, l'entreprise publie l'offre d'emploi sur Internet. Ensuite, elle examine les CV des candidats. Puis, elle invite les meilleurs à l'entretien. Enfin, l'entreprise choisit la personne la plus adaptée. »
\end{textebox}
\end{exercice}

\begin{exercice}[Transformation de phrases]
\begin{enumerate}
\item Réécris la phrase \zh{只要下雨，我们就不去。} \emph{Zhǐyào xià yǔ, wǒmen jiù bú qù.} (« Tant qu'il pleut, nous n'y allons pas ») avec la structure \zh{只有……才……}, puis explique en français la différence de sens entre les deux phrases.
\item Combine ces trois éléments en une seule phrase avec \zh{既……又……} : \zh{这个工作 / 工资高 / 离家近} (« ce travail / le salaire est élevé / c'est proche de la maison »).
\item Mets les mots dans le bon ordre pour former une phrase correcte :\\
\zh{才 / 只有 / 能 / 找到 / 你 / 工作经验 / 有 / 好工作}\\
\emph{cái / zhǐyǒu / néng / zhǎodào / nǐ / gōngzuò jīngyàn / yǒu / hǎo gōngzuò}
\end{enumerate}
\end{exercice}

\begin{exercice}[Corrige les phrases fautives]
Chaque phrase contient une erreur de structure. Souligne l'erreur et réécris la phrase correctement.

\begin{enumerate}
\item \zh{只要多练习，才会进步。} \emph{Zhǐyào duō liànxí, cái huì jìnbù.}
\item \zh{只有努力，就能成功。} \emph{Zhǐyǒu nǔlì, jiù néng chénggōng.}
\item \zh{他既会说英语，又会说中文也。} \emph{Tā jì huì shuō Yīngyǔ, yòu huì shuō Zhōngwén yě.}
\item \zh{最后我起床，首先我刷牙。} \emph{Zuìhòu wǒ qǐchuáng, shǒuxiān wǒ shuā yá.}
\end{enumerate}
\end{exercice}

\begin{exercice}[Traduction (version)]
Traduis les phrases suivantes en français.

\begin{enumerate}
\item \zh{只要你不放弃，就一定能找到工作。}\\
\emph{Zhǐyào nǐ bú fàngqì, jiù yídìng néng zhǎodào gōngzuò.}
\item \zh{只有学好外语，才能在大公司工作。}\\
\emph{Zhǐyǒu xuéhǎo wàiyǔ, cái néng zài dà gōngsī gōngzuò.}
\item \zh{雅温得既大又热闹，我很喜欢这个城市。}\\
\emph{Yǎwēndé jì dà yòu rènao, wǒ hěn xǐhuan zhège chéngshì.}
\end{enumerate}
\end{exercice}

\courssub{Je me prépare au Bac}

\begin{exercice}[Compréhension de texte et questions de langue]
Lis le texte suivant, puis réponds aux questions.

\begin{textebox}[Le rêve de Mballa]
\zh{姆巴拉拉是杜阿拉的一个中学生。他想以后在一家中国公司工作。他知道，只有学好中文，才能实现这个梦想。所以他首先每天学习生词，然后练习写汉字，其次听中文歌，最后周末和老师练习说话。学习中文既不容易，又很有意思。只要他不放弃，就一定会成功。}\\
\emph{Mballa shì Dù'ālā de yí ge zhōngxuéshēng. Tā xiǎng yǐhòu zài yì jiā Zhōngguó gōngsī gōngzuò. Tā zhīdào, zhǐyǒu xuéhǎo Zhōngwén, cái néng shíxiàn zhège mèngxiǎng. Suǒyǐ tā shǒuxiān měitiān xuéxí shēngcí, ránhòu liànxí xiě Hànzì, qícì tīng Zhōngwén gē, zuìhòu zhōumò hé lǎoshī liànxí shuōhuà. Xuéxí Zhōngwén jì bù róngyì, yòu hěn yǒu yìsi. Zhǐyào tā bú fàngqì, jiù yídìng huì chénggōng.}\\
« Mballa est un élève de Douala. Il veut travailler plus tard dans une entreprise chinoise. Il sait que ce n'est qu'en apprenant bien le chinois qu'il pourra réaliser ce rêve. Alors, d'abord il apprend chaque jour du vocabulaire, ensuite il s'entraîne à écrire les caractères, puis il écoute des chansons chinoises, et enfin le week-end il s'entraîne à parler avec son professeur. Apprendre le chinois est à la fois difficile et très intéressant. Tant qu'il n'abandonne pas, il réussira sûrement. »
\end{textebox}

\textbf{Questions de compréhension} (réponds en français) :
\begin{enumerate}
\item Où habite Mballa et quel est son rêve ?
\item Quelles sont les quatre étapes de sa méthode d'apprentissage ? Cite-les dans l'ordre.
\item Selon le texte, quelle condition est nécessaire pour réaliser son rêve ?
\end{enumerate}

\textbf{Questions de langue} :
\begin{enumerate}
\item Relève dans le texte une phrase avec \zh{只有……才……} et traduis-la.
\item Relève dans le texte la phrase avec \zh{既……又……} et explique le rôle de cette structure.
\item Transforme la phrase \zh{只要他不放弃，就一定会成功} en phrase avec \zh{只有……才……} et explique la différence de sens.
\end{enumerate}
\end{exercice}

\begin{exercice}[Thème et version]
\textbf{Thème.} Traduis en chinois (caractères et pinyin) :
\begin{enumerate}
\item « Ce n'est qu'en apprenant le chinois que tu pourras travailler dans cette entreprise. »
\item « Tant que tu travailles sérieusement, tu réussiras. »
\item « Ce marché de Yaoundé est à la fois grand et animé. »
\end{enumerate}

\textbf{Version.} Traduis en français :
\begin{enumerate}
\item \zh{首先你要写简历，然后去面试，最后等公司给你打电话。}\\
\emph{Shǒuxiān nǐ yào xiě jiǎnlì, ránhòu qù miànshì, zuìhòu děng gōngsī gěi nǐ dǎ diànhuà.}
\item \zh{只有多练习，才能说好中文。}\\
\emph{Zhǐyǒu duō liànxí, cái néng shuōhǎo Zhōngwén.}
\end{enumerate}
\end{exercice}

\begin{exercice}[Expression écrite guidée]
\textbf{Sujet :} Tu prépares le baccalauréat. Décris, en chinois, les étapes de ta préparation.

\textbf{Consignes :}
\begin{itemize}
\item Rédige un texte d'au moins \textbf{60 caractères chinois} (environ 5 à 6 phrases).
\item Utilise obligatoirement la suite \zh{首先……然后……其次……最后……}.
\item Utilise au moins \textbf{une fois} \zh{只有……才……} ou \zh{只要……就……}, et \textbf{une fois} \zh{既……又……}.
\item Écris ton texte en caractères, puis recopie une phrase de ton choix en pinyin.
\end{itemize}

\textbf{Aide vocabulaire :} \zh{复习} \emph{fùxí} (réviser), \zh{做练习} \emph{zuò liànxí} (faire des exercices), \zh{休息} \emph{xiūxi} (se reposer), \zh{考试} \emph{kǎoshì} (examen), \zh{紧张} \emph{jǐnzhāng} (stressé), \zh{有信心} \emph{yǒu xìnxīn} (avoir confiance).
\end{exercice}

\newpage

\coursec{Corrigés des exercices}

\begin{corrige}[Exercice 1 — QCM : choisir le bon connecteur]
\begin{enumerate}
\item \textbf{Réponse B (才).} Phrase complète : \zh{只有努力学习，才能考上好大学。} \emph{Zhǐyǒu nǔlì xuéxí, cái néng kǎoshang hǎo dàxué.} « Ce n'est qu'en étudiant avec effort que tu pourras entrer dans une bonne université. » Justification : \zh{只有} exige impérativement \zh{才} dans la proposition principale (condition nécessaire).
\item \textbf{Réponse B (就).} Phrase complète : \zh{只要天气好，我们就去踢足球。} \emph{Zhǐyào tiānqì hǎo, wǒmen jiù qù tī zúqiú.} « Tant qu'il fait beau, nous allons jouer au football. » Justification : \zh{只要} se combine avec \zh{就} (condition suffisante).
\item \textbf{Réponse C (既……又……).} Phrase complète : \zh{这个工作既工资高，又离家近。} \emph{Zhège gōngzuò jì gōngzī gāo, yòu lí jiā jìn.} « Ce travail a à la fois un salaire élevé et est proche de la maison. » Justification : on cumule deux qualités du même sujet, donc \zh{既……又……}.
\item \textbf{Réponse B (首先……然后……).} Phrase complète : \zh{面试的时候，首先要介绍自己，然后要回答问题。} \emph{Miànshì de shíhou, shǒuxiān yào jièshào zìjǐ, ránhòu yào huídá wèntí.} « Lors de l'entretien, il faut d'abord se présenter, puis répondre aux questions. » Justification : il s'agit d'une suite d'étapes ordonnées dans le temps.
\end{enumerate}
\textbf{Point de vigilance :} ne jamais croiser les paires : \zh{只有……就……} et \zh{只要……才……} sont fautives.
\end{corrige}

\begin{corrige}[Exercice 2 — Vrai ou faux ? Justifie ta réponse]
\begin{enumerate}
\item \textbf{F.} \zh{只要} doit s'associer à \zh{就}, pas à \zh{才}. Correction : \zh{只要多练习，就会进步。} \emph{Zhǐyào duō liànxí, jiù huì jìnbù.} « Tant que tu t'entraînes beaucoup, tu progresseras. »
\item \textbf{V.} La structure \zh{既……又……} cumule correctement deux caractéristiques (grand + célèbre) du même sujet \zh{这家公司}. « Cette entreprise est à la fois grande et célèbre. »
\item \textbf{F.} L'ordre logique est inversé : on se lève avant de manger ; la suite des connecteurs doit suivre l'ordre chronologique réel. Correction : \zh{首先我起床，然后吃饭。} \emph{Shǒuxiān wǒ qǐchuáng, ránhòu chīfàn.} « D'abord je me lève, ensuite je mange. »
\item \textbf{V.} Structure \zh{只有……才……} correcte : la condition (savoir parler chinois) est présentée comme nécessaire pour travailler dans cette entreprise. « Ce n'est qu'en sachant parler chinois qu'on peut travailler dans cette entreprise. »
\end{enumerate}
\end{corrige}

\begin{corrige}[Exercice 3 — Vocabulaire : complète le tableau]
\begin{center}
\begin{tabularx}{\linewidth}{|X|X|X|}
\hline
\rowcolor{popblueL} Caractères & Pinyin & Traduction \\
\hline
首先 & \emph{shǒuxiān} & d'abord, en premier \\
\hline
然后 & \emph{ránhòu} & ensuite, puis \\
\hline
其次 & \emph{qícì} & ensuite, en second lieu \\
\hline
最后 & \emph{zuìhòu} & enfin, finalement \\
\hline
既……又…… & \emph{jì……yòu……} & à la fois... et... \\
\hline
只有……才…… & \emph{zhǐyǒu……cái……} & seulement si... alors... (condition nécessaire) \\
\hline
只要……就…… & \emph{zhǐyào……jiù……} & tant que... alors... (condition suffisante) \\
\hline
\end{tabularx}
\end{center}
\textbf{Point de vigilance :} attention aux tons de \zh{既} \emph{jì} (4\up{e} ton), différent de \zh{即} \emph{jí} (2\up{e} ton) ; et de \zh{才} \emph{cái} (2\up{e} ton).
\end{corrige}

\begin{corrige}[Exercice 4 — Associe les caractères au pinyin]
\begin{enumerate}
\item \zh{努力} → \textbf{b.} \emph{nǔlì} (faire des efforts)
\item \zh{成功} → \textbf{d.} \emph{chénggōng} (réussir, succès)
\item \zh{进步} → \textbf{a.} \emph{jìnbù} (progresser, progrès)
\item \zh{面试} → \textbf{c.} \emph{miànshì} (entretien d'embauche)
\item \zh{工资} → \textbf{e.} \emph{gōngzī} (salaire)
\item \zh{经验} → \textbf{f.} \emph{jīngyàn} (expérience)
\end{enumerate}
\textbf{Point de vigilance :} ne pas confondre \zh{经验} \emph{jīngyàn} (expérience) et \zh{经历} \emph{jīnglì} (vécu, traverser) ; attention aussi au 4\up{e} ton de \zh{试} \emph{shì} dans \zh{面试}.
\end{corrige}

\begin{corrige}[Exercice 5 — Texte à trous : le recrutement]
Texte complété :

\zh{首先，公司在网上发布招聘信息。然后，他们看申请人的简历。其次，他们请最好的人来面试。最后，公司选择最合适的人。}

\emph{Shǒuxiān, gōngsī zài wǎngshang fābù zhāopìn xìnxī. Ránhòu, tāmen kàn shēnqǐngrén de jiǎnlì. Qícì, tāmen qǐng zuì hǎo de rén lái miànshì. Zuìhòu, gōngsī xuǎnzé zuì héshì de rén.}

« D'abord, l'entreprise publie l'offre d'emploi sur Internet. Ensuite, elle examine les CV des candidats. Puis, elle invite les meilleurs à l'entretien. Enfin, l'entreprise choisit la personne la plus adaptée. »

\textbf{Justification :} l'ordre suit la chronologie du recrutement : publication de l'offre (étape 1 → \zh{首先}), lecture des CV (étape 2 → \zh{然后}), convocation à l'entretien (étape 3 → \zh{其次}), choix final (étape 4 → \zh{最后}). \zh{其次} ne peut pas ouvrir la série : il suppose qu'une première étape a déjà été annoncée.
\end{corrige}

\begin{corrige}[Exercice 6 — Transformation de phrases]
\begin{enumerate}
\item Transformation : \zh{只有不下雨，我们才去。} \emph{Zhǐyǒu bú xià yǔ, wǒmen cái qù.} « Ce n'est que s'il ne pleut pas que nous irons. »

\textbf{Différence de sens :} avec \zh{只要……就……}, la pluie suffit à bloquer la sortie (condition suffisante : dès qu'il pleut, on ne va pas ; mais s'il ne pleut pas, rien n'est garanti). Avec \zh{只有……才……}, l'absence de pluie devient la condition \emph{nécessaire} du départ : s'il pleut, c'est impossible. La seconde formulation est plus restrictive : elle présente « ne pas pleuvoir » comme le passage obligé.
\item \zh{这个工作既工资高，又离家近。} \emph{Zhège gōngzuò jì gōngzī gāo, yòu lí jiā jìn.} « Ce travail a à la fois un salaire élevé et est proche de la maison. » Les deux caractéristiques portent sur le même sujet \zh{这个工作}, placé une seule fois avant \zh{既}.
\item Ordre correct : \zh{你只有有工作经验，才能找到好工作。} ou plus naturellement : \zh{只有有工作经验，你才能找到好工作。} \emph{Zhǐyǒu yǒu gōngzuò jīngyàn, nǐ cái néng zhǎodào hǎo gōngzuò.} « Ce n'est qu'en ayant de l'expérience professionnelle que tu pourras trouver un bon travail. » Schéma : \zh{只有} + condition + sujet + \zh{才} + résultat.
\end{enumerate}
\end{corrige}

\begin{corrige}[Exercice 7 — Corrige les phrases fautives]
\begin{enumerate}
\item Erreur : \zh{才} après \zh{只要}. Correction : \zh{只要多练习，就会进步。} \emph{Zhǐyào duō liànxí, jiù huì jìnbù.} « Tant que tu t'entraînes beaucoup, tu progresseras. » (\zh{只要} appelle \zh{就}.)
\item Erreur : \zh{就} après \zh{只有}. Correction : \zh{只有努力，才能成功。} \emph{Zhǐyǒu nǔlì, cái néng chénggōng.} « Ce n'est qu'en faisant des efforts qu'on peut réussir. » (\zh{只有} appelle \zh{才}.)
\item Erreur : \zh{也} en fin de phrase, redondant avec \zh{既……又……}. Correction : \zh{他既会说英语，又会说中文。} \emph{Tā jì huì shuō Yīngyǔ, yòu huì shuō Zhōngwén.} « Il sait parler à la fois anglais et chinois. » En chinois, on ne cumule pas deux marqueurs d'addition comme en français (« à la fois... et... aussi »).
\item Erreur : ordre chronologique inversé (on se lève avant de se brosser les dents) et connecteurs mal placés. Correction : \zh{首先我起床，最后我刷牙。} ou plus naturel : \zh{我首先起床，然后刷牙。} \emph{Wǒ shǒuxiān qǐchuáng, ránhòu shuā yá.} « D'abord je me lève, ensuite je me brosse les dents. »
\end{enumerate}
\textbf{Point de vigilance :} l'erreur n°1 et n°2 (croisement \zh{只要……才……} / \zh{只有……就……}) est la faute la plus fréquente des francophones : mémorise les deux paires comme des blocs inséparables.
\end{corrige}

\begin{corrige}[Exercice 8 — Traduction (version)]
\begin{enumerate}
\item \zh{只要你不放弃，就一定能找到工作。} → « Tant que tu n'abandonnes pas, tu pourras sûrement trouver un travail. » (Condition suffisante : ne pas abandonner suffit à garantir le résultat.)
\item \zh{只有学好外语，才能在大公司工作。} → « Ce n'est qu'en apprenant bien une langue étrangère qu'on peut travailler dans une grande entreprise. » (Condition nécessaire : la langue est le passage obligé.)
\item \zh{雅温得既大又热闹，我很喜欢这个城市。} → « Yaoundé est à la fois grande et animée, j'aime beaucoup cette ville. »
\end{enumerate}
\textbf{Point de vigilance :} \zh{放弃} \emph{fàngqì} = abandonner ; \zh{热闹} \emph{rènao} = animé, bruyant de vie (faux ami : ce n'est pas « chaud » au sens de la température, malgré \zh{热}).
\end{corrige}

\begin{corrige}[Exercice 9 — Compréhension de texte et questions de langue]
\textbf{Questions de compréhension :}
\begin{enumerate}
\item Mballa habite à Douala (\zh{杜阿拉}). Son rêve est de travailler plus tard dans une entreprise chinoise (\zh{他想以后在一家中国公司工作}).
\item Les quatre étapes, dans l'ordre : d'abord (\zh{首先}) il apprend chaque jour du vocabulaire ; ensuite (\zh{然后}) il s'entraîne à écrire les caractères ; puis (\zh{其次}) il écoute des chansons chinoises ; enfin (\zh{最后}) le week-end il s'entraîne à parler avec son professeur.
\item La condition nécessaire est d'apprendre bien le chinois : \zh{只有学好中文，才能实现这个梦想} (« ce n'est qu'en apprenant bien le chinois qu'il pourra réaliser ce rêve »).
\end{enumerate}
\textbf{Questions de langue :}
\begin{enumerate}
\item Phrase relevée : \zh{只有学好中文，才能实现这个梦想。} \emph{Zhǐyǒu xuéhǎo Zhōngwén, cái néng shíxiàn zhège mèngxiǎng.} Traduction : « Ce n'est qu'en apprenant bien le chinois qu'il peut réaliser ce rêve. »
\item Phrase relevée : \zh{学习中文既不容易，又很有意思。} \emph{Xuéxí Zhōngwén jì bù róngyì, yòu hěn yǒu yìsi.} Rôle : \zh{既……又……} cumule deux caractéristiques du même sujet (apprendre le chinois) : « à la fois pas facile et très intéressant ». Les deux qualités coexistent sans contradiction.
\item Transformation : \zh{只有他不放弃，才会成功。} ou mieux : \zh{只有不放弃，他才会成功。} \emph{Zhǐyǒu bú fàngqì, tā cái huì chénggōng.} Différence de sens : avec \zh{只要……就……}, « ne pas abandonner » suffit à garantir la réussite (condition suffisante, optimiste) ; avec \zh{只有……才……}, « ne pas abandonner » est présenté comme la condition indispensable, sans laquelle la réussite est impossible (condition nécessaire, plus exigeante).
\end{enumerate}
\textbf{Barème indicatif (sur 20) :} compréhension 6 pts (2 pts par question) ; questions de langue 9 pts (3 pts par question : relevé exact 1 pt, traduction ou transformation correcte 1 pt, explication grammaticale 1 pt) ; qualité de l'expression française 5 pts.
\end{corrige}

\begin{corrige}[Exercice 10 — Thème et version]
\textbf{Thème (propositions de traduction) :}
\begin{enumerate}
\item \zh{只有学好中文，你才能在这家公司工作。}\\
\emph{Zhǐyǒu xuéhǎo Zhōngwén, nǐ cái néng zài zhè jiā gōngsī gōngzuò.}\\
Point de méthode : « ce n'est qu'en... que » se rend par \zh{只有……才……} ; le verbe modal \zh{能} (pouvoir) se place après \zh{才}.
\item \zh{只要你认真工作，你就会成功。}\\
\emph{Zhǐyào nǐ rènzhēn gōngzuò, nǐ jiù huì chénggōng.}\\
Point de méthode : « tant que » se rend par \zh{只要……就……} ; \zh{认真} \emph{rènzhēn} = sérieux, consciencieux.
\item \zh{雅温得的这个市场既大又热闹。}\\
\emph{Yǎwēndé de zhège shìchǎng jì dà yòu rènao.}\\
Point de méthode : « à la fois... et... » se rend par \zh{既……又……} ; le sujet n'est dit qu'une fois, avant \zh{既}.
\end{enumerate}
\textbf{Version :}
\begin{enumerate}
\item « D'abord tu dois rédiger ton CV, ensuite aller à l'entretien, et enfin attendre que l'entreprise t'appelle. »
\item « Ce n'est qu'en t'entraînant beaucoup que tu pourras bien parler chinois. »
\end{enumerate}
\textbf{Barème indicatif (sur 20) :} thème 12 pts (4 pts par phrase : choix du connecteur 2 pts, ordre des mots 1 pt, caractères et pinyin 1 pt) ; version 8 pts (4 pts par phrase : sens global 2 pts, rendu précis du connecteur 2 pts).
\end{corrige}

\begin{corrige}[Exercice 11 — Expression écrite guidée]
\textbf{Proposition de rédaction modèle} (environ 90 caractères) :

\zh{下个月我要参加毕业考试。首先，我每天复习汉语生词和语法。然后，我做很多练习，特别是翻译练习。其次，我周末和同学一起复习，我们互相帮助。最后，考试以前，我要好好休息。复习既紧张又有意思。只有认真准备，我才能考好。只要我不放弃，就一定会成功！}

\emph{Xià ge yuè wǒ yào cānjiā bìyè kǎoshì. Shǒuxiān, wǒ měitiān fùxí Hànyǔ shēngcí hé yǔfǎ. Ránhòu, wǒ zuò hěn duō liànxí, tèbié shì fānyì liànxí. Qícì, wǒ zhōumò hé tóngxué yìqǐ fùxí, wǒmen hùxiāng bāngzhù. Zuìhòu, kǎoshì yǐqián, wǒ yào hǎohao xiūxi. Fùxí jì jǐnzhāng yòu yǒu yìsi. Zhǐyǒu rènzhēn zhǔnbèi, wǒ cái néng kǎohǎo. Zhǐyào wǒ bú fàngqì, jiù yídìng huì chénggōng!}

« Le mois prochain, je passe l'examen de fin d'études. D'abord, je révise chaque jour le vocabulaire et la grammaire du chinois. Ensuite, je fais beaucoup d'exercices, surtout des exercices de traduction. Puis, le week-end, je révise avec mes camarades et nous nous aidons mutuellement. Enfin, avant l'examen, je dois bien me reposer. Réviser est à la fois stressant et intéressant. Ce n'est qu'en me préparant sérieusement que je pourrai réussir. Tant que je n'abandonne pas, je réussirai sûrement ! »

\textbf{Vérification des consignes :} suite \zh{首先……然后……其次……最后……} présente et dans l'ordre ; \zh{既……又……} utilisé (\zh{复习既紧张又有意思}) ; \zh{只有……才……} et \zh{只要……就……} tous deux employés ; plus de 60 caractères.

\textbf{Barème indicatif (sur 20) :} respect des consignes et longueur 4 pts ; correction des quatre connecteurs de séquence 4 pts ; emploi correct de \zh{既……又……} 3 pts ; emploi correct d'une structure de condition 3 pts ; correction générale (ordre des mots, caractères, tons du pinyin) 4 pts ; richesse du vocabulaire et cohérence 2 pts.

\textbf{Points de vigilance :}
\begin{itemize}
\item Ne pas écrire \zh{只有……就……} ni \zh{只要……才……} : vérifier chaque paire avant de rendre la copie.
\item \zh{首先} peut se placer avant ou après le sujet, mais \zh{其次} ne doit jamais ouvrir la série.
\item Dans \zh{既……又……}, les deux éléments cumulés doivent être de même nature (deux adjectifs ou deux verbes) et porter sur le même sujet.
\item Compter ses caractères : la consigne des 60 caractères est éliminatoire si elle n'est pas respectée.
\end{itemize}
\end{corrige}

\newpage

\coursec{Fiche bilan}

\courssub{Fiche bilan — Leçon 31 : Les connecteurs logiques}

\begin{popfigure}
\begin{tikzpicture}[
    mindmap,
    grow cyclic,
    every node/.style={concept, execute at begin node=\small\bfseries},
    concept color=popblue,
    level 1/.append style={level distance=4.5cm, sibling angle=360/7},
    level 2/.append style={level distance=3cm, sibling angle=360/4},
    text=white,
    font=\small
]
\node[concept, scale=1.2, align=center]{Connecteurs\\logiques (Leçon 31)}
    child[concept color=poporange] { node[concept, align=center]{Séquence\\首先…然后…\\其次…最后…} }
    child[concept color=popgreen] { node[concept, align=center]{Cumul\\既…又…} }
    child[concept color=poppurple] { node[concept, align=center]{Condition\\nécessaire\\只有…才…} }
    child[concept color=poppink] { node[concept, align=center]{Condition\\suffisante\\只要…就…} }
    child[concept color=popgold] { node[concept, align=center]{Comparaison\\FR / CN} }
    child[concept color=popmuted] { node[concept, align=center]{Pièges\\francophones} }
    child[concept color=popdark] { node[concept, align=center]{Production\\écrite/orale} };
\end{tikzpicture}
\legende{Carte mentale des connecteurs logiques : séquence, cumul, conditions nécessaire et suffisante, contrastes français/chinois et points de vigilance.}
\end{popfigure}

\courssub{Lexique clé (mots-outils et adverbes de liaison)}

\begin{center}
\begin{tabularx}{\linewidth}{|X|X|X|X|}
\hline
\rowcolor{popblueL} \textbf{Caractères} & \textbf{Pinyin} & \textbf{Nature} & \textbf{Traduction} \\ \hline
首先 & shǒuxiān & adv. & tout d'abord, en premier lieu \\ \hline
然后 & ránhòu & adv. & ensuite, puis \\ \hline
其次 & qícì & adv. & en second lieu, de plus \\ \hline
最后 & zuìhòu & adv. & enfin, en dernier lieu \\ \hline
既 & jì & conj. & à la fois (corrélatif) \\ \hline
又 & yòu & adv./conj. & aussi, également (corrélatif) \\ \hline
只有 & zhǐyǒu & conj. & seulement si (condition nécessaire) \\ \hline
才 & cái & adv. & seulement alors (marque la conséquence) \\ \hline
只要 & zhǐyào & conj. & du moment que, il suffit que (condition suffisante) \\ \hline
就 & jiù & adv. & alors, dès lors (marque la conséquence) \\ \hline
\end{tabularx}
\end{center}

\courssub{Règles de grammaire à connaître par cœur}

\begin{center}
\begin{tabularx}{\linewidth}{|>{\raggedright\arraybackslash}X|>{\centering\arraybackslash}X|>{\raggedright\arraybackslash}X|}
\hline
\rowcolor{popblueL} \textbf{Structure} & \textbf{Schéma / Exemple} & \textbf{Emploi / Nuance} \\ \hline
\textbf{Séquence ordonnée} & \zh{首先……，然后……，其次……，最后……} & Énumérer des étapes chronologiques ou logiques. L'ordre est fixe : \zh{首先} ouvre, \zh{然后} et \zh{其次} développent, \zh{最后} conclut. Ne pas inverser \zh{其次} et \zh{然后}. \\ \hline
\textbf{Cumul de qualités} & \zh{既……又……} (Sujet + 既 + Adj1 + 又 + Adj2) & Deux caractéristiques \textbf{simultanées}, souvent positives. Adjectifs monosyllabiques ou bisyllabiques, de même valence. \\ \hline
\textbf{Condition nécessaire} & \zh{只有……，才……} (Condition unique $\rightarrow$ Résultat) & « Ce n'est que si… que… ». \cle{才} est obligatoire, placé devant le verbe de la principale. \\ \hline
\textbf{Condition suffisante} & \zh{只要……，就……} (Condition large $\rightarrow$ Résultat automatique) & « Du moment que… », « Il suffit que… ». \cle{就} marque la conséquence immédiate ou automatique. \\ \hline
\textbf{Distinction 只有 / 只要} & \zh{只有你努力，才能成功。} vs \zh{只要你努力，就能成功。} & \textbf{只有…才} : condition \textbf{unique et indispensable} (sans elle, impossible). \textbf{只要…就} : condition \textbf{suffisante} (elle suffit à produire le résultat). \\ \hline
\end{tabularx}
\end{center}

\begin{exemplebox}[Deux conditions, deux nuances]
\zh{只有他来了，我才走。} \emph{Zhǐyǒu tā lái le, wǒ cái zǒu.} « Ce n'est que s'il vient que je partirai. » (sa venue est indispensable)\\[2pt]
\zh{只要天气好，我们就去海边。} \emph{Zhǐyào tiānqì hǎo, wǒmen jiù qù hǎibiān.} « Du moment qu'il fait beau, nous allons au bord de la mer. » (le beau temps suffit)
\end{exemplebox}

\courssub{Méthodes express}

\begin{itemize}
    \item \textbf{Ordonner un récit ou un processus} : identifier les étapes $\rightarrow$ associer \zh{首先 / 然后 / 其次 / 最后} $\rightarrow$ vérifier l'ordre logique (chronologie ou importance).
    \item \textbf{Décrire une personne ou un objet (cumul)} : choisir deux adjectifs compatibles $\rightarrow$ les insérer dans \zh{既…又…}. Attention : pour la négation, on n'utilise pas *\zh{既不…又不…} mais \zh{既不…也不…} (« ni… ni… »).
    \item \textbf{Traduire une condition (français $\rightarrow$ chinois)} :
    \begin{enumerate}
        \item Repérer « seulement si / à condition que / il faut que » $\rightarrow$ \zh{只有…才…}.
        \item Repérer « du moment que / il suffit que / pourvu que » $\rightarrow$ \zh{只要…就…}.
        \item Placer \cle{才} ou \cle{就} juste devant le verbe (ou l'auxiliaire) de la proposition principale.
    \end{enumerate}
    \item \textbf{Éviter l'erreur d'ordre des mots} : en chinois, la subordonnée introduite par \zh{只有} ou \zh{只要} précède \textbf{toujours} la principale contenant \zh{才} ou \zh{就}, contrairement au français où la conséquence peut ouvrir la phrase (« Alors je partirai, seulement s'il vient »).
\end{itemize}

\begin{attention}[Erreurs fréquentes des francophones]
\begin{itemize}
    \item *\zh{如果只有…} : ne pas cumuler \zh{如果} et \zh{只有} ; l'un des deux suffit.
    \item *\zh{只要…才…} : mélange interdit. \zh{只要} s'accorde avec \zh{就}, \zh{只有} avec \zh{才}.
    \item Ne pas placer \zh{才} ou \zh{就} en tête de phrase : ce sont des adverbes préverbaux.
    \item Pas de virgule entre \zh{只有} et le sujet, ni entre \zh{才} et le verbe.
\end{itemize}
\end{attention}

\begin{aretenir}[Checklist avant le Bac]
\begin{itemize}
    \item $\square$ Je sais écrire et ordonner \zh{首先、然后、其次、最后} sans les confondre.
    \item $\square$ Je construis correctement \zh{既…又…} avec deux adjectifs de même valence (souvent positifs), et \zh{既不…也不…} pour la négation.
    \item $\square$ Je distingue \textbf{condition nécessaire} (\zh{只有…才…}) et \textbf{condition suffisante} (\zh{只要…就…}) sur un exemple camerounais (réussite au baccalauréat, obtention d'un visa pour la Chine).
    \item $\square$ Je place \textbf{obligatoirement} \zh{才} et \zh{就} devant le verbe ou l'auxiliaire de la principale, jamais en début de phrase.
    \item $\square$ Je ne mets pas de virgule entre \zh{只有} et le sujet, ni entre \zh{才} et le verbe.
    \item $\square$ J'évite les gallicismes : *\zh{如果只有…} (redondant), *\zh{只要…才…} (mélange interdit).
    \item $\square$ Je sais transformer une phrase en \zh{只有…才…} en phrase avec \zh{除非…否则…} (« à moins que… sinon… », niveau avancé).
    \item $\square$ Je rédige un paragraphe de 5 à 6 phrases utilisant au moins trois de ces connecteurs pour décrire mon parcours d'orientation ou une recette (ndolé, riz cantonais).
    \item $\square$ À l'oral, j'enchaîne les connecteurs de séquence avec la bonne intonation, en marquant une courte pause après chaque connecteur.
    \item $\square$ Je reconnais les caractères clés \zh{首、既、才、要} et leurs clés (radicaux) : \zh{首} (clé 首, « tête »), \zh{既} (clé 旡), \zh{才} (clé 扌, « main »), \zh{要} (clé 覀, variante de 西).
    \item $\square$ Je traduis correctement « Ce n'est que s'il vient que je partirai » $\rightarrow$ \zh{只有他来了，我才走。} (\emph{Zhǐyǒu tā lái le, wǒ cái zǒu.}).
\end{itemize}
\end{aretenir}

\findecours

\end{document}
